% ============================================================
%  H. Brøchner, Problemet om Tro og Viden
%  A Historical-Critical Treatise
%  Copenhagen: P. G. Philipsen, 1868.
%  TRANSLATION (English)
%  Translated from the Danish transcription (transcription.tex).
% ============================================================
\documentclass[12pt,a4paper]{book}
\usepackage[utf8]{inputenc}
\usepackage[T1]{fontenc}
\usepackage{libertinus}
\usepackage{libertinust1math}
\usepackage{textalpha}   % Greek words set as glyphs (τέλος, πίστις, γνῶσις…) for parity with the Danish
\usepackage[english]{babel}
\usepackage[protrusion=true,expansion=true]{microtype}
\usepackage{setspace}
\usepackage{enumitem}
\usepackage[margin=1.2in]{geometry}
\usepackage{fancyhdr}
\setlength{\headheight}{28pt}
\addtolength{\topmargin}{-13.5pt}
\usepackage{hyperref}

\hypersetup{
  pdftitle={The Problem of Faith and Knowledge},
  pdfauthor={H. Brøchner},
  hidelinks
}

\pagestyle{fancy}
\fancyhf{}
\fancyhead[LE,RO]{\thepage}
\fancyhead[RE]{\textit{The Problem of Faith and Knowledge}}
\fancyhead[LO]{\textit{\leftmark}}

\setstretch{1.3}

\title{\textbf{The Problem of Faith and Knowledge}\\[6pt]
  \large A Historical-Critical Treatise\\[6pt]
  \normalsize\textit{Problemet om Tro og Viden}}
\author{H.\ Brøchner\\[4pt]
  \small Translated from the Danish}
\date{Copenhagen: P.\ G.\ Philipsen, 1868}

\begin{document}
\maketitle
\thispagestyle{empty}
\newpage

\tableofcontents
\newpage

%% ============================================================
%%  PREFACE (pp. 1--7)
%% ============================================================
\chapter*{Preface}
\addcontentsline{toc}{chapter}{Preface}
\label{ch:preface}

\noindent This work reproduces, in all essential respects unchanged,
the content of the lectures I delivered at the University of Copenhagen
in the autumn semester of 1867 on the relation between faith and
knowledge, between religion, philosophy, and ethics. These lectures
encompassed the historical-critical part of the investigation and
formed the preparation for an attempt at an independent resolution of
the problem under discussion, which will be published in the course of
this year.

The historical-critical part of the investigation provides, beyond a
brief sketch of the problem's genesis, a detailed critique of those
proposed solutions that have acquired a special significance in our
literature --- among them one, the solution originating with
S.\ Kierkegaard, that was executed with an originality, a seriousness,
and a depth that will preserve for it this significance at all times.
Through the critique of these proposed solutions --- which, in the
course of their development and justification, have stood in a
sustained polemical or mediating relation to earlier and foreign
contributions, and which can partly be regarded as results of
previously dominant tendencies, and have partly substantially
accomplished the critique of such tendencies --- the most significant
of the older and contemporary proposals have been illuminated, and in
part the complete preconditions have been supplied for a final judgment
upon them. This historical critique will then find its supplement in
the work that attaches itself to this one.

The present state of the problem in our literature is marked by the
conflict between, on the one side, a form of theology that borrows in
an eclectic fashion from newer philosophical hybrid formations the
weapons with which it will defend the union of faith and knowledge,
and, on the other side, the anti-theological-religious standpoint,
which, by means of the principled opposition between faith and
knowledge, aims to secure the validity of religion and science in
their mutual independence and the possibility of their peaceful
coexistence within consciousness --- while at the same time, by virtue
of this same principled opposition, denying theology all right to bear
the name of science, a name theology claims for itself by demanding to
be a ``science of faith'' [\textit{Troesvidenskab}].

The philosophizing theology, which maintains the possibility of
appropriating and expressing the content of faith in the form of
knowledge, and which therefore not only wishes to claim the name of
science for theology but assigns to it, as the concluding form of
revelatory science, primacy among the sciences and judicial authority
over them, has had as its historical point of departure the speculative
system that, through an ambiguous extension of the concept of
``religion,'' brought religion in its specific sense into an identical
relation with philosophy, from which it was to be distinguished not by
any difference of content, but only by a difference of form. After
this ambiguous unity had to be abandoned in the struggles that
developed following the emergence of Hegel's philosophy of religion,
the theology influenced by philosophy, exchanging the designation of
speculation for that of the Christian, nevertheless preserved the
tendency toward mediation; but it now rests primarily upon a different
philosophical orientation --- and in particular upon the system that,
by introducing a redoubling into philosophy, creates room for a ``free
thinking'' which borrows from its close relative, the imagination, the
``magic cloak'' [\textit{Fjederham}] enabling it to soar above the
abysses of mystery; and it cautiously inserts a qualification that cuts
off even the ``higher science's'' consequences whenever thought,
despite its ``freedom,'' involuntarily comes to follow its own
essential laws.

The anti-theological theory, which likewise presupposes Hegelian
philosophy, but whose relation to that presupposition is originally
negative and polemical, had as its historical incitement also the
protest --- lodged by a philosophy that presents itself at once as the
consequence of and as the opposition to that speculative doctrine of
mediation --- against the unity of religion and philosophy. The most
significant representative of this tendency is Feuerbach --- a man
whose writings are known only very incompletely and judged only very
superficially by those in our literature who have criticized him.
Conceiving religion's standpoint as practical in opposition to
science's theoretical standpoint, and determining the former ever more
sharply through the emphasis on its alogical element, he arrived at
the irreconcilable opposition as his view of the relation, and
demanded, in the interest of genuine cognition and genuine ethical
action, that religion be absorbed into science and ethics. Religion as
such was to disappear --- as something alogical, as an imperfect
transitional stage in the development of the human spirit's
self-consciousness --- leaving behind only its name in an improper and
figurative sense.

In agreement with Feuerbach regarding the opposition, but disagreeing
as to its grounds and its application, S.\ Kierkegaard draws the
distinction between religion and philosophy in religion's interest.
The human spirit, which in the restlessness of existence cannot grasp
the Eternal in the objective certainty of cognition, and which
acknowledges the limits of its own cognition, finds only in the
religious relation, in faith, a subjective existential certainty
regarding what is its deepest interest; it finds there, and only
there, its true reconciliation. The validity of knowledge is not
denied, but objective knowledge is bounded to a domain in which the
spirit's deepest interests cannot find their satisfaction. Faith and
knowledge thus each receive their own territory, and through this, and
through knowledge's subordination, the compatibility of both within
the same consciousness is to become possible.

This Kierkegaardian conception of the problem of the relation between
faith and knowledge has been taken up by Professor R.\ Nielsen and
presented in such a way that some have wished to place him on a par in
originality with the theory's originator, attributing to him the
resolution of a task of the same creativity as the one Kierkegaard
strove to resolve. How little this view is warranted will be
demonstrated in this work. It will emerge that the problem has not
been advanced by Professor Nielsen; that the form he has given the
solution has introduced only uncertainty and vacillation into the
decisive determinations; that the attempt at a psychological and
metaphysical grounding of the separation of the principle of faith
from the principle of knowledge as absolutely heterogeneous with it
has failed and leads to irresolvable contradictions; and that the same
holds for his grounding of the proposition that the absolutely
heterogeneous principles can be united within the same consciousness.

When, then, alongside S.\ Kierkegaard --- the genuine representative
of the anti-theological-religious standpoint --- such detailed
attention is devoted to Professor Nielsen, it is not because the
theory has received from him a deeper grounding or a richer
development, but because it has through him exchanged the character of
an esoteric doctrine for an exoteric one, penetrated into the broad
circles of ``the cultivated,'' and in these circles is in the process
of changing from a question of knowledge into a question of faith, and
of being fixed as dogma by virtue of an unappealable
\textit{ipse dixit}.

As representative of the philosophizing theology, Bishop Martensen
appears in this work; but since the argument is closely tied to what
is given in his writings, the universal, the typical in the standpoint
is nowhere lost from view, nor forgotten in favor of the individual.
It has not been necessary to give direct consideration to other
theological attempts in our literature.

The ultimate task, whose resolution is prepared through the critique
--- which immediately can yield only a negative result --- is the
vindication of a standpoint that has not found expression in our
discussions, and that has been advanced elsewhere only in forms that
cannot satisfy. This standpoint maintains such an opposition between
science --- philosophy in particular --- and the positive forms of
religion, between knowledge and faith understood as revelatory faith,
that knowledge and this faith cannot be brought to unity within the
same consciousness; --- it maintains the impossibility of a genuine
and substantial ethics, one that encompasses the whole, real human
life, developing on the basis of positive religion with its
supernatural law of action; --- and it maintains the inner,
indissoluble unity between genuine cognition and genuine ethical
action, the necessary dependence of the latter upon the former. But
while it maintains this opposition, it asserts with equal
decisiveness, on the other side --- distinguishing the essential
determination of the religious from the phenomenal forms of the
positive religions, in which it sees only a partial, imperfect,
finitizing expression of the former, and thus distinguishing what is
religious in the religions from the religions in their historical
forms --- the inner unity of the religious, the philosophical, and the
ethical, and through this inner unity the inner unity of human life:
the condition for its true and complete reconciliation, the
reconciliation within reality, the reconciliation of undivided human
nature with itself in its essential ground. In affirming the
essentiality and necessity of the religious relation for the human
being, whose self-consciousness is first fully actualized therein,
through God-consciousness, and who according to the fundamental
relation of its nature has never been and never will be without
religion; --- in further seeing ethical action and genuine cognition as
gaining a higher significance through their unity with the religious
fundamental relation; --- and in allowing genuine knowledge and genuine
ethical action to become an adequate form of actuality for the
religious, it conceives this knowledge and this action as containing
within themselves a Highest and Unconditioned, in which what is
essential in the religions is preserved. This standpoint we designate
as \textit{the religiously reconciled humane consciousness}.

It differs from the position occupied by Feuerbach not merely in that
its metaphysical fundamental concept is a different one --- one upon
which a cognition and an action bearing the stamp of ideality and
universal validity can be built, and from which the concept of human
freedom can be vindicated and a religious relation grounded to the
Absolute, with which the human being is in essential unity --- but
also in that it combats theology and dogma not only in the interest of
ethics and science, but first and foremost in the interest of religion
itself, in the interest of the infinity of the religious relation.

It differs from the old Hegelian school's conception of religion's
relation to cognition, in that it does not see the religious merely as
a form of knowledge, but as the expression of a fundamental relation
encompassing the whole of human nature in its concreteness, and in
that it does not let the medium of spiritual life be an illusory
abstract element of eternity, but reality with its real conditions and
limits.

With the anti-theological-religious theory it concurs in its
conception of theology's relation to science. Since it rests upon
scientific cognition and, in undivided human nature, assigns primacy
to cognition, it must be its task to vindicate the clear concept of
science and science's autonomy, grounded in its own nature, against
equivocations and against attempts, based on equivocations, to
subordinate science to the dominion of a foreign, heterogeneous
authority. But while there is thus, in an essential point, an
agreement between it and that theory, it enters into a decided
opposition to the latter when the relation between faith and
knowledge, between religion, philosophy, and ethics, is to be
determined. This opposition concerns Kierkegaard's conception, with
respect to which it will be shown that it undermines the certainty of
all cognition and allows the ethical to disappear into an abstract and
reality-less religious sphere. But it concerns even more precisely the
form in which the theory has been elaborated and presented by Nielsen
--- a form that Kierkegaard would certainly not recognize as his own.
The form in which Nielsen develops the theory leads, despite all the
more or less adroit turns that are made in order to avoid this
consequence, to an actual dualism in human nature, to a division of
the life of the spirit from its very ground by means of the ``absolutely
heterogeneous'' principles, to the inevitable self-dissolution of
consciousness; it corrupts the ethical at its root and opens the path
to the religious for all that is superstitious and absurd. With the
conviction that the absolute heterogeneity Nielsen posits between the
principle of knowledge on the one hand and the principles of faith and
action on the other is ``a misunderstanding that is ruinous for
religiosity, morality, and all genuine character development,'' I have
subjected his theory to a thoroughgoing critique, and sought to
demonstrate its untenability, its deficient grounding, its
self-contradictions, its ruinous consequences.

In publishing this contribution to the investigation of a problem that
through the centuries has had, and will continue to have, the most
decisive significance for the human spirit, it is my hope that it may
occasion a renewed, serious, and conscientious discussion --- a sharp
and open conflict between principles. Such a conflict, even if carried
on with sharp weapons, can only serve the interest of the cause; and
one must therefore hold dear the combatant for whom the cause, and not
some egoistic self-regard, is the highest consideration. A firm and
deep personal conviction I shall always respect; by sound reasons I
shall always be just as open to influence as I intend to be impervious
to anything other than these.


%% ============================================================
%%  HISTORICAL MOMENTS (pp. 8--31)
%% ============================================================
\chapter{Historical Moments Illuminating the Emergence of the Problem}
\label{ch:historical}

\section{The Preparation of the Problem in Antiquity}
\label{sec:antiquity}

It is first among the Greeks that a definite approximation to a
clarification of the problem of the relation between faith and knowledge
comes into view. In the Orient, e.g.\ in India, one becomes aware,
where a thinking relatively independent of religion develops, of the
discrepancy between its results and the doctrines of religion, and an
opposition forms between a heterodox and an orthodox doctrinal edifice;
but the distinctive character of Oriental philosophizing, its lack of
thoroughly achieved independence from religious doctrine, of the
capacity to develop its content in the form of the concept, and of a
clear consciousness of the principle of philosophy, brings it about
that no reflection is reached concerning the principled ground of the
difference---that is, no consciousness of the essential determination
of faith and knowledge. The conscious difference is only the factual
difference of content.

In Greece, where the knowing spirit becomes conscious of itself in
philosophy, and, in unfolding the content of its knowledge, connects it
through a single principle, it can, with this content of knowledge---of
whose principle and whose connection with the principle it is
conscious---enter into relation with religion as it is given in the
forms of popular religion, and the problem is intimated; indeed it even
seems, at the close of classical antiquity, to come forward in its
decisive form. It will, however, become apparent that even where the
question of the relation between the doctrines of philosophy and of
religion, and of religion's significance for the spirit's cognition, is
most strongly accentuated, one does not get beyond intimations, and that
what prevents this is the specific distinctive character of the Greek
spirit.

Already in the earliest period of Greek philosophy we see philosophy
standing in a critical-polemical relation to popular religion and its
myths. Thus the founder of the Eleatic doctrine of unity, Xenophanes,
combats the anthropomorphistic conceptions of the gods in the myths, and
his polemic is repeated in other thinkers of that remote time, but
without our finding in any of them a genuine reflection on the
significance and nature of popular religion. In the culmination period
of Greek philosophy, in Plato and Aristotle, the critical relation to
the mythical emerges in connection with a clearer reflection on
philosophy's position with respect to the myths and cult of popular
religion. Here Aristotle's conception in particular becomes
significant. What is distinctive for him, as in general for Greek
philosophy up to his time, is that philosophical thought, resting and
satisfied in itself, passes judgment upon the positively religious from
out of itself, without entertaining doubt about thought's authority. In
the doctrines of religion Aristotle acknowledges as true only what is
confirmed by philosophy; the mythical clothing, on the other hand, he
denies truth without hesitation, tracing it back to conscious invention
or unconscious anthropomorphizing, and the forms of cult have
significance for him only as what is legally valid in the state.

The situation presents itself differently in the post-Aristotelian
period, and in this period it passes through two essentially different
stages. As Greek philosophy passes beyond its culmination point and
becomes conscious of the inner discord that it cannot reconcile, there
first appears a greater attentiveness to the religious, in that the
relation to it, both the acknowledging and the polemical, acquires a
more essential significance for thought that no longer rests securely in
itself. We thus see in the Stoics an attempt to rationalize the myth,
and we see this attempt bring about a counter-effect of the mythical
upon philosophical theory, a counter-effect that becomes evident in the
Stoics' theodicy and in their doctrine of divination. On the opposite
side we see in the Epicureans a rational opposition to myth, a struggle
to free the spirit from the disturbing influence of the fear of the
gods, in which Epicurus sees the essence of religion; and in the later
Academics, e.g.\ Carneades, as in the later Skeptics and in the
renewal of Cynicism, we find an opposition directed partly against the
theological attempts, partly against popular religion.

A new stage in the development of the relation, and one that in this
connection is of essential interest, begins around the time when
Christianity appears, historically prepared by the development of the
Greek and Oriental spirit, in Neo-Pythagoreanism and in the concluding
forms of the older philosophical schools. The more thought was
exhausted, the more the spirit theoretically and practically lost its
resting point and through skepticism was brought to despair of its own
and of reality's truth, the more was truth---the truth of cognition and
of action---sought outside the human spirit in the transcendent divine,
which through intermediate beings reveals itself to the spirit; and
since such divine revelations were seen in the religions, in them a
starting point for knowledge was sought. And the relation to
revelation---to that occurring through intermediate beings or already
given in the religions---was then posited in Neoplatonism as conditioned
by religious purifications and by religious practices, in general by a
religious life and a religious disposition. Philosophy, which supports
itself on this content of revelation and would draw truth from it, is
then transformed into theology, and in Neoplatonism, which assigns to
philosophy the task of justifying the content of myth, and, in
assimilating the idealized myth, of taking the place of the religions as
a universal religion, we see within paganism the parallel to
Christian-medieval Scholasticism.

At this point one has come as close to posing the problem of the
relation between faith and knowledge in its decisive form as can be
reached among the Greeks at all. Here there has been reflection on the
spirit's---when left to itself---lack of capacity to grasp what is
decisive in the last instance in both practical and theoretical
respects; philosophy, the unfolding of autonomous thought, has been
placed alongside divine revelation, and it has been declared as the
result of this juxtaposition that the relation to the Absolute, to
absolute truth, cannot be reached through knowledge, but only through
the religious life and the religious disposition. Yet the problem in its
sharpness does not emerge even here; one reaches only the intimation,
and the decisive reason is that the Greek spirit apprehends itself
exclusively under the determination of knowledge, that it lets knowledge
constitute the essence of the human being---whereby from the outset the
possibility is cut off of a principle coordinate with knowledge in its
own distinctive character, and still more of a valid opposition to
knowledge; but only through the distinction and the opposition can the
problem be brought forward and the reconciling unity be recognized.

That the Greek spirit becomes conscious of itself only in the
determination of knowledge can be made evident from a double point of
view. If we first fix our gaze on Greek ethics, it is immediately
striking that in it the concept of virtue becomes the central concept,
while the concept of duty recedes. But if we then investigate how virtue
is determined, we find it determined as knowledge; and since for the
Greeks virtue expresses the activity of the essence according to its
undisturbed nature, the proposition that virtue is knowledge becomes
equivalent to the proposition that knowledge is the essential
determination of the human being.

To the same result we are led when we consider the way in which the
greatest Greek thinkers determine the divine. For Plato the highest is
the Idea, more precisely the Idea of the Good, but the Idea is the
objectified and hypostatized thought-content, which receives from Plato
the determinations \textit{nous} and \textit{phronesis}. In Aristotle
the divine is determined as \textit{nous} and more precisely as
\textit{noesis noeseos}, which in eternal, uninterrupted energy relates
thinkingly to itself and from which all determination of will is
excluded. But the essential determination of the divine is the ideal of
the human spirit, and it becomes also for these thinkers the highest
expression of the essential determination of the human spirit. For Plato
the human soul is closely related to the Idea, indissolubly connected
with it, and for Aristotle the human spirit in its highest form as
\textit{nous poietikos} is in unity with the divine spirit itself,
whose essence is exhausted in thought.

How the Greek consciousness, despite the attempts that are made in that
direction, does not definitively get beyond the determination of
knowledge, becomes apparent when we more closely consider the relation
between knowledge and will and between knowledge and faith.

As for the first relation, even where the volitional side of the human
spirit is more strongly accentuated---which already occurs in Aristotle,
and still more in the last period of Greek philosophy---the distinctive
nature of will is never properly recognized, and the attempts to
demonstrate the source of this distinctiveness fail, so that will is
continually reabsorbed by knowledge. Indeed even where, at the close of
Greek philosophy, among the Neoplatonists, it is declared that genuine
cognition is conditioned by the right disposition---where, that is, a
primacy of will over thought is formally asserted---even there the
movement is only a circular movement from knowledge back to knowledge,
and knowledge shows itself in the last instance as the sole original.
The right disposition, namely---that disposition which turns away from
the finite and as such the untrue and vanishing---is immediately
determined by a knowledge of this fundamental determination of the
finite, such that cognition, in being determined by this disposition,
relates only to itself through will, and in its determinateness as
concrete cognition has as its actual condition and presupposition the
abstract cognition of the fundamental relation. The relation of will to
cognition also becomes clear through the abstraction in which will
remains: will as will receives no distinctive positive content; it moves
only in the generality of the negative; its highest goal becomes the
spirit's ecstatic absorption into the indeterminate infinite.

What holds for the relation of cognition to will holds also for its
relation to faith. When we speak of faith, we must make a definite
distinction between the concept's application in the earlier period and
in a later one. In Plato, \textit{pistis} is placed alongside
\textit{doxa} and stands as something lower in opposition to the true
form of knowledge; in Aristotle, \textit{pistis} (in the
\textit{Rhetoric}) acquires the meaning of means of persuasion. In the
meaning of something higher in relation to knowledge, the concept of
faith is first used in the concluding period of Greek philosophy. In
Philo, who is at least as much a Jewish believer as a Greek thinker, the
concept appears in the meaning of the trusting reception and
acknowledgment of the content of revelation, and after him the concept
is preserved as the concept of that which is higher than knowledge. For
Iamblichus the first condition for genuine knowledge of God becomes the
faith that nothing is impossible for God. For Proclus, faith designates
the immediate relation to the divine unity, in which union with it
through the contemplation of its pure being is given. It here forms the
opposition to knowledge, which is always determined by the plurality of
reflection and which is said to be unable to reach the highest. And yet,
despite these determinations, despite the fact that faith, as the
higher, is distinguished from knowledge, faith falls under knowledge;
the movement here too is only a circular movement from knowledge to
knowledge. In the Neoplatonists the indeterminate, infinite unity to
which faith relates is precisely nothing other than the negation of
determinate knowledge posited through the most extreme abstraction by a
knowledge that has come to despair of itself---a negation that exists
only through thought, through knowledge, and that in its negativity is
determined by it. And in Philo the relation of devotion to the
self-revealing divine in its highest and true form becomes the devoted
absorption, occurring in ecstasy, into the abstract unity; but this
indeterminate unity is the last residuum of abstracting thought, an
object of negative knowledge.

The Greek consciousness can thus reach the point of seeing the
difference between the content of knowledge unfolded in the system of
philosophy and the determinations of religion unfolded in myth and cult,
and of becoming conscious of the principle of knowledge; but it cannot,
in its monism of knowledge, recognize the principle of religion in its
distinctive character, through which it could show itself as something
different from or opposed to philosophy and yet valid---just as it
cannot, owing to its deficient conception of the movement of the
concept, of the nature of spiritual development, while making knowledge
its sole point of view, see the principle of religion as the principle
of a subordinate stage in the spirit's teleological development toward
philosophical cognition as the spirit's true and adequate form.

\section{The Emergence of the Problem in Christianity}
\label{sec:christianity}

\noindent Whereas within Greek philosophy---and only within philosophy
could the problem be hinted at among the Greeks---only a preparation of it
was reached, but since the sharply defined form of it cannot be found even
where all the conditions seem to be present and where the decisive point is,
as it were, touched upon, the Christian era brings about the conditions for
its emergence. Yet here too they develop only successively, and in the
development we can distinguish several stages.

\subsection{The New Testament's Pronouncements on the Relation}
\label{ssec:nt}

\noindent What, as the new doctrine is proclaimed, must at once come to the
fore is the opposition between it and that which it is to supersede. This
opposition is an opposition between the content of the Gospel of the coming
of the Kingdom of God, which is proclaimed and appropriated in faith by
Christians, and the conviction of the non-Christian world, which has its
common ground in a mode of thought in which the Greek and the Eastern spirit
met. Inasmuch as that conviction, in its opposition to the Gospel, is
represented most immediately by the cultivated, the knowing, the opposition
becomes an opposition between the world's wisdom and the wisdom of Christ
revealed to the simple (the wisdom of God); but this form of the opposition
too, when ``wisdom'' is taken in an objective sense, points back essentially
to the opposition in content. A new content, formerly hidden from mankind,
is revealed, a counsel of God for which the fullness of time has come: a
μυστήριον. The concept of the mystery points to the new content, which
embraces also what is concealed in the future (thus the future destiny of
the Jewish people, the resurrection and the transformation); but it points
likewise to that which becomes manifest only to those in whom the Spirit
works. Inasmuch as the concept embraces this latter signification, the
ground of the opposition, its principle, is thereby intimated; but since
there can be no question of the believing consciousness, foreign to
knowledge, seeking the ground in a principle of knowledge, the ground of the
opposition between the non-Christian and the Christian consciousness is
sought in a defect of disposition, of will, which, so far as that
consciousness has an acquaintance with the faith-content of the latter,
becomes a lack of will to accept the content of revelation; and this lack of
will is seen in turn as grounded in the ``carnal'' and the ``psychical,''
which form the opposition to the ``pneumatic.'' Behind this determination of
the opposition lies the opposition of principle which, in its connection
with the problem, first comes clearly and sharply forward in
Scholasticism: the opposition between the natural and the spiritualistically
supernatural. The concept of the natural embraces the thought darkened and
obscured by sin, and for this thought the higher supernatural truth is
therefore inaccessible, which becomes accessible to faith only through a
divine revelation.

The Christian consciousness's relation to its object is designated partly as
πίστις and partly as γνῶσις, and γνῶσις is conceived as the consummation of
faith. The view that would set faith and knowledge in unity and maintains
the possibility of a science of faith has leaned upon the determination of
γνῶσις. But in the New Testament the designation γνῶσις still has only a more
indeterminate meaning; it here means an acknowledging cognizance, not a
cognition through thought. γνῶσις stands to πίστις as the more concrete to
the more abstract, as the unfolded to that which is as yet only beginning,
but it falls within the same sphere, does not become a form of cognition, of
knowledge, for its content is that which surpasses knowledge. But in the
designation γνῶσις there lies a possibility for that conception, and this is
seized upon in Gnosticism and in the older Church, in which the spirit, in
its urge for cognition, would extend the domain of cognition to that which
by its essence cannot be cognized, and therefore lays the signification of a
comprehending cognition into the promised γνῶσις.

\subsection{The Patristic Period's Relation to the Problem}
\label{ssec:patristic}

\noindent By the patristic period we understand the Christian antiquity that
was the time of the production of dogma, and as such we distinguish it from
Scholasticism, the time of the Christian Middle Ages, the period of the
elaboration of dogma, of the formal appropriation of dogma.

The patristic period unfolds dogma in essential completeness into a coherent
Christian edifice of doctrine, and it does this by employing the forms of
ancient philosophy, attaching itself from the beginning especially to the
Platonic and Stoic philosophy, later preferentially to the restored
Platonism. By this unfolding of a Christian edifice of doctrine, which as a
totality can be set against the heathen systems, an essential condition is
given for a sharper apprehension of the problem. Among the Church Fathers a
distinction is drawn between the believing cognition, a Christian γνῶσις,
whose relation to πίστις is now more fully determined, and the cognition of
philosophy, the heathen γνῶσις; but reflection on the relation does not lead
to a definitive result, it does not become clear whether, in the transition
from πίστις to γνῶσις, one is to think only of an ordered combination of the
faith-content into a coherent unity, whose point of departure remains
untransformed into cognition, or of an ἐπιστήμη in the stricter
sense.\footnote{Note the vacillation in the application of the concept that
shows itself in Clement of Alexandria. His conception of πίστις points back
to Philo and has affinity with that of the Neoplatonists.} The opposition
therefore remains, even in the patristic period, in spite of the more
developed thought, still essentially an opposition with respect to content,
and even this opposition cannot, because the consciousness of the matter of
principle is unclear, attain its full sharpness.

This is grounded in the patristic period's relation to antiquity. The men
who in the Eastern and Western Church represented the Christian gnosis were,
by the bonds of language and nationality, by the whole inherited culture, so
closely grown together with classical antiquity that the unconscious
influence of its thought and view of life could not be warded off. While,
therefore, proceeding from the immediate religious consciousness, one set the
opposition between the content of the Christian faith and that of heathen
knowledge, the after-effect of the Greek conception of knowledge prevented
the attainment of full clarity with respect to the principle of the
opposition; and during the unfolding of the content of dogma there pressed
in, hidden and irresistibly, an admixture from ancient philosophy and the
ancient view of the world and of life, and it pressed in even into that
which was central for Christianity, as in the doctrine of the God-man and of
sin, and the boundary-determination between the oppositions became uncertain
for reflection.

\subsection{The Development of the Problem in Scholasticism}
\label{ssec:scholasticism}

\noindent Scholasticism has for its presupposition the unfolding of dogma
into an essentially completed whole; its task is, by the aid of a reflecting
thought, to appropriate dogma to consciousness by finding, if not the
\textit{ratio} of dogma, yet the \textit{rationes} of the dogmas. When we are
therefore to designate its age according to its striving, we designate it as
the time of the appropriation of dogma. Just as by this determination it
distinguishes itself from the patristic period, so it distinguishes itself
from the latter by the altered conditions of the time. The dependence on
antiquity that was grounded in the immediate connection of nationality had
ceased, and the production of dogma was so far concluded that a system,
embracing the totality of faith and knowledge, could now be attempted. In the
service of the Church and bowing beneath its authority, thought strives to
erect such a Christian edifice of doctrine, in which the content of faith and
the content of thought are to be fused into a whole. The manner of the
joining becomes clear when we consider the great systems that were produced
in the flourishing period of Scholasticism. In these we see the ancient view
of the world, which gathered itself together in the representation of the
harmoniously ordered divine Κόσμος, as it were embraced by the frame of the
Christian. The view of world-reality, which had to be taken from philosophy
given as a historical presupposition, for which the world was a reality, and
which for antiquity was the whole, now became, conjoined with the Christian
faith-content, a part, fitted into a larger whole. The content of thought
taken over from antiquity filled out only the middle part of that
world-drama, of which the first and last act belonged to Christianity. Before
the natural universe was set the nature-free, almighty, creating God, after
it the Christian heaven, the supernatural kingdom of grace. But in this
composition the beginning and the end became, by Scholasticism's own
admission, something inaccessible to thought, hence in essence heterogeneous
with that which filled out the interval; the specifically Christian content
therefore presents itself, in spite of the attempted joining, as something
incomprehensible to the content of thought. And in seeking the ground of the
incomprehensibility, one finds it in the incongruence between natural thought
and supernatural revelation. Faith's ``higher cognition,'' which owes its
existence to a heightening of the natural faculties brought about by a divine
miraculous working, has only the name in common with the cognition of
knowledge, which rests upon those natural faculties, which do not stand in a
continuous developmental relation to that which is posited by the miraculous
heightening. The significance of natural cognition for the faith-content
becomes only this: in negative form to assert it as the not-impossible, not
positively to ground it. But in spite of the difference of essence, which
leaves behind only the likeness of the name, the relation between faith and
knowledge is nevertheless determined as a relation between two forms of
cognition. Nevertheless there lies in the heterogeneity of the oppositions an
inducement to trace them back to different functions of the spirit, and such
a tracing back we find, for example, in Duns Scotus, when he posits faith as
coming into being only through an act of will and therefore lets theology,
which rests upon faith and whose aim is not most immediately cognition but
blessedness, the enjoyment of God, stand in a closer relation to the will
than to the understanding, inasmuch as it becomes not a speculative but a
practical doctrine that is to lay hold of our will.

When we follow the problem through the time of Scholasticism, we find a
development that is essentially conditioned by that attempt at union, by which
the consciousness of the relation must necessarily be clarified, and the
development becomes partly a development within the same point of view: that
of knowledge, partly a development by which a new point of view for the
opposition is intimated.

In Scotus Erigena, who forms as it were a transitional link between the
Fathers and Scholasticism, whose distinctive character he intimates, while
he, like the former, bound by the relation to antiquity, answers the problem
through a postulate of the harmony between the activities of spirit sprung
from the same source, faith and reason---a harmony in which reason becomes
the overarching and primitive. But in the time of Scholasticism proper the
development proceeds in the direction that the cognition of reason is
subordinated to the cognition of faith, and as the difference posited in the
subordination is sharpened more and more, there is reached at last an
abstract opposition between faith and knowledge, between the truth of
religion and that of science, so that the at first merely quantitative
difference becomes a qualitative one, and finally an irreconcilable
opposition; and as this development approaches its conclusion, it is that
reference is made to a principle which, as principle of faith, is another than
the principle of cognition and of knowledge in general.

In the first period of Scholasticism, Anselm posits faith as point of
departure and norm, but demands a knowledge that takes up the faith-content
into itself, and he ascribes to comprehended faith a higher worth than to the
uncomprehended. Among the great representatives of Scholasticism in its
culminating period, in Thomas Aquinas and Duns Scotus, a higher cognition of
faith is distinguished from the cognition of the reason not illumined by
grace, but Thomas maintains, alongside the necessity of faith, its agreement
with reason, inasmuch as the latter too is a revelation of God, and Duns
Scotus maintains, by his doctrine of the proportionality between the
oppositions, the connection between supernatural and natural cognition, even
though he must admit that we cannot reach a comprehension of faith, which, as
indicated above, he lets have its human origin from an act of will. In the
concluding period of Scholasticism, finally, in consequence of the absolute
opposition between the supernatural and the natural, the difference between
theology's higher cognition, which is given through the faith brought about by
a spiritual creation, and natural knowledge is so sharpened that worldly
science is not even to be able to give an indication of the divine truth, a
preparation for our higher ends. It is now pronounced, for example, by William
of Occam, that something can be true in theology that is false in philosophy,
and herewith Scholasticism has itself pronounced its judgment upon the union
that was its essential task; it has itself sundered the two elements and the
two principles that it would force together.

\section{The Problem's Resolution in Modern Philosophy and Theology}
\label{sec:modern}

\noindent In the Christian Middle Ages, in which philosophy was the handmaid
of theology, thought governed by the authority of the Church, the resolution
of the problem had necessarily to bear essentially the Church's stamp, and
the individual voices that rose in behalf of an independent thought and
asserted the supreme right of reason had to be drowned out and silenced. The
relation shows itself otherwise from the time when Scholasticism itself
dissolved the contradictory unity of the Middle Ages, when philosophy and
theology each received its own truth and the emancipated philosophy now
autonomously unfolded its truth: the truth-content of the human reason
undisturbed in its essence. Against theology's answer to the problem there now
stood philosophy's, and the more the latter gained independence and asserted
itself unassailed by the Church, the more theology's apprehension of the
problem was determined by philosophy's.

Modern philosophy has come into being under an opposition to the unfree,
formal, and reality-less scholastic thought bound by the authority of the
Church, and to the Church's doctrine itself, an opposition that already in the
period of transition developed into a passionate struggle in which philosophy
had its martyrs. Taken as a whole, modern philosophy stands in opposition to
the religious view of life of the Christian Middle Ages. The apprehension of
the essentiality of the natural and of the autonomy of thought and its
validity as a source of truth divides philosophy decisively from
Christianity's view of the world and of life. Yet there arises, in the
relation to Christianity and, in consequence, in the answer to the problem of
the relation of knowledge to faith---for which this determinate form of
religion was essentially taken as representative---a strongly marked
difference between the two main directions into which this philosophical epoch
has developed: the realistic and the idealistic.

The realistic direction had necessarily, in accordance with its conception of
nature, to mark itself off at once decisively from the Christian, and it has
consistently preserved its relation and constantly sharpened it in form,
inasmuch as it passes from a formal accommodation, defensively prescribed by
external circumstances, through an avowed indifferentism, to an ironic mockery
and thence to a sharper polemical relation that gradually assumes the
character of a passionate hatred. While thought, which thus successively gives
nature the preponderance to such a degree that it is itself absorbed by nature
conceived as matter, and, unexplained by it, becomes a riddle to itself,
judges, as thought, in the name of reason, religion to be the irrational, the
unconditionally pernicious for man's progress and man's happiness; the source
of religion is determined as ignorance, faith becomes superstition,
irrationality, even with an admixture of the morally reprehensible, and the
authority to judge is awarded solely to that thought which pronounces the
immutable laws of nature, of moved matter, of the One-Existing, and asserts
them in their application to the human, which becomes only a dependent member
in the all-embracing nature.

In comparison with the thoroughgoing uniformity of realistic philosophy's
position toward the present problem, the idealistic line of development
presents a far greater manifoldness of forms of the relation. The consistent,
conscious carrying through of the whole epoch's fundamental thought had
necessarily, within this line too, everywhere to bring the opposition to
light; but the abstract likeness that obtains between idealism---particularly
where it develops as panlogism---and the Christian view of the world, and the
fact that, with a single exception, all the thinkers who represent idealism
were personally attached to Christianity, has been able to call forth a
striving to set religion and philosophy, faith and knowledge, in harmony, and
this striving comes forward in various nuances.

The opposition, which throughout the whole of modern philosophy is
apprehended essentially from the point of view of knowledge, of cognition, is
determined either as an opposition between the true cognition, that of reason,
and an untrue cognition, that of imagination (Spinoza: \textit{intellectus}
--- \textit{imaginatio}), an opposition that essentially coincides with the
opposition between the determination of infinity and that of finitude, between
the eternal idea and the vanishing historical form (Schelling's earlier
philosophy);---or as the opposition between the goal and truth of cognition
and a stage of development in cognition that is as yet only imperfect and
relatively true, and that in the teleological development is taken up into the
perfect truth, the truth of philosophy (Hegel, when the Logic's determinations
of the relation between form and content are consistently carried through). To
this latter form of opposition are to be traced back, in principle, the
attempts that have been made to bring the religious doctrines, by a
reinterpretation---moral or speculative---under the domain of the
philosophical cognition of reason; for while in these very attempts there lies
the presupposition of the teleological developmental connection, there lies
therein at the same time an acknowledgment of the necessity of raising that
which appears in merely relative truth to the form of the pure absolute truth.
That, individually, the consciousness of this connection may be more obscure,
that the reinterpretation may be conceived as touching only something
unessential, an external difference of form, makes, in principle, no
difference to the matter.

The endeavor to bring faith and knowledge into harmony, which comes forward
within certain forms of development of modern philosophy, assumes---when we
leave aside the superficial notion of a ``religion of reason,'' a ``natural
religion,'' of which Christianity is supposed to be the proper
expression---either the form that a quantitative difference is conceded in
favor of revelation (a supra-rational that is yet not to be a
contra-rational), or the form that an essence-, content-untouching difference
of form is posited in favor of philosophy. The latter nuance was represented
chiefly by the so-called right side of the Hegelian school, which leaned upon
Hegel's philosophy of religion, whose restoring tendency, at least in
expression, had led to the relation between form and content developed in his
Logic becoming obscure in its application to religion's relation to
philosophy. From that school this mediating conception was then carried over
onto theology's domain in the speculative dogmatics.

The intimations of a collision-preventing distinction between the domains of
faith and knowledge, which in this period are given within philosophy, attain
no further significance. When Spinoza gives religion a practical task, and by
this determination would ward off the collision with philosophy, it is an
accommodating turn that cannot be maintained over against his philosophy's
determinations of the relation between the theoretical and the practical. And
when Leibniz somewhere intimates the same distinction, his utterances lose
their significance alongside the general determinations of his doctrine, and
alongside his endeavor to assert the theoretical truth of the dogmas.

Within the domain of theology the sharply marked opposition between faith and
knowledge comes forward essentially in the period of the Reformation and
particularly in Luther. Thought's domain is, according to his conception, the
worldly, but in the domain of the religious it can avail nothing without
denying itself. He does indeed sometimes seek to assert for thought a
significance as negatively determining, inasmuch as he ascribes to it the
capacity to decide what the divine is not; but he cannot consistently carry
this concession through. Among the old Lutheran dogmaticians the opposition is
already softened. During theology's later development the influence of the
mediating philosophy makes itself felt, directly or indirectly, and the
difference becomes, where it is not wholly leveled out (as in the
rationalistic identification of natural religion and Christianity), a
quantitative difference, even where the designations sound qualitative---one
that in reality is in principle excluded by theology's claim to be a science.
The attempt made by Schleiermacher to assert, in accordance with what had
already earlier been intimated by Lessing, under the opposition between
religion and philosophy, the validity of religion by giving each of the
opposites its own peculiar domain, by setting each of them in relation to a
peculiar faculty of spirit, has only a subordinate significance in comparison
with the attempt at distinction made later by Kierkegaard. We rest this
judgment on its significance partly on Schleiermacher's defective conception
and arbitrary delimitation of knowledge, partly, and chiefly, on the unclear
and contradictory in his position generally, inasmuch as he, on the basis of
an---albeit inconsistent---pantheistic philosophical standpoint, would assert
the validity of the positively religious, particularly of the Christian, and
out of the contentless and purely negative conception of the idea of the
Absolute conditioned by his philosophical standpoint, determines the essence
of religion and the organ of religion, whereby it must precisely become
impossible to reach the concrete dogma and the historical that is
indispensable for positive religion.

Against the striving, represented by the Hegelian philosophy of religion and a
part of his school, to maintain the essential unity of religion and philosophy
under the difference of form apprehended as essential, an objection was raised
both from the side of religion and from the side of philosophy. From the
former a protest was lodged against the dissipation of that which is peculiar
to religion through its transference into philosophy's concepts and into the
form of the conceptual process; from the latter was asserted, in accordance
with the Hegelian logical development of the concept, the inseparability of
the difference of form from a corresponding difference of content. The
recognition of the impracticability of the attempt at mediation then led to
the attempt to solve the problem through a distinction of spheres. In Germany
Feuerbach makes this attempt in the interest of knowledge. The character of his
attempted solution has been indicated in the foregoing in its essential
fundamental features; a critique of it will, partly directly, partly
indirectly, be given in the work that is to follow upon this one. Among us the
problem entered a new phase through Kierkegaard, by the attempted solution
which, while like Feuerbach's it rests upon a distinction of spheres, aims at
the union of the principles in the same consciousness through the recognition
of both.

From this attempted solution derives all that has a real significance in the
Nielsenian, and this relation between them can already here, at the close of
this brief survey of the essential stages in the history of the problem, be
made clear by a preliminary juxtaposition of the two with respect to the
points that become of significance in this place.

In Kierkegaard the religious is the original and moving personal interest and
the goal toward which from the beginning aim is taken; but in the development
of the standpoint knowledge is the point of departure, and it is through
knowledge's self-delimitation, through the recognition of the barrier that the
relation of existence sets for the existing knower, that the necessity of
faith and of the relation of faith is acknowledged. The validity of knowledge
is then not annihilated, but is acknowledged as relative, within its
determinate boundary; faith becomes, as the highest form of certainty for the
existing one, the absolutely valid, that which alone gives the existing one
the reconciling certainty, as though wrested out of uncertainty, of that which
satisfies the deepest urge of personality. Knowledge can, precisely on account
of its delimitation, not disturb the content of faith, nor comprehendingly
appropriate it; the only thing it can do in relation to faith is to become
serving, inasmuch as it guards the paradox, preserves it as such in its
sharpness and inaccessibility. Conversely faith lets knowledge remain
undisturbed in its delimited domain, and in accordance with this relation the
two principles with their spheres can find their place in the same
consciousness without disturbing collision.

In Nielsen too knowledge becomes the point of departure of the development, in
him too knowledge delimits itself, so that, through the recognition of its not
merely relative but absolute boundary, it is led to acknowledge in its
validity the faith-principle absolutely heterogeneous with itself. Knowledge
clarifies the relation between the faith-principle and the knowledge-principle
and their spheres, and it is capable of this inasmuch as it, reflectingly,
embraces its opposite, and so makes that too its object, not only that which it
grounds. Each of the spheres that are determined by the principles absolutely
heterogeneous with one another unfolds in its inner infinity without outwardly
restricting the other and therefore without coming into conflict with the
other; and thus, precisely on account of the heterogeneity, the union of the
spheres and their principles in the same consciousness does not divide the
unity of consciousness.

In the decisive fundamental features the two forms of the theory thus agree.
Were a difference to be sought, in the determinations indicated, it would be
this, that Nielsen would set both spheres as infinite and unrestricted, both
principles, in accordance with their absolute heterogeneity, as absolutely
valid, and that he would posit a relation of coordination, whereas Kierkegaard
clearly pronounces the merely relative validity of knowledge and the finitude
of its sphere, and definitely sets knowledge in a relation of subordination to
faith. But on closer consideration such a difference would show itself as
illusory.

As regards first the infinity of the spheres, which would have to follow
immediately as a consequence of the absolute validity of the principles, and
which would have to imply that all reality, reality in its completeness, found
its place within each of the two spheres, it cannot be maintained. With
respect to the sphere of knowledge, which seems to receive such an infinity
through the fact that knowledge also has for its object that which it merely
reflects, the infinity would in any case become an infinity with respect to
extent, not to content. By the very separation between that which is
assimilated---that grounded and comprehended by thought is assimilated---and
that which is merely reflected, there is set in the content of knowledge a
boundary that points to an essential delimitation in knowledge as human
knowledge, even though this delimitation is postulated as abolished in the idea
of knowledge. But the essential delimitation, which comes forward immediately
through the delimitation of the content of knowledge, must inevitably lead back
again to a delimitation with respect to extent, on account of the dialectical
relation between extent and content and their necessary reciprocal determining.
It is immediately evident that that which is merely reflected in knowledge---not
only when knowledge is taken in the strict sense as appropriating (assimilating)
knowledge, and only in this sense is it properly knowledge---is known only
according to its abstract being; but that, even when knowledge is taken in the
extended meaning, it too is known only abstractly, because concrete knowledge
exists only through the cognition of the connection, but for reflecting
knowledge the connection of the reflected must withdraw itself, inasmuch as its
form---on account of the absolute heterogeneity in the principles---becomes
absolutely incommensurable with that of knowledge. And what holds with respect
to the infinity of the sphere of knowledge has its counterpart with respect to
that of the sphere of faith. This too cannot be maintained in the indicated
meaning. Within faith there is found no relation that corresponds to the
relation of reflection in knowledge and in accordance with which the content of
knowledge could be taken up into it. And even if such a relation obtained,
faith would not have all reality for its object, for the corresponding reason to
that for which knowledge did not have it. The sphere of faith can thus not
assert its infinity. Only this result is sought here; in what follows it will be
shown in its place that the factual delimitedness of the spheres, however much
the self-delimitation is accentuated, will yet lead to a collision at the
boundary with the sphere that is touched therein.

How the delimitation in both spheres is to be conceived is easy to see when we
take a concrete object. If we take thus the natural, then the law of nature and
the causal connection determined by the law of nature, which are both
determinations of knowledge and as essential objects of knowledge have validity
as reality, do not exist for faith, which sees things only as created, in their
dependence on God. Conversely knowledge does not cognize the form of connection
of the creaturely relation of dependence, which, as essential object of faith,
according to the positing of form, would likewise have to have validity as
reality. The total reality, whose being in knowledge and faith was to give them
infinity, is thus reduced for each to an incomplete reality, for that which is
object for knowledge and faith is only the connectionless things, things as it
were ``an sich,'' but not the relation in which they stand to their ground, not
the connection with this and with the homogeneous reality, hence not that which
makes things what they are; and this connection, which determines the concrete
reality, belongs only to knowledge or to faith, and falls outside the sphere of
the common reality, but as a reality; for only in it are things what they are.
But this really means: in reality the reality is doubled and the doubled reality
is divided between faith and knowledge, so that the apparently common becomes
only a homonym. And this doubling of reality will show itself as contradictory
and impracticable.

As regards the second point emphasized above: the coordination of the spheres,
this is in Nielsen, in spite of the determination of the principles as
absolutely heterogeneous, only apparent; in reality there obtains in him, as in
Kierkegaard, a subordination, so that the inequality between the theory's two
forms of presentation here too disappears. The subordination of the principle of
knowledge and the sphere of knowledge under that of faith is in Nielsen not only
indirectly posited in several ways: by the determination of the will, which is
brought into special relation to faith, as the properly constitutive in the
human essence; by the conception of the essential determinations of theoretical
and practical cognition and their relation to personality, and by the conception
of ``the religious ethics' superiority over the rational''; but it is expressly
pronounced in the designation of the truths of religion as supra-scientific
truths.

In the points that have an essential significance for the development of the
problem, the Nielsenian theory thus coincides with Kierkegaard's, and the
originality belongs to the latter alone. Yet there arises, in the critique of
the theory, occasion to keep its forms of presentation apart from one another in
order to illuminate the additions it has received in Nielsen through his attempt
to set the fundamental determination of the relation between faith and knowledge
in connection with his psychological and metaphysical theory, and the
contradictions and inconsistencies from which it suffers in him, but which
Kierkegaard knew how to avoid.

\begin{center}
---
\end{center}

After what has been developed in the foregoing, the essential standpoints that
have been taken up in the answering of the problem are accordingly the
following:

\begin{enumerate}[label=\arabic*)]
\item The religious standpoint, which sets faith, as the relation of
  appropriation to the divine revelation conditioned by God's supernatural
  working, as absolutely raised above natural knowledge, which is conceived as
  that which not only cannot cognize the content of revelation, but must deny
  it.
\item The standpoint of theology and of a single philosophical form, which
  posits a quantitative difference in favor of faith\footnote{To this standpoint
  the scholastic distinction of \textit{duæ veritates} in reality goes back.}.
\item The standpoint of pure rationalism, on which the difference is set forth
  by the notion of a religion of reason, a natural religion, with which
  Christianity is identified.
\item The standpoint of mediation within idealistic philosophy, on which a
  difference of form is posited in favor of philosophy.
\item The standpoint of realistic and of consistent idealistic philosophy, on
  which, the relation being determined from the point of view of knowledge,
  knowledge receives an absolute preponderance over faith, and becomes the sole
  organ of truth.
\item The distinction of the spheres in order to set aside the religious.
\item The distinction of the spheres in order to bring about the union of both
  in the same consciousness.
\end{enumerate}

Of these standpoints, as already earlier indicated, the theological and that
which through the distinction of the spheres aims at their reconciliation will
become the object of a separate, detailed critique; while the others, so far as
it is at all necessary, will be illuminated through the critique of those, or
through the developments in the work that forms the continuation and conclusion
of this one.


%% ============================================================
%%  PRELIMINARY CONCEPTUAL DETERMINATIONS (pp. 31--43)
%% ============================================================
\chapter{Preliminary Conceptual Determinations Illuminating the Problem}
\label{ch:concepts}

In raising the question of faith's relation to knowledge, we take faith not in
the vague, effaced meaning---as immediate knowledge, immediate certainty of
being---in which the word is used by Jacobi and by theologians who, in the
interest of ambiguity, adopt his indeterminacy, but we take the word in a
specifically religious meaning, in the meaning of the faith that relates itself
to religion, more determinately to positive religion. The determination of the
concept of faith therefore becomes inseparable from the concept of religion.

This latter concept we apprehend first in its ideal meaning, in the meaning of
that which is the moving power in the development of the phenomenal forms of the
religious, the goal toward which the historical development of the religions
strives. We reach the concept of religion according to its ideal meaning by
taking our point of departure from what is the aim of the religious relation.
This we can determine in a single word as the reconciliation of the spirit, its
reconciliation in itself and with its principle, the ground of its essence. The
human spirit seeks in religion liberation, peace, rest, reconciliation; it seeks
it for itself as a totality, according to the full determinateness of its
essence; only thus is it reconciled in itself. The religious relation must
therefore be a relation in which the human spirit, according to all its
essential fundamental activities, comes to rest, a rest that includes
reconciliation and satisfaction for cognition and feeling as for the will. It is
the undivided human spirit's urge that leads to the religious relation, and when
this urge, while reconciliation is sought, comes forward most nearly and most
immediately under the form of a single fundamental activity, then the
satisfaction of this activity's urge, which here as it were represents the
striving of the other activities, must go in unity with the satisfaction of
their urge too: the reconciliation of the whole human spirit must be present
under the reconciliation of the single function of spirit. But the full
reconciliation, the reconciliation of the spirit determined according to its
totality with itself and with its own and existence's principle, can take place
only through that relation to the principle by which the spirit, in
consciousness of its essential unity with this, in free surrender sets it as the
principle of its \emph{whole} being. This is to say, in other words: the
religious relation must be determined as the \emph{absolute} relation to the
Absolute. This determination includes a threefold. Partly it points back to the
fact that the Absolute must be in essential unity with our essence; for only to
that with which we are in essential unity, to our Absolute, can we stand in an
absolute relation, in a relation embracing at once our whole cognizing and our
whole life. Partly it points to the fact that that to which we stand in the
religious relation must be the truly infinite; for only in the infinite can the
human spirit absolutely find rest. Partly, finally, there lies therein the
consequence that the religious relation, precisely as the absolute, cannot fall
outside the relative and particular relations, whereby it would itself have to
become relative, but must enclose these relations in its concrete, embracing
unity.

But when the concept of the religious relation is thus determined according to
its ideal and universal meaning, there shows itself an incongruence between this
concept and the positive, historical religions stamped as particular forms of
the life of spirit. Through all of these there passes, as the moving power of
the development, that striving of the human spirit, called forth by the
fundamental relation to the principle, after the truly religious, i.e. the
absolute relation embracing the whole human essence and human life; it is this
striving that, during the historical development of the religious forms, leads
them forward from the lowest forms of the nature-religions to the highest of the
spirit-religions; but through all the historical, positive religious forms the
ideal concept of religion comes forward in a determinateness of relativity and
finitude: the positive religions show themselves as \emph{finite} expressions of
the essence of the religious. The finitude becomes recognizable through the
forms of the human spirit's activity in the religions, through the conception of
the divine conditioned thereby, and through the determination of the human's
relation of interaction with this.

We seek first the determinateness of finitude in the \emph{theoretical} form of
the religious consciousness. This is, at the stage of the positive religions,
even though the religious consciousness as human consciousness is worked through
by thought, not that of the comprehending thought developed to clarity and
coherence, infinite according to its essence, but essentially that of immediate
feeling and of representation, and from this the determinateness of finitude is
inseparable. But the peculiarity of the form of consciousness must necessarily
determine the peculiarity of the content of consciousness. The spiritual content
can indeed have an imperfect form, inadequate to the content according to its
essential determinateness, without the content being transformed for
consciousness, when it has already been present for consciousness in the
adequate form, and from this has consciously been carried over into that. The
form then becomes only a clothing, symbol, allegory. But in the development of
mankind the form of religion constantly precedes the comprehending cognition:
the Absolute is present in its forms before it is present for consciousness in
the form of thought. It is thus not set over from a more adequate form into the
representational form of religion: this is for consciousness not a clothing, but
the \emph{only} form for the content; it counts therefore as the true and
adequate form. The content can thus not remain untouched by the peculiarity of
the form: it exists only as the thus-formed content for consciousness, and a
separation of it from this form becomes impossible at the original religious
standpoint. Such a separation takes place only through an awakening reflection,
and this comes forward at the religious standpoint only in relation to other,
most nearly to subordinate, religious standpoints. In relation to its own content
the religious consciousness is, by the nature of the case, not critical: it
finds itself again in the thus-formed content.

The specific form of the religious consciousness thus acquires an essential
significance, and here particularly the effect of the representational form is
to be emphasized, inasmuch as the content of the religious consciousness, which
in the immediacy of feeling is in a closed-together, intensive, but unfolded
unity, is first unfolded into determinateness through representation, becomes a
concrete object for consciousness.

The peculiarity of representation is determined by its psychological position
between sensation and thought, between which it forms the middle term.
Representation has, as middle term between the two, a moment of both: it combines
the receptive dependence of sensation with the freedom of thought, its
singularity with the universality of the concept. That which is represented is
loosened from the immediately existing external, set over from the external space
into the inner space of consciousness; it has begun to be distinguished in
relative universality from external singularity and its contingency; but
according to content it is still determined thereby. From this it follows that
that which is in representation is touched by the form of \emph{sensibility}, that
accordingly, when the Absolute is to be apprehended through representation, its
universal essence must be intuited in the form of a single existence determined
by a reflex of sensibility, its eternal life-process, embracing and idealizing
the development in time, as a process determined by the finitude of time, whose
single moments fall separated outside and after one another; that the forms of
the pictorial, of the mythical and the historical, must become inseparable from
the apprehension of the Absolute.

While the form of feeling made the content finite by letting it remain without
unfolding, the form of representation thus brings, through its relation to
externality, a determinateness of finitude into the religious consciousness's
apprehension of the Absolute. It comes forward in that form of distinction by
which the transcendence of the divine is determined; in the apprehension of God
under the form of personality, in which the peculiarity of the sensibly
determined single personality (the individual) shines through; in the historical
in the religions, in which the divine receives a finite limitation of time and
space, and in the apprehension of God's working as a working through the single,
separated acts of the miracle\footnote{The miracle, the specific mark of positive
religion, is not the highest expression of omnipotence, but the making-finite of
omnipotence, the expression of an illusory infinity of omnipotence instead of the
true infinity of omnipotence, hence in reality the negation of omnipotence.
Omnipotence is made finite precisely by the negation of the law of nature, which
in the miracle is set as the expression of a working of the almighty will's
other, which is opposed to the working of omnipotence.}, the discontinuous
expression of arbitrariness.

To the determinateness of finitude in the religious consciousness, apprehended
from its theoretical side, there corresponds, and by it is conditioned, a similar
determinateness of finitude in it, seen from the \emph{practical} side. The
relation of thought to the object is a relation of freedom and infinity. The
individual is free in the relation, inasmuch as it, through the object, in which
thought as comprehending finds itself again, relates itself to itself, and in
this relation it has an inner infinity. And as free it lets the object subsist in
its freedom, inasmuch as it, when it dissolves the form of being of the object
into thought, when it seeks itself in the object, does not seek and find again a
subjective, individual thought, but the objective, universally valid thought,
which constitutes at once the essence of the object and of the individual. It is
thus, precisely because the relation as a relation of freedom is a relation of
infinity, uninterested, unegoistic in the relation to the object. Of another kind
becomes the relation where representation is the form of consciousness. The
representing subject is in a finitude-distinction from its object conditioned by
the nature of representation; it therefore asserts itself according to its
singularity, does not yet meet with the object in the universal, does not stand
in that relation of freedom and infinity to it. Applied to the relation to the
Absolute, this entails the following consequence. Inasmuch as the believing
individual relates itself to the Absolute apprehended in the form of
representation, it stands as a finitely determined single personality over against
a divine single personality, from which in reality the determinations of finitude
are not removed. But this relation, in which the finitude is abiding, must
precisely for that reason acquire the character of open or hidden egoism. Even if
the relation is designated as a relation of love, it holds that love's surrender
is only absolutely free of egoism where the love relates itself to the Absolute
cognized in its true infinity\footnote{cf. S. Kierkegaard:
\textit{Kjærlighedens Gjerninger}.}, and only where love's surrender is
absolutely free of egoism is love the true love and reconciliation through love,
the true unity with God, possible through it. Where the determinateness of
finitude abides, this unity cannot be reached, the reconciliation cannot be
apprehended as the individual's reconciliation in its ground of essence, but it
must be apprehended as a reconciliation with another being, which by its form of
singularity preserves the distinction from man. This transcendent being's demand,
its commandment, then becomes determining for the religious relation by which the
reconciliation is to be brought about; to the individual's fulfilling or
not-fulfilling of the demand there attaches a reward and punishment, which, like
the commandment, is determined not out of a relation of essence, but by the
unconditioned divine will, and therefore becomes a motive of egoism; and even
where the representation of blessedness most nearly touches the thought of unity
with God, a determination of finitude presses into it, partly directly in the
apprehension of blessedness as reward, partly indirectly through the
representation of time, through the representation of the futurity of blessedness,
through the distinction of the consequence of the action from the action itself.
On the whole the abiding determinateness of finitude shows itself in this, that
the religious relation becomes a merely particular relation; for this is one with
its becoming a finite relation.

If we, instead of taking the point of departure from the apprehension of the
religious consciousness's theoretical determinateness and from this drawing the
consequences for its practical relation, wished immediately to hold to the
religious consciousness's practical determinateness, and to assert the view that
the point of departure had to be taken from this, because the will was the
decisive thing for the religious relation, this would yet not lead to any other
result. For we would then, inasmuch as that emphasis on the will rests upon an
opposition of principle between it and cognition, receive as the determining
thing the will in its isolation; but the will in its isolation is, as will later
be shown in detail, only unclarified, egoistic drive, hence a form of finitude
of the will\footnote{However much the will's autonomy and independence of
principle are accentuated, its making-finite is yet not avoided when it is
apprehended in isolation, for precisely the abstract distinction between the
universality-form of knowledge and the singularity-form of the will makes both
finite.}, even though it be adorned with the highest names. Where the will, in
postulated autonomy, is set as the constitutive in the human essence, and as that
which in its distinction is to express the highest spiritual, it precisely reaches
no further than to the imperfect form of the will, which forms the parallel to the
form of representation in the theoretical domain. Only in unity with the true
thinking does the will reach its truth. Where it unconsciously follows its own
law, it is blind drive, and where it is determined by a conscious motive of will
not raised to thought (``a will-cognition''), it always has in itself a
determinateness of affect; it is determined by craving and passion, and this
practical determinateness of finitude inevitably reacts upon consciousness's
determining of the object, which shows itself especially clearly in the praying
person's apprehension of the divine, which is to be moved by the prayer. The form
of affect that comes forward in religion is essentially fear and hope, and the
fact that fear and hope in religion are determined by the representation of a
\emph{superhuman} being (an expression applicable to all the most diverse
religious forms) does not free them from unclarity and egoism, and it does not do
so because the divine, apprehended through a form that is not that of the concept,
has in itself the determinateness of finitude.

For the determination of the concept of faith in relation to the phenomenal
religious forms the presuppositions are now given. Inasmuch as the divine acquires
the character of a transcendent and supernatural, of something distinguished from
the human essence, the relation of communication of the divine thus determined to
the human spirit is apprehended as a relation of revelation, represented in such a
way that revelation is not an inner communication of essence proceeding from the
essential unity with free necessity, but an arbitrary communication, conditioned
by external means, of a content that, as communication about God's essence, is
limited by God's arbitrariness, and, as communication about God's commandment,
which is to determine human action, has its validity through the will of the
subject of revelation, bound by no necessity of essence. Revelation, taken in this
meaning, demands its own peculiar instruments, and it grounds a \emph{positive}
religion, which contains a law for cognition and action.

These determinations with respect to the concepts: revelation and transcendent,
supernatural subject of revelation, have their validity not only for the
spirit-religions, but also for the nature-religions, even though, by the nature of
the case, they come forward less clearly in these. In the nature-religions too,
through the involuntary personification of the nature-power, a transcendence is
brought about, so that the nature-power becomes not only a \emph{superhuman}
power, i.e. a power higher in degree than the human, but in a certain sense a
\emph{supernatural} power, i.e. a power that in its hypostatization is formally
loosened from the real natural and by fantasy is freed from nature's, i.e. from
its own barriers; a power that, through fantasy determined by fear and hope,
transcends itself, abolishes its own conformity to law---which, at this standpoint,
means: the conformity to law that consciousness here apprehends: the relative
conformity to law of the customary. From this subjectivized transcendent there is
seen also in the nature-religions a revelation as proceeding, either through the
stirrings of the religious in the single individuals, or through peculiar
individuals as media, or through the supposed signs and workings of the divine.

When religion as positive, historical religion thus becomes equal to revealed
religion in the meaning indicated above, and is constituted by the relation,
through a content of revelation, to this content's transcendent-supernatural
source, then faith as religious faith acquires hereby its closer determination.
Faith as religious faith thus becomes the relation of acknowledgment---the
theoretical and practical---to revelation, and the relation of acknowledgment
included therein to the subject of revelation.

In this determination of the essence of \emph{faith}, which apprehends it from the
standpoint of the positive religions, there lies implicitly how this faith's
relation to knowledge and to the action resting upon knowledge must be determined
from the standpoint of the religions and from that of science. Inasmuch as the
religious consciousness standing at the stage of positive religion apprehends its
content as \emph{divine} revelation, it must, when it is to determine the
relation---as we have seen above with respect to Christianity---set the
faith-content as incommensurable with the human power of knowledge, and this
incommensurability acquires its more determinate stamp by being traced back to the
opposition between infinite and finite, between supernatural and natural, between
divine wisdom and human cunning, between the doctrine of the Holy Spirit and the
contrivance of the human spirit corrupted by sin. For the distinction between the
divine and the human immediately set by the form of the religious consciousness
must, inasmuch as reflection, determined by the representational form, apprehends
it, acquire the shape of opposition; and inasmuch as faith now acquires a divine
object, and likewise---since the object is subject of revelation and revelation is
that which conditions faith---a divine principle and origin, it must necessarily,
when the relation is to be determined, set itself in a relation of opposition to
knowledge with its finite content and natural principle. The form of transcendence
and of single personality, in which the divine is apprehended, includes at once
that the content of revelation, as having the divine for its object, cannot be
communicated to men except in accordance with a free divine creative decision, and
that it must be a higher than human thought.

Between faith and knowledge there thus forms itself an opposition, unabolishable in
accordance with the separation between the divine and the human, with respect to
the cognition of the divine essence; but this opposition embraces likewise the
opposition between revealed religion's supernatural law of action and the law of
action that is determined by the natural cognition of the human spirit's essence
and essential relations.

The difference between faith and knowledge would also in reality be determined as
an opposition if the relation between revelation and reason were determined by the
apparently quantitative difference that revelation's content was to have a wider
extent than knowledge's, was to surpass knowledge, and \emph{abidingly} to be a
supra-rational, without its yet being expressly determined as a contra-rational.
For that which distinguishes itself from knowledge and within the sphere of
cognition---and thus the relation is apprehended here---will limit knowledge,
becomes knowledge's opposite and in reality abolishes knowledge. For knowledge is
knowledge only as that which according to its concept is closed in itself, which
holds the whole of existence within itself, its principle is a universal principle.
Were a more embracing than the sphere of knowledge set, then the latter, as going
into the more embracing concatenation---the infinite concatenation of the ``divine
truths''---would have to be determined as an organic member in that; but from this
it would follow that it was first apprehended according to its truth when it was
apprehended as such, and that knowledge's apprehension of its content would be
false, because it apprehends that which is only a member as a whole, and otherwise
it cannot cognize it, because revelation's supra-rational content cannot become a
content of knowledge. Only faith's cognition would then be the true cognition,
knowledge's would be not only incomplete, but essentially untrue.

Apprehended from the standpoint of the believing consciousness, the relation
between faith and knowledge will thus show itself as a relation of opposition, but
in the same way the relation must be determined when it is seen from the standpoint
of knowledge. Inasmuch as science cognizes the limitation of finitude, conditioned
by the positive religions' relation to representation, in their apprehension of the
divine and its communication to the spirit, and vindicates for itself as its
specific form the infinity-asserting, comprehending cognition, it sets the
opposition as an opposition between faith's ``cognition'' and knowledge's
cognition, inasmuch as, even when the religious is essentially apprehended from the
side of the will, the will-form is seen as standing in an indissoluble connection
with a determinate form of cognition, so that the moment of cognition accordingly
becomes essential also from this standpoint.

The given conceptual determinations thus lead us to the preliminary result with
respect to the relation between faith in the meaning of positive religion and
knowledge, that the difference, both from the standpoint of faith and from that of
knowledge, must be determined as an essential and unabolishable opposition, which
must lead to the necessity of a choice between them, and that the unity of
knowledge and faith becomes possible only when the limitation of finitude is
removed from faith and religion, and these are apprehended according to their
ideal determination of essence\footnote{How the relation between science and
religion according to the latter's ideal concept, between knowledge and the faith
corresponding to that concept of religion, is to be apprehended, will be developed
in the work that is to follow upon this one.}.

A departure from this determination of the relation is now given at the two
standpoints that will particularly become the object of a critical examination.
From the side of theology, while a positive religion is maintained as truth's
absolute form, the difference is, through the claim to the possibility of the
union of faith and knowledge in a science of faith, set as a merely quantitative
one. From the standpoint that distinguishes the opposites in order through the
distinction to reconcile them, the opposition between faith and knowledge is
conceded, but against the assertion that a choice between them is necessary because
the kind of opposition makes impossible their union in the same consciousness, the
assertion of the possibility of such a union is set, while at the same time
objection is made to religion's truth being judged from the standpoint of knowledge
or subjected to science's judgment, inasmuch as it is asserted that its truth rests
upon another principle than that of knowledge and follows its own, higher law.

Through the critique of these two standpoints' answer to the problem the validity
of the determination given above will then be demonstrated.


%% ============================================================
%%  THEOLOGY'S ANSWER TO THE PROBLEM (pp. 43--118)
%% ============================================================
\chapter{Theology's Answer to the Problem}
\label{ch:theology}

\section{Theology's Task as a Science of Faith}
\label{sec:science-of-faith}

Theology's general position toward the problem---a position that is here
accordingly distinguished definitely from that of the immediate religious
consciousness---is to be more closely determined in the following way. Theology
as a science of faith becomes precisely the expression of an essential belonging
together of faith and cognition. It does not indeed limit faith to cognition, but
sets cognition as an essential determination in faith. It does not deny the
difference between faith and knowledge, both when knowledge is apprehended as
natural knowledge and when it is apprehended as the knowledge, determined by
faith, that would comprehend revelation, but it accentuates the possibility and
necessity of giving faith an expression in the form of science. It thus acquires
the peculiar task: inasmuch as it sets revelation as a transcendent being's free
communication, which no human cunning could discover and which none can dissolve
into concepts, and demands that human knowledge give up its autonomy in favor of
revelation's authority, of determining revelation and God's essence and working
in such a way that they are drawn out of the purely incomprehensible, and on the
other hand of letting the concept of knowledge be so transformed, be it by a
``heightening'' or by an extension, that it can give place to that which yet can
never and shall never in the proper sense be comprehended, and of making possible
a science that follows another law than that of autonomous knowledge, also
embraces that which is unappropriable for this knowledge, and yet does not cease
to be science.

The possibility of this task's solution is what the critique of theology is to
investigate. This critique will, as earlier indicated, attach itself closely to a
determinate theological attempt that in our literature has acquired a peculiar
significance and among us can be regarded as the proper representative of the
theological endeavor; but inasmuch as the critique attaches itself most closely
to this, it will everywhere trace its treatment of the decisive points back to
that which is typical for theology as theology and show it to be expressive of
this, so that the critique, through the single theological attempt, relates
itself to theology in general.

We first present in coherence this theological theory, which has been developed
by Bishop Martensen in his various writings, in its essential fundamental
features, in order thereby to indicate the points of attachment for the critique.

The task of the science of faith is this: to unite faith and knowledge in an
organic unity, so that on the one hand faith is asserted in its primacy, on the
other hand its essential connection with knowledge comes forward. Faith has the
primacy because it, as the immediate, is the first, while knowledge, as the
mediate, is the derived, and because faith is far richer than knowledge: a new
life in the soul, a relation of existence, and can never be exhausted by the
comprehending knowledge; but its unity with knowledge must be asserted because
faith, which in itself embraces will, feeling, and cognition, is itself in
knowledge: the central knowledge, so that knowledge becomes the common
determination for science and faith. Precisely this relation makes a
``Christian,'' ``religious,'' ``believing'' science possible. That faith has the
determination of knowledge in itself follows from the fact that God has revealed
himself to faith, and that his love and grace have made themselves known under
the form of wisdom, as also from the fact that the believer, in virtue of the
Christian idea of truth contained in the living faith, must be able to say: I
know what I believe. Were faith not already in itself a knowledge, then it would
not be the faith of \emph{truth}, and the great opposition between truth and lie,
pronounced on every page of Holy Scripture, would be without significance. Were
faith not itself a knowledge, then the believer could not confess his faith and
could not reject and combat the perverse and erring doctrines. This knowledge in
religion itself is what was designated as the primitive, the central knowledge,
and, where it comes forward with special vigor, it can be designated as the
central intuition, which has in itself the drive to unfold itself into concept.
From this proceeds the scientific, the reflecting and speculating (intuiting in
the mirror of the idea and seeking possibilities) cognition.

In accordance with this relation between faith and knowledge, the science of
faith must stand in a relation of dependence, but in a free relation of
dependence, to revelation. Christianity's absolute truth must be acknowledged as
given prior to and independent of all speculation; cognition shall be a cognition
in faith and out of faith (\textit{credo ut intelligam}). The holy thought of
wisdom must in the believing consciousness become the principle of thought, and
human thought is able only through it to search out the depths of revelation, to
scrutinize for connection and ground in the Christian representations, to strive
to bring forth a spiritual counter-image of the eternal wisdom of revelation.

The Christian science of faith would present the Christian intuition as a
doctrinal concept coherent in itself. It is not content with an explicative
comprehension, but seeks, in accordance with an inner drive, a speculative one,
which does not remain standing at presenting the connection in the given, but
asks about possibility and ground. Every «\textit{Ita}» contains a hidden
«\textit{Quare}», which under the thorough explication cannot but come forward and
call for a higher speculative comprehension. Yet the science of faith, since its
object is not a mere truth of reason, but a positive truth, a truth of faith,
cannot develop its content or demonstrate its truths with the mathematical and
logical necessity peculiar to the merely rational science: there arises a
boundary, not indeed for the certainty, but for the clarity; but this is not a
defect in the science of faith; it points precisely to its higher worth.

This becomes clear both when we look to what is this science's proper organ in the
human spirit, and when we look to its objective principle, to the concept of God
that determines it.

As regards the first point, faith is determined as a will's appropriation in love
of the divine revelation. The will is accordingly that function of the human
spirit to which it has its essential relation, but the will is, just as it is the
first in the universe, the properly constitutive in the human personality: man's
fundamental being is will. But thus the will becomes the highest, that which has
the primacy in relation to the other functions of spirit, which are dependent on
it. That what is grasped in faith through the will is not commensurable with
thought, with the concept, therefore shows only that the former is the higher and
richer.

And to the same result we are led when we consider the science of faith's
objective principle. The principle of the science of faith is a higher than that
to which the merely logical thinking attains; it is not ``reason's dead idol,''
but an ethical principle, a real principle of personality, a supernatural
principle. The ethical concept of God is the higher, which does not exclude, but
includes, the logical and physical. It is the concept of the free, willing,
personal God, of the not merely logical but holy, absolute truth. The personal
God, who is not the idea-bound, but the idea-free, just as he is the nature-free,
reveals as creating and miracle-working his freedom and omnipotence, and in his
revelation through history the same freedom of personality and of will makes
itself felt, in opposition to nature's and the merely logical's necessity, and it
shows itself in its highest form where God not only reveals himself through his
workings, but reveals himself in his personal appearance in the historical Christ.

With this concept of God there enters another concept of science than where there
is question only of a merely logical and physical knowledge: there is reached the
concept of a positive philosophy, the concept of a science of revelation, in which
theology becomes the central science. Of the science of revelation in general
holds that which was said above of theology, that it does not demonstrate its
truths with logical and geometrical necessity, and that it can bring it only to a
relative comprehension of the highest truths; but this is precisely the mark of
the higher and real science, of the higher cognition of reality. The science of
revelation and particularly theology must, by the nature of the case, be richer
in intuition than in concept, but it is intuition that gives the dialectic its
worth; for intuition is the organ for the whole and reaches further and deeper
into the very innermost of things, to their first beginnings and roots, than the
merely dianoetic thinking. The science of faith is therefore science in the
preeminent sense, while the natural sciences can in no way in the stricter sense
be called sciences, and while the mere philosophy of reason, the negative
philosophy, is only propaedeutic, a child's lesson in comparison with that higher
wisdom. The latter's peculiar scientific character and higher worth lie precisely
in the relation to the free, personal God and his revelation, to the holy will
that is the measure of what things in the last meaning are. Just as revelation
has proceeded from the divine will's freedom, not governed by an abstract logical
or physical necessity, so too its apprehension is conditioned by a free act of
will, which determines thought, is not determined by it. The divine revelation
can, as a revelation of freedom, a revelation of personality, not be apprehended
in the forms of logical necessity, not be dissolved into mere conceptual
determinations; it can, as a communication of life, not be taken up through the
forms of thought alone. In all revelation of freedom there is something
incommensurable with the purely logical thinking, and therefore this thinking
inevitably lacks the eye for the historical that characterizes the higher science.
But when in the cognition of the revelation of freedom one goes beyond the domain
of the merely logical, we are yet not led out onto the domain of the lawless, for
the divine freedom is no groundless arbitrariness; it contains a higher necessity,
but one that lets itself be cognized only voluntarily. The free divine omnipotence
is, as creating and miracle-working, determined by love and wisdom, hence by an
ethical necessity of freedom, and necessity in the general meaning acquires again
its relative validity, inasmuch as a physical becomes the condition for the
ethical will in God, and inasmuch as God is indeed nature-free, but not
nature-less, inasmuch as there is in God an eternal physis that stands in
essential relation to the will. In this relation there is also given a point of
attachment for thought, by which the spiritual God's creating of a material
nature, and his immediate miraculous working, even though the how of creation and
of the miracle eternally withdraws itself from human comprehension, is yet raised
above the purely incomprehensible, in a certain measure becomes the object of a
cognizing.

Thus, then, the possibility of the science of faith is grounded in a higher
cognition determined by faith and will, a cognition for which revelation's truth
opens itself so far as the divine can be grasped by the human spirit's frailty,
while the knowledge that by the human spirit's own power and own means, by mere
reason, would comprehend the divine revelation, must distort and deny its truth.
The higher cognition is the consummation, the redemption of mankind's general
cognition of reason, a higher potency of reason's revelation. Revelation in Christ
abolishes the laws of reason in their abstraction in order to confirm them in the
Christ-wisdom that is the fullness of the law of reason. And the higher cognition
includes in itself the norm and the corrective for the lower; by it the higher
science is the ruling science, and thus the right relation is brought about in the
world of thought: οὐκ ἀγαθὸν πολυκοιρανίη, εἷς κοίρανος ἔστω.

\section{The Relation of Will to Knowledge and to the Essence
  of the Human}
\label{sec:will-knowledge}

What is justified in this theory, that which has an unconditional claim to
acknowledgment, is the endeavor to bring the human spirit's highest forms of
expression, its most essential fundamental relations, into inner harmony. The goal
toward which there is striving is the spirit's highest: the spirit's true
reconciliation. Even though, therefore, the means by which what is aimed at is
sought to be reached are insufficient, even though the way that is followed cannot
lead to the goal, this theory will yet constantly be able to make claim to a
respectful, unprejudiced, and penetrating critique, all the more as it comes
forward in a form that is not only fine and full of spirit, but bears the
unmistakable stamp of personal conviction.

\begin{center}
---
\end{center}

The critique of the theory must from the beginning be directed against a point
that is the whole theory's scientific center of gravity; against its determination
of the will's relation to knowledge and to the essence in the human and in the
divine.

For on this determination, according to which the will is set as the primitive,
the essence-grounding in the human and in the divine personality, and as including
in itself knowledge as a governed moment, there rests on the one hand the
apprehension of a science that goes beyond the boundary of merely logical
knowledge, and yet not only, by the moment of knowledge that the will includes,
asserts for itself the character of science, but precisely by its overstepping of
that boundary acquires a higher worth;---and on the other hand there rests
thereon the apprehension of God as in his free willing raised above knowledge's
abstract necessity and above the essence-determined freedom that is only a
\textit{libera necessitas}, and yet as in his being and working remaining the
object of a relative knowledge. When the will's primacy is asserted, then
knowledge's necessity is governed, its binding law subjected to the subject's
infinite freedom and its individual urge, then the relation of cognition is
dependent on a relation of will and the will is entitled, as choosing the truth,
to correct cognition's consequences. When the will is set as the primitive in God,
then God is the creating and miracle-working, the freely self-revealing, and with
this concept of God the concept of the demanded higher science, of its coming into
being through a free communication from its principle, is immediately given.

The point of departure for the critical investigation must be taken in the
psychological, whose apprehension becomes in general decisive for the
determination of faith's and knowledge's relation. In the apprehension of the
psychological, the development in Martensen attaches itself most closely to the
single philosophers who stand near to theology. It is particularly the not very
clear thinker J. H. Fichte from whom the determination of the psychological
fundamental relations is drawn, and what he has set forth in the introduction to
his still unfinished Psychology is reproduced in part with his own words. It is
said\footnote{Martensen: Om Tro og Viden pp. 45 f. 123 f. 129.} to be the result
of the deeper psychological observation, a result in which J. Böhme, Fr. Baader,
Schelling, and the younger Fichte agree, that not thought or reason, but the will
is the first, innermost, most original, and most central in the human soul, that
both our theoretical and our practical consciousness has its psychological root in
a pre-conscious will, which manifests itself in these two main forms, that
accordingly self-consciousness and thought itself is a form of the will, namely
the form of mirroring, in which the willing being objectifies itself, itself and
its world. But inasmuch as self-consciousness and thinking are in themselves
nothing other than the mere act of mirroring, they presuppose a really willing
subject that mirrors itself therein. This real subject, whose fundamental being is
will, is the substantial, consciousness is only quality, predicate, attribute for
it: it is that through which the real being comes to clarity about itself and its
relations, that through which the pre-conscious, potential will actualizes itself.
---All observation is also to lead us back to desire, drive, and instinct as the
first psychological, hence to an obscure will that includes reason and cognition,
as that out of which consciousness develops. And if we ask about the fundamental
relation of cognition and will in the actual consciousness, it is not cognition
either, but the ethical will that has the primacy; not the ethical will as a
merely practical, cognition-less will, but the ethical will that has cognition,
the contemplative, as its own moment. And as in man the will is the first, not
thought, so it is also in objective existence, in the universe, and thus, as
theism precisely develops, will, not knowledge, is the first in God, and the
ethical idea, the good, which is the will's idea, is the highest, the absolute
idea, to which nothing is coordinated, but all subordinated.

The untenable in the reported psychological theory shows itself as soon as the
main points in it are more closely illuminated.

1.~Inasmuch as the will is set as the primitive, there is a going back to the
concept of a ``pre-conscious will'' (Fichte's ``vorbewußter Wille''), which is to
be that out of which the spiritual life's various forms develop. But what is called
``pre-conscious will'' is, precisely as pre-conscious, i.e. as not yet will, not
will; for will is---as Martensen himself concedes by letting the will be
actualized through self-consciousness---only as conscious will. Therefore that will
is also to be only an ``obscure'' will. But when the concept of will in the
stricter sense becomes inapplicable, the question must be raised: what it is in the
determination will that becomes applicable to that primitive working. The answer
must be: it is the determination of energy; but energy as a subject's, a self's
energy we designate as selfhood-energy. Energy is not a predicate of a presupposed
subject, for thus apprehended the subject would become the energy-less, its
connection with the predicate contradictory; energy is an essence-determination,
expression for the very essence as ideal-real\footnote{The ``real, willing
being'' is in Martensen and Fichte a contradictory reminiscence of Herbart's
immutable, inactive, and unconscious Real. ``Reality'' here becomes for Martensen
a meaningless determination; it could at most designate the monadic, but thereby,
in this connection, nothing would be gained.}. But the essence-energy, the
selfhood-energy, comes forward first through the organic forming; through this it
forms the conditions for the will, but just as fully and just as originally the
conditions for cognition: the selfhood-energy therefore stands in its first form
of appearance no closer to the former than to the latter, and inasmuch as the will
and cognition are realized, energy becomes the determination for both, the concept
embracing both, and we speak with the same right of an energy of cognition as of an
energy of will; in cognition and will energy realizes its inner possibility in two
mutually supplementing fundamental forms. That one has transferred the will's
determination to energy and designated it as pre-conscious will, also at the stage
where the will would be only physically (organically) working, derives from a
misunderstanding of the will's relation to energy and to real working. While one on
the one hand overlooks cognition's determination as energy and the cognitive
activity's unconscious working upon the real in the organism, one overlooks on the
other hand that inasmuch as the will as will exercises a conscious working upon the
real, this its activity presupposes the selfhood-energy's organ-forming activity,
which continues also after the will as will has come to appearance, is independent
of the will, and precisely thereby distinctly separates itself from it. Inasmuch as
one calls the original energy will, one slips in the representation of a being
working as it were behind our essence, whose activity is represented in analogy
with the conscious will's working upon an external through the organism.

2.~If there thus takes place a confusion of concepts in the determination of the
primitive in the \emph{essence}, then the determination of that which for
observation is the psychologically first suffers from a similar confusion. The
first for observation becomes, when one thereby understands the first psychological
functions that go beyond the merely organic forming, not desire and drive. Instinct
one could at a pinch use as designation for this first, but instinct would then,
just as in animal life, from which the designation is taken, embrace intelligence's
and will's most rudimentary working, the unconscious and half-conscious functions
of cognition, and the unconscious and half-conscious functions of will. The first
psychological is that to which desire and drive point back. Desire and drive are a
striving to hold fast or to remove a determinate state; but this state is only a
state of the self as felt, through the organic sensation or the sense-sensation.
The psychologically first is accordingly sensation: cognition's rudimentary form,
just as drive is the will's rudimentary form. This relation is conceded by
Martensen in so far as it is said that in the relation of reflection, i.e. in the
relation to the objective, it is to be correct to set the will as the second,
cognition as the first. But in the relation to the subject itself the will is to be
the first. This distinction can, however, not be maintained. The self's relation to
itself is constantly a relation to itself as a member in an existence, as
determined by the objective. Even in sensation's most primitive form as organic
sensation, in which the self relates itself to its own side of reality, the relation
to the objective, which is set as such in the sense-sensation, is implicitly
present. But the organic sensation is the first form of the self's being-for-itself.

3.~The apprehension of self-consciousness and thought as a mere reflection, which
is asserted by Martensen, is impracticable. They are to be the real, willing
being's mirroring, in which it objectifies itself. But mirroring expresses
reflection. In consciousness's mirroring the real being is set as reflected in
itself. But such a reflection is the absolutely inexplicable when reflection is not
originally set as an essence-determination. Without this the essence would not
become clear to itself in reflection, would not objectify itself for itself. The
reflection would merely become a reflection in another, apprehended by a supposed
outside-standing consciousness. The essence would become the obscure-to-itself in
spite of the supervening predicate, which as this subject's predicate would become
a purely incomprehensible, and the willing that expressed the essence would, when
reflection's determination inseparable from the actual will was taken away, be only
a working similar to one material's mechanical working upon another material. If, on
the contrary, reflection is set as an original essence-determination, then the
essence is originally determined by reflection, i.e. unconscious thinking, by the
universal self's being-for-itself. That apprehension of thought and consciousness
as the originally reflection-less being's mirroring in itself would likewise make it
absolutely inexplicable how a consciousness of another and a comprehension of this
other can take place. Only where reflection becomes an original determination is
there given, in the self's doubling, in the determination of otherness thereby set
within the self, the connection to a relation to the otherness outside the self.

4.~By the determination Martensen gives of the psychological fundamental relation,
he cannot avoid entangling himself in contradictions. It is a contradiction when
there is talk of a common root for the will and cognition, when they are accordingly
coordinated, and the will is yet set as the higher unity. It is likewise a
contradiction when the will, which is to be the primitive, as will is to be
actualized by self-consciousness; for that by which the actualization takes place
must be the first, just as it must be a substantial. It is a contradiction that the
will, not reason, is to be the first; for thus the will as reason-less would have to
ground reason. It is finally a contradiction that Martensen pronounces that
knowledge and will cannot be separated, that the one cannot be without the other (p.
40), and that he everywhere urges an \emph{essential} significance for the
theoretical alongside the practical, but yet would deprive thought and
self-consciousness of substantiality and reduce them to being merely quality and
reflection.

Against the psychological determinations given by Martensen we then set, supported
on the foregoing developments, the following:

The fundamental determination of the human essence is a selfhood-energy that
reveals itself in knowledge and will.

Knowledge and will both stand in that relation to the essence, to the substantial,
that they cannot be determined as qualities of a presupposed ``real'' essence, but
must be apprehended as forms of being for the subjective essence, as substantial
forms of essence.

In the development of the human soul's life the psychologically first is
cognition's first rudimentary form.

In accordance with that which marks the human essence's peculiarity in relation to
the animal's: the essence's universality, the form of subjective being-for-itself
conditioned by universality, which, determined by the universal relation to itself,
comes forward in the energy of self-consciousness and thought, becomes the form of
being principal according to concept for the human essence.

The significance and power that are ascribed by Martensen to the will, in
accordance with its posited primacy, in relation to cognition (as governing
thinking and as choosing truth), cannot be maintained; it rests essentially on a
confusion of concepts, partly of the will and the energy of cognition, partly of
the will and the human essence determined as self by energy.

From these psychological determinations it will then follow as a consequence that
there cannot belong to the will, inasmuch as it determines itself as ethical, an
autonomy and a primacy in the meaning in which Martensen ascribes it to it. The
will becomes ethical will first as determined by cognition: the cognition of the
good is the condition for willing it\footnote{cf. Martensen's statement (Om Tro og
Viden p. 26): the existing good is the condition for our willing of the good, for
our striving after it.}.

The ethical will is an essence-willing, but the essence is principally the being
disposed to self-consciousness, thinking according to its concept: an
essence-willing is therefore determined by cognition, by self-conscious thought.

\section{Will and Knowledge in God}
\label{sec:will-god}

The same fundamental relation that was posited with respect to the human essence
is, as already stated, also posited with respect to the divine. The concept of God
to which philosophy attains through the cognition of nature and reason, and which
expresses the absolute principle of reason and of the thought-determined reality,
is not, even though it determines the Absolute as absolute subjectivity, to be the
true concept of God, not to be the expression for the actual God, and it is not so
because God is in it determined by knowledge, and along a logical way. In his
truth and reality God is to be apprehended only when he is apprehended as the God
whose essence is will, and when the ethical concept of God resting upon this
apprehension, reached through a cognition determined by the will, which determines
him as the \emph{Good}, as the personal God, supersedes the
logical-physical\footnote{By the logical concept of God Martensen understands (cf.
Om Tro og Viden, p. 11) ``the concept of God as impersonal idea, all-intelligence,
Νοῦς, ontological subjectivity, or the like''; by the physical ``the concept of
God as an all-producing nature, \textit{natura naturans}, substance, etc.''; by the
ethical: ``the concept of God as personality and will, as the God of the good, of
love, of grace, and the God of the not merely logical but holy truth''.}, which it
is yet not to exclude, but to preserve in itself as a moment. The ethical concept
of God, which, in opposition to the logical-physical, apprehends God as a
supernatural principle, sets determinations of will, such as: absolute freedom,
holiness, love, as this supernatural principle of spirit's determinations; it
determines God as the personal creator-God. In this determinateness God becomes
the object of the higher science, which Martensen gives the name of science of
revelation or philosophy of religion.

As regards this apprehension of the will's relation to God's essence and the
concept of God conditioned thereby, it is now at once evident that the apparent
grounding by which, for man's part, the will was to be asserted as the principal,
and which without further ado seems to be supposed to hold for the divine, can
acquire no application to the latter, onto which the human's relations could not
immediately be transferred. The obscure, pre-conscious will, which in man is set as
presupposition for the conscious will and for thinking, acquires no place in the
divine essence; for even though there is talk of a physis, a ground of nature, in
God, God is yet not to be thought as coming-to-be out of this obscure,
consciousness-less ground, but as eternally being in his free knowledge of
himself, in the knowledge that has its reality in the Trinity. Knowledge and will
have then the same eternity in God, nor can there, when this relation is regarded,
be ascribed to the will a primacy of concept; for the will in its clarity would
precisely have God's knowledge of himself as conceptual presupposition, and if a
higher third were to be sought, it would be to be sought, not in the will, but in
the divine self's energy, by which it eternally sets its being and sets it as
determined by omniscience and omnipotence.

If, then, the will's primacy in the divine is to be asserted, it must happen along
another way: it must be asserted by showing that inasmuch as the will \emph{as}
will, that is, in its original autonomy, not as determined by cognition, is the
principle of the ethical, it as such must be the highest, because the ethical,
which expresses freedom, is the higher in relation to the logical and physical,
which do not reach beyond necessity.

\section{Critique of the Concept of God}
\label{sec:god-concept}

That the will as will was the principle of the ethical would be demonstrated if it
were shown that the ethical was higher than the logical in such a way that it could
not be taken up by the latter, but yet included it, in the same way as it is to be
higher than the physical, which it is to include as a moment. But a development by
which this was shown is not given in Martensen, although he in many turns of phrase
pronounces that the ethical concept of God is the highest and the absolutely true;
and by the determinations he gives, the postulated relation transforms itself
arbitrarily: that which was postulated to be a higher goes back to that in
comparison with which it was to be the higher, indeed, it even goes down below this.
And this is in him nothing accidental, but is grounded in the nature of the case.
This is here to be demonstrated with respect to the ethical's relation to the
logical; it will later be demonstrated with respect to its relation to the physical.

\subsection{The Relation of the Ethical Concept of God to the Logical}
\label{ssec:ethical-logical}

When the ethical concept of God is set by Martensen as higher than the logical and
physical, this is not to be apprehended as if it comprehended the logical and the
physical in a logical-physical unity that as such becomes the higher, but in such a
way that it, inasmuch as it includes in itself the logical-physical, is a higher
than this very unity. But from this follows the immediate consequence, which is
really only a tautological paraphrase of the relation, that the essence of the
ethical cannot be determined by an immanent development of thought. Were the essence
and relation of the ethical determined by such a one, it would itself go in under
the logical and become only its highest concept; logic would, as metaphysics, also
be the ethical's metaphysics, include the fundamental concepts of ethics, naturally
in logical determinateness of form.

In order to maintain the presupposition, the ethical must be set as transcending the
logical according to its \emph{whole} extent, and this could be apprehended either
in such a way that another relation to reality was posited in the ethical than in
the logical, or in such a way that the ethical was given a content richer than that
of the logical.

To the first apprehension we are pointed when Martensen sets the relation to
existence as the decisive thing. The ethical is to be a higher than the logical
because it points to existence, while the logical is to be indifferent to existence,
to have only ideal, not factual being. But this opposition and these determinations
rest on an unclarity. The logical, as subjective, formal logical, has only a
``thought-being.'' The logical, apprehended in unity with the ontological, expresses
that which has ideal-real being. ``Factual existence,'' which, in opposition to this
being, would be only external, finite being, it does not have, but it grounds that
which exists in this form. And the being of the ethical is precisely a being of the
same kind as that of the logical-metaphysical; the ethical concept of God has as
concept the same sphere as the logical. Were the existence of the ethical of another
kind, it would be only that finite, sensible existence, and it will later show
itself, in the development of the concept of personality, that Martensen precisely
cannot free himself from the representation of this. But if now the being of the
ethical is of the same kind as that of the logical, then there could accordingly
not, through the relation to the latter, be secured for the will the demanded
primacy.

As regards the second apprehension, it will not be able to be maintained, neither
when we look to the relation that must be set in the eternal principle between its
forms of activity, nor when we look to the relation of human cognition. For if we
look to the first-named relation, it lies in the concept of the eternally actual God
that there must be set in God an eternal unity and congruence of knowledge and
will---both concepts apprehended thus as they can be transferred to the
Absolute---, that nothing can be set by God's will that is not transparent in his
knowledge, which must be thought as that which makes his will manifest, while this
has his knowledge for its content. And if we look to the second relation, it shows
itself that the fact that the ethical is determined as richer with respect to
content must lead to the ethical's supra-logical content having to be set as that
which can only be grasped through another function of spirit than the logical: by a
``free thinking,'' as Schelling expresses it; by ``a knowledge that has gone through
the will,'' as Baader designates it; by a willing knowledge, a knowledge chosen by
the will, as Martensen would say. But the possibility of this relation cannot be
demonstrated. To point to the will as to a principle absolutely heterogeneous with
knowledge, in the way Nielsen apprehends it, is cut off by the presuppositions. Were
that relation therefore to be made thinkable, it would have to happen in such a way
that the logical thinking at once determined its own boundary and pointed out beyond
it, and at the same time indicated the organ that was to supersede it, and stated a
negative norm with respect to that which fell outside the boundary, a norm that
secured against the arbitrary and absurd being set as a higher truth. But such a
self-delimitation of thought by thought itself is the impossible. As expression for
the essence-energy, thought is a reality, every one of its activities is a
confirmation, a corroboration of this reality, a self-confirmation: an activity by
which it was to negate itself is therefore a self-contradictory. The relation would
then have to be determined in such a way that not thought as such, but the thinking
subject out of its totality of essence, embracing also other functions, delimited
thought. But from the psychological developments given in the foregoing it follows
that this subject in its unity must enclose the fundamental functions as harmoniously
united, as determined by the same goal, and that knowledge, cognition, as a whole, is
a fundamental function, and that which in the human spirit has ideal priority. Nor
along this way, then, could that posited relation be asserted. Were one to let the
ethical as the higher come forward by a break, and yet at the same time hold it fast
as that which included knowledge, then the knowledge that the ethical will included
would be only homonymous with the knowledge that preceded the break, and is knowledge
in the strict sense. The relation would then in reality be the relation between two
absolutely heterogeneous principles, which as such could not reasonably be set as
standing in the relation of super- and subordination to one another that Martensen
posits.

It thus shows itself to be not accidental that the postulated relation is not
demonstrated as the true one: a demonstration of it is an impossibility. And if we
were to suppose that the will's primacy in the indicated meaning were asserted, it
would imply that our knowledge's necessity was not only suspended for the higher
domain's part, but in general made impossible. Were the will asserted as the
essence-grounding, then the concept of essence would be absorbed in that of the will,
the necessity of essence would then be abolished in the will's freedom, which, as
transcending the necessity of essence, would become \emph{absolute} arbitrariness
(«\textit{bonum est quia deus vult}»); knowledge would in God, as expression for the
will preceding the essence, determining the essence, only express the absolute
arbitrariness's consciousness of itself and of that which is arbitrarily set by the
arbitrariness; and in man it would in the same way lack every ground of necessity.
And when thus the will's primacy abolished knowledge's necessity, this would imply
that in reality it could never become an object of knowledge, when knowledge is taken
in the proper meaning.

When, for the indicated reasons, it is impracticable to furnish a proof for the set-up
relation, then it becomes inevitable that the relation transforms itself where there
is made, most nearly in the form of thetically set-up propositions, an attempt to give
closer determinations with respect to the ethical concept of God's higher worth than
that of the logical. The essence with its necessity, which becomes knowledge's
determination, acquires involuntarily the primacy in relation to the will. The
demonstration of this becomes so far the less essential, as the developments given by
Martensen could suffer from accidental defects; but it is yet of interest, after the
general relation has been demonstrated above, to see how the natural relation
involuntarily makes itself felt, in spite of the point of departure's conflict
therewith.

We here emphasize those determinations of the relation between the ethical and the
logical, between the will and knowledge, that come forward inasmuch as the relation
between the Good (one) and the good is to be indicated, and we emphasize them also for
the reason that they include moments for the apprehension of the concept of the
freedom that was to raise the will above knowledge. The relation between the good and
the Good points back to those relations, inasmuch as the good, as ``the eternally
existing good,'' points to the theoretical, the Good, in accordance with the indicated
apprehension of the constitutive, to the will. The development of this relation is in
Martensen vacillating and uncertain, and leads precisely to the opposite result of the
intended. It is said\footnote{Martensen: Om Tro og Viden, p. 67 f.} that the good---which
is demanded with necessity---can only exist as a free will, and therefore must be
determined as the Good. ``Now,'' it is continued, ``God as the Good is doubtless the
one existing with necessity; but were he only this, then we had not yet come out beyond
the physical and logical; for a God who is only good with necessity (the good nature) is
not the in-truth ethical God. But the good as his essence's necessity demands precisely
that he with freedom realizes it, or that he, as the one who does not have a law outside
himself or above himself, but who is to himself the holy law, as the absolutely
autonomous, actualizes, brings forth himself as the Good, and in so far as he brings
forth the good outside himself, does not do it by natural necessity, but by freedom's
necessity.'' ---Precisely in the given determinations lies that, when there is to be
talk of a law at all, the essence-determination and the essence-necessity becomes the
first, the presupposition for the will; accordingly the essence becomes the holy law:
the good becomes law for the Good. And when there is talk of an actualizing, then there
lies in the concept of actualization a working, hence an activity not yet determined as
``the Good's'' will, hence an energy determined by the essence-necessity, which acquires
its expression in knowledge, by which the ethical will is, and which, inasmuch as it is
the ideally grounding, becomes the principal. And this essence-necessary working can,
since God is not to be coming-to-be, and since he is eternally knowing, not be brought
by Martensen in under ``the obscure will's'' determination.

To the same result lead also the subsequent developments, in which the ethical concept
of God is to be shown as the truth of the logical and physical. As the absolutely
willing, God is at the same time to be the absolutely knowing, who knows himself in an
eternal ideal-system. But the ethical God is not to be the in-the-idea-included,
idea-bound, and so to speak in-the-idea-absorbed God. His knowledge is to determine
itself as wisdom, which is precisely to designate the absolutely idea-free in his
existence in opposition to heathendom's idea-bound, merely logical God, and to designate
his knowledge not merely as a logical, but as an ethical, a willing knowledge. ---For
when God is set as the \emph{absolutely} knowing, he accordingly also knows himself in
the determination of the Good, and inasmuch as in God there can be no talk of a merely
reflecting or of an inadequate knowledge, the divine will goes in under knowledge's
necessity, is seen as the form for the essence's \textit{libera necessitas}. And when
knowledge is determined as wisdom, then wisdom's \emph{content} and wisdom's \emph{goal}
lead back to the same overarching of knowledge, however much there is talk of the
difference between eternal ideas and eternal counsels. The determination of God's
relation to the idea becomes very unclear and the idea raises itself, against what is
intended, as the highest. The idea expresses a logical, but yet not only are the ideas
to have their principle of unity in God, although the absolute idea, which yet would
surely have to be the ideas' principle of unity, is not to be God; but an idea: the
will's idea, is to be the expression for the ethically determined God, accordingly for
that which is above the logical, and the absolute personality is apprehended as the
highest \emph{ontological} determination. The will's idea is to be that under which all
is subordinated. It is to include the determination of existence, accordingly precisely
that determination which is to characterize the actual God in distinction from the
logical concept of God.

The result of what has been developed is accordingly: that the attempts to determine
the relation between will and knowledge in the divine essence through the relation
between the ethical and the logical must lead back to the unity of the will and its
freedom with knowledge and its necessity, so that the logical comes to embrace also the
concept of the ethical God, and that what has been set forth about the will's primacy
and what stands in connection therewith is not demonstrated and cannot be demonstrated.

\subsection{The Freedom of the Divine Will}
\label{ssec:divine-freedom}

What has been developed with respect to the relation between will and knowledge in
connection with the relation between the ethical and the logical can acquire a new
illumination from the fact that the set-up concept of the will as principle of
freedom is analyzed, and thereby the posited freedom's kind and relation to
knowledge's necessity is clarified, since it will thereby show itself why this
freedom cannot assert for the will the demanded precedence in relation to that
determined by necessity. The result of this analysis is in essentials anticipated in
the determinations won through the foregoing developments.

We here, as in the earlier developments, leave aside the question: how far the
transference of the psychological determination will onto God is at all justified,
and whether this concept must not necessarily in this transference onto the Absolute
undergo a transformation; a question to which we come back in the development of the
concept of God's personality. We here raise, inasmuch as we apply the concepts will
and knowledge to the divine, the question: how would that freedom be determined by
which the will was to be asserted as the principal, as the superordinate in relation
to knowledge and its necessity? The freedom that was to raise the will to this
significance would have to be a freedom loosened from the determination of necessity,
for this determination, as determination by freedom, would lead it back to the
essence's and knowledge's in their ground identical necessity. But such a freedom
would, as the \emph{absolutely} necessity-less freedom, be a \textit{liberium
arbitrium} that has no ground in the essence, and therefore becomes only groundless
arbitrariness. That it is in fact apprehended thus will become clear when the doctrine
of creation and the doctrine of the miracles are analyzed. But a freedom of this kind,
which moreover, by annihilating all of knowledge's necessity, would annihilate
knowledge and become a ``non-intelligent'' freedom, would not be able to assert for the
will a higher validity. Just as it is itself a lower than the freedom determined by the
necessity of reason, than that essence's \textit{libera necessitas} which expresses the
self-transparent essence's law of freedom and reveals itself equally in knowledge and
the will, so it would make the will, not the highest, but a lower; it would set a
sublimate of human arbitrariness, a mere fantasy-image of the will, in place of the
will's true concept: to be essence-willing, expression for the total essence's energy,
in which knowledge and will are in unity, and which therefore can express itself only
law-determinedly, in free necessity.

Through that concept of freedom the will, which is designated as ``the only principle
that can begin anything,'' is set not only as knowledge's, but as the essence's
\textit{prius}; in its aseity it at once sets the essence and transcends
it\footnote{The distinction of will and essence becomes clear in the doctrine of
creation out of nothing. Inasmuch as God creates by the will out of nothing, he does
not create out of the essence, and the distinction between will and essence is set. The
will sets being, sets the essence, but transcends being and essence.}. Instead of being
an utterance of essence, bound by the essence's law of freedom and the essence's
content, it accordingly becomes an unlimited will borne by nothing, which proteus-like
transforms itself and can will everything possible\footnote{This finds its expression in
the representation of an ``all-possibility,'' and of God as ``the lord of
possibilities.''} i.e. all \emph{the possibilities of fantasy}; for this will is in its
lawlessness precisely determined only by the barrierless fantasy's lawless omnipotence.
Were the will itself set as essence, this would either lead to the will, against the
presupposition, acquiring the determination of necessity in itself, or---when the
freedom in that meaning was maintained---to freedom, that is, arbitrariness, in its
unboundness by all conditions, hence also by that set by itself, negating this in every
utterance of freedom, hence negating the will identical with the essence, whose
determination it itself is, so that the will accordingly, while it was to be an
essential, i.e. an abiding, unabolishable, would only be in the incessant change that is
only self-abolition.

Through that concept of freedom, which stands in close connection with the ``ethical''
concept of God, it thus becomes impossible to assert the will's precedence in the
divine. We are again pointed to the necessity of setting will and knowledge in eternal
unity and congruence in the divine, of setting the ethical and logical in unity, and it
becomes evident that the endeavors to assert for the will that higher worth, inasmuch as
they must aim at excluding all determination of necessity from it, must precisely
acquire the opposite result of the intended, namely: to draw the will and that which is
called the ethical and freedom in God down below that which was to be the subordinate:
below knowledge, below the logical and below the essence-necessity. But the name of the
ethical and of freedom is also only used here by a misuse; that which is so named does
not even reach to the boundary of the ethical, does not touch the true concept of
freedom, in general not the sphere of the concept.

\subsection{God's Personality}
\label{ssec:gods-personality}

What has been said concerning the apprehension of the concepts treated is confirmed
by the apprehension of God's personality determined thereby.

In order first to be able to apply this concept to God at all, the objection must be
met that personality, as presupposing limitation by an other, hence finitude, is
inapplicable as a determination for the Absolute. This objection is essentially met
only through postulates. It is to be invalid because the fact that besides God there
are created beings is not to be a barrier that violates the concept of the perfect
being. Self-limitation is to be inseparable from the perfect nature. ``God mirrors
the fullness of his essence in the inner infinity of his self-consciousness and
thereby in truth takes it into possession. For an all-perfect being that does not
know of its perfection lacks a very essential perfection. God limits his power
inasmuch as he calls forth out of his eternal depth of life a world of created
beings, to which he gives, in a derived sense, to have life in themselves; but
precisely by being the power in a free world, he reveals the inner greatness of his
power. For not that power is the true one which tolerates no stirring of freedom
outside itself, because it itself would immediately be all and work all; but that
power is the true one which creates freedom, and which nonetheless is able to make
itself all in all.''

These determinations, which are not reached through a deduction, but only thetically
set forth, link together immediately self-abolishing oppositions (e.g. ``a power that
creates freedom,'' ``in a derived sense to have life in themselves''), and transfer
concepts from one sphere to another where precisely that is lacking which gives them
their determinateness. Self-consciousness is postulated of God as \textit{conditio
sine qua non} for perfection. But self-consciousness, such as we know it in
\emph{man}, that is, such as we know it at all, is possible only through the
opposition to otherness, hence through limitedness. But in God it is to precede the
limitation by an other and nevertheless be self-consciousness. The limitation
necessary for human self-consciousness is then yet brought into God \emph{afterward}
by the fact that the self-conscious God limits himself by calling forth out of his
eternal depth of life that world of created beings, to which he gives, in a derived
sense, to have life in themselves. He limits himself, and this limitation becomes not
only a limitation by a spiritual, by a life of freedom centralized in itself, which
has the possibility of setting itself in a relation of opposition to God's will, but
a limitation by a material, by an other opposed to God's spiritual essence, and yet
it is not to do harm to the infinity: the infinity is to be \emph{inner} infinity.
But by what then does God's essence distinguish itself from every single spiritual,
indeed, from every single living creature, which also has inner infinity? \emph{Inner}
infinity is relative abstract infinity, relative abstract relation to itself through
its abolished other, or else possibility's infinity; but can this be an adequate
determination of God? is it a determination that removes finitude from him? These
questions we can only answer in the negative.

But inasmuch as it does not succeed in freeing the determination of personality from
limitation, finitude, the consequence becomes that God is by that determination
constantly drawn down into the sphere of finite representation, that the representation
of the sensibly determined individuality slips in under the concept of God's
personality. This representation of the sensible single individual Martensen cannot
and will not keep away from the apprehension of the divine. He says\footnote{Martensen:
Om Tro og Viden, p. 47.}: ``the personal God must inevitably be thought analogous with
the psychological consciousness,'' and by ``analogous'' is here understood, as appears
from the subsequent developments, more than a mere analogy: there is understood thereby
the identity of form. The necessity of holding fast that finite representation lies in
the double that is aimed at through the determination of God as personal. There is aimed
at, in the first place, to give God a being \emph{outside our thought}: precisely the
determination of personality is to hold God in the single existence's distinction from
the human, to assert the transcendence of representation. But this distinction by an
``outside,'' by which emphasis is laid on God's not merely being in our thought, our
spirit, not merely having the being of the ideal, makes precisely God \emph{finite} by
giving his being a determinateness of \emph{sensibility}, even though this is again
formally negated. There is aimed at, in \emph{the second} place, to free God from the
law of necessity, to find a point of attachment for the apprehension of will and freedom
criticized above, and since this cannot be found above the universal, it is sought
\emph{below} it in the representation of the single being, which in the finitude that
the relation of consciousness intimates has a possibility for the will's arbitrary
loosening from the universal of the necessity of reason,---admittedly only on the
condition that the will unconsciously falls in under the universal of natural necessity;
and the determination of the human single being is therefore transferred onto the divine.
Precisely for the reason that this double is sought in the determination of God's
personality, the concept of an idea of subjectivity, of an absolute subjectivity, that
is, of an absolutely concrete subjectivity, does not satisfy---one that does not, like the
human, in abstract self-consciousness set itself off from its object of consciousness and
its content of consciousness, but in its being-for-itself sets and idealizingly
comprehends the totality of the real. Without one's giving account of the transformation of
the concept that must take place in the transference from the human personality, which is
bound to the finitude-form of consciousness, to the absolutely infinite, there is demanded,
in opposition to the concept of an absolute subjectivity, an ``actual'' subject, a ``real
subject'';\footnote{cf. Martensen: Om Tro og Viden pp. 38, 49, 53.} but actual and real
become for Martensen synonymous with limited, finite, with an existing in the form of the
human will and self-consciousness, although this last is by him explained only as a form
of mirroring.

The same apprehension also makes itself felt in the closer determinations of this divine
personality's relation to man in revelation, and of the organ for man's apprehension of the
personal God.

Concerning the relation of revelation it is said\footnote{Martensen: Om Tro og Viden, p. 93
f.} that if God is the personal, willing God, then we can cognize neither more nor less of
him than he himself will reveal to us. ``It already holds in the relation between man and man
that a personality can \emph{only} be cognized by its free \emph{actions}, and that every
closer acquaintance rests on the fact that it itself will come to meet us and open itself to
us in free communication; and this holds in the most eminent sense of the divine personality.
\textit{A priori} construct him we cannot; for only the logically necessary, but not the free
spirit, lets itself be constructed \textit{a priori}. Even the human genius cannot be
constructed \textit{a priori}, but only cognized when its wondrous reality is given. Only
after the reality can we cognize that such a thing is possible, since its possibility far
surpasses the mere possibilities of reason.'' ---It is here striking that the representation
of the separated single existence is again brought in, and that it is therefore forgotten
that God's essence is determined as love, love accordingly as a necessary communication of
essence, which, precisely according to Martensen, must also be apprehended as a communication
to cognition. By the skewed analogy with the genius it is likewise overlooked that there is,
in relation to God, no question of constructing \textit{a priori} a historical future, but of
comprehending a universal, eternal fundamental relation, a relation of essence. When there is
then added: ``no revelation of personality lets itself be dissolved into mere thought-
determinations. It can only be grasped by intuition like everything individual and original,
and here there is indeed question of the absolute single being,''---then it is forgotten that
the ``\emph{absolute single being}'' is \textit{eo ipso} \textit{ens universalissimum}, and
when this last determination is overlooked, then there comes no comprehension at all, but only
a sensible intuiting that relates itself to a sensibly determined. But the character of the
``intuition'' by which the higher relation of experience is to be realized, and which is to
acquire its significance in the grounding of the higher science, will be more closely
illuminated in the critique of the latter's concept.

\subsection{The Ethical Concept of God in Relation to the Physical}
\label{ssec:ethical-physical}

The juxtaposition of the \emph{logical} concept of God with that which is designated
as the ethical has accordingly shown the superiority of the first, has shown that the
logical, apprehended in its unity with the metaphysical, embraces also the
essence-determinations of the ethical---something Martensen himself involuntarily comes
to concede by using the determinations the absolute idea, the idea of the good;---while
that which was called the ethical concept of God, on account of the finite
representation of personality, will, and freedom, does not reach the in-truth ethical,
and, instead of becoming a higher than the logical, remains in the alogical, for which
reason also the attempts to deduce this ``concept of God's'' determinations rationally
must fail. And when we now pass over from the consideration of the ``ethical'' concept
of God's content in relation to the logical to the consideration of its relation to the
physical, then the essence-connection between the logical and the physical will entail
that the result becomes an analogous one. When nature is seen as a living unity, as
teleologically striving toward spirit, as that whose essence is spirit, and whose
concept first finds its fulfillment in spirit, then nature is seen as the objectively
existing rational, the logical as immanent in it, and logic as metaphysics expresses at
once nature's and spirit's reason-determined essence, becomes at once nature's logic and
spirit's logic. And in man the natural becomes the will's foundation, and, as the
reason-determined real, the basis of the ethical, that which is to give the free,
conscious spirit's action fullness and reality, to become, clarified through cognition,
content for the ethical will.

In Martensen's apprehension of the physical, nature's reason-determinateness is not
overlooked either. Reason is in nature, and nature does the works of thought. But nature
is yet only the blindly existing: blind in the subjective sense, in so far as nature is
not self-thinking, blind in the objective sense because nature exists only with
necessity, not with freedom, and every unfree necessity is a blind, so to speak
fatalistic, necessity. And the physical concept of God, which has this blindly existing
nature for its substratum, no dialectic is to be able to bring into true unity with the
logical, because the logical God is the purely ideal, indifferent to existence. Existence
and idea are to have their actual point of unity only in a will: only in the ethical God
is their perfect unity to be found.

In the ethical concept of God the physical is therefore to be included as a moment and as
a necessary moment, but in such a way that the ethical, just as it transcended the logical
according to its whole extent, also in its absolute freedom raises itself above all the
physical, again negates it as a condition.

\subsection{The Duality in the Conception}
\label{ssec:duality}

In these determinations there shows itself a duality in the determination of the
relation.

When the ``ethical concept of God'' is to be made comprehensible, made the object of a
cognition, then that side must be accentuated that the physical is a necessary moment,
a necessary condition for the ethical, and the physical then acquires essential
validity, is seen as determining the ethical.

When, on the contrary, the accent falls on the ``ethical'' representation of God, on the
representation of freedom in the meaning indicated above, then the essentiality of the
natural is again abolished. God becomes nature-less and it becomes impossible, from the
representation of this God, to reach the determination of the physical.

The duality becomes at once clear in the determinations Martensen gives inasmuch as he is
\emph{more closely} to indicate the relation between the ethical and the physical in God.
We can leave aside such statements, perhaps according to form more accidental, as e.g.
that God's ethical will cannot be without a nature in relation to which it can determine
itself as ethical\footnote{Martensen: Om Tro og Viden. p. 51.}, a statement that would
point partly back to an impracticable analogy with the human, partly to the doctrine, not
approved by Martensen, of a coming-to-be God who frees himself out of an obscure ground to
spirit and personality. We hold to other statements that, as definitively determining the
relation, cannot have this character of accidentality. It is said\footnote{Same place. p.
68.} that God, as the absolutely knowing and willing, is at the same time to be the
absolutely able. God's ethical will is not to be thinkable without an eternal nature,
physis, in God, without which his will would be without instruments and potencies of
production; it is not to be thinkable without an eternal fullness of power. ---In these
determinations there is ascribed to the physical a real validity and an essential
significance for the will and the ethical, precisely by the fact that the will, without
that physis, is set as being without possibility of production, impotent for bringing-forth
and in general for working\footnote{If the will is seen as that by which the divine
essence's being is eternally set, then this setting is in the highest sense a working, and,
more determinately, a working with respect to that to whose essence a physis is to stand in
necessary relation; the will's working can accordingly not here either be thought without
``instruments,'' that is, without an eternally existing physical, when it is not to be
thinkable without a physical with respect to that existing outside God.}. And there lies in
them the further-reaching consequence that when the will, in order to work, presupposes
instruments, and these instruments accordingly cannot be thought as set by the will, since
in this setting it would precisely have worked without instruments; then the instruments
must acquire equal originality with the will, and, when the determination physis is taken
seriously, we get either the same dualism that Martensen reproaches Nielsen with, or we get
a natural God, for whose physical necessity the will becomes the form of activity. But
precisely the opposite is aimed at, and therefore the following determinations recall that
on which these consequences rest. For it is said, in immediate connection with what is
stated above, that that eternal fullness of power is an all-possibility that relates itself
to the ethical in God as the serving, subordinate, to the superordinate, ruling, and royal,
as God's inauthentic other-essence to his authentic innermost essence. ``God is, as the God
of freedom, the lord of possibilities, and can freely reveal himself, which does not hold of
the logical and physical God, whose possibilities blindly pass over into their consequences.
God is the absolutely nature-free (not nature-less), and as such he is in himself the perfect
spirit, the perfect personality, the true God, who has the logical and physical for
attributes, predicates, qualities, is a living unity of different moments.''

Inasmuch as the fullness of power is set as equal to an all-possibility, inasmuch as it is
determined as the serving, subordinate, hence becomes that which is to be mastered by
something that falls outside the fullness of power, inasmuch as God's freedom in the meaning
developed above is appealed to, and inasmuch as God is determined as the absolutely
nature-free, as absolutely, unconditionally disposing over the physical, his instrument, then
the determination physis is deprived of its reality, for inasmuch as the divine will, whose
absolute priority accordingly also in relation to this physis continues to be the
presupposition, unconditionally disposes over the instruments, then the will is not determined
by their nature, it can, while by using them it apparently ascribes to them reality as
instruments, substitute in their place everything possible, and by this substitution their
reality is negated. And when the concept of the ethical God is identified with the concept of
the ``supernatural'' principle, then it becomes still clearer that physis in God becomes only
a semblance of nature, a fantasy's sublimation of the actual nature, which is taken up only
for the benefit of the analogy with the human, and for the benefit of an illusory deduction.
And yet the divine essence is to be unity, true unity of the different moments. This becomes
only a postulate: the reality of nature and the necessity of nature must always become
incomprehensible from such a representation of God.

\subsection{Creation}
\label{ssec:creation}

This consequence will confirm itself when the ``ethical,'' personal God is set in
relation to actual nature, and when a deduction of the doctrine of creation and of the
miracles is attempted.

If we first analyze the representation of creation such as the religious consciousness
apprehends it, then there lie in it the following determinations:

\begin{enumerate}[label=\arabic*)]
\item Creation is creation out of nothing; but that which is created out of nothing,
  that is, in its origin is nothing, is in its being continuingly a nothing, not
  \textit{natura}, but \textit{creatura}. By the representation of creation the concept
  of nature is therefore negated.
\item Inasmuch as the created is set by the in-itself perfect God, determined as spirit,
  and set in the determination of a material, i.e. in a determination absolutely opposed
  to the purely spiritual God's essence, then creation is not grounded in God's essence's
  necessity, but in a will unbound by the essence, in an absolute arbitrariness, whose
  working, inasmuch as it falls purely outside the necessity of essence and sets that
  which is opposed to the essence, is the absolutely incomprehensible.
\end{enumerate}

Against these consequences, which sharply indicate the simple conceptual relations,
Martensen makes objection. He maintains against the first partly that a physis in God is
the presupposition of creation, and that, inasmuch as this physis becomes equal to an
all-possibility, the nothing out of which the world is created is to be one with God's
will's eternal possibilities;---partly that creation is not to set a nothing, but a
something that produces and develops itself in given independence. Against the second
consequence he objects that the ground of creation must not be sought in an absolute
arbitrariness, but in the ethical, in the ethically determined omnipotence.

As regards the objections against the first consequence, the representation of a physis
in God has already above been shown in its inner contradiction. This representation, such
as it has come forward in its originality in Jacob Böhme, has proceeded from a feeling of
unsatisfaction at the representation of nature's origin from a nature-less God, from the
feeling of the nothing-explaining in the explanation of nature's coming-to-be by an
absolutely groundless divine act of will. In order to explain nature, one then places
nature in God; one indirectly pronounces that nature can only be explained by nature
itself, that it has essentiality and self-validity, one accordingly unconsciously betrays
a doubt with respect to the representation of creation. But the involuntary acknowledgment
of nature one again takes back, precisely by placing nature in God, by exchanging the
actual nature for its image in fantasy, from which then the former is to be explained. It
is the nature metamorphosed by fantasy, the nature which, though ``wahrhaftig und
eigentlich,'' is yet ``himmlisch und geistlich,'' that one places in God, and inasmuch as
one makes it the ground for the real nature, one lets on the one hand the metamorphosis
conceal that which deviates from the religious apprehension, on the other hand the
likeness of the name bring about the semblance of an explanation of the actual nature's
existence. But the spiritualized, dematerialized, denatured physis in God explains nothing
with respect to the real nature's being, and an actual nature-determinateness in God would
let the representation of creation be superseded by the concept of an eternal necessary
unfolding of essence.

When nature in God is then made into an all-possibility, and the nothing out of which the
world is created is set as equal to God's will's eternal possibilities, then this
identification becomes an exchanging of two entirely \emph{different} concepts. The eternal
possibilities in the eternally actual God must indeed for God be eternally realized
possibilities, i.e. realities, hence as far as possible from nothing. Were they to be
possibilities that first acquired reality by the finite existence, then there would lie
therein consequences with respect to God's relation to this existence that Martensen would
least of all be able to accept.

The objection that creation does not set a \emph{nothing}, but a \emph{something} that
\emph{produces} and develops itself in \emph{given independence}, abolishes itself by its
immediate combining of opposite, mutually negating determinations. From what consideration
of natural existence this objection has proceeded will be illuminated in what follows.

The objection against the \emph{second} consequence drawn above is carried out in Martensen
in the following way\footnote{Martensen: Om Tro og Viden. p. 59 f.}. It is to be entirely
superficial and unjustified to trace creation back to an absolute act of arbitrariness of
the mere omnipotence and accordingly to set it as the incomprehensible. Creation is to
presuppose unconditionally the ethical concept of God, and without this it is to be entirely
unreasonable to speak of creation; there is to be no question of the mere
omnipotence---a pure physical determination---but of the omnipotence determined by love and
wisdom, hence ethically. In love's ethical ``necessity'' the ground of creation is to be
sought; through it creation is to become comprehensible, not indeed in such a way that one
pretends to fathom its how, but in such a way that it is seen in its necessary connection
with the determination of God's essence.

But this appeal to the ethical becomes completely useless, and would already for the reason
have to become so that, as shown above, the ethical in Martensen's apprehension does not in
reality acquire the physical for its moment. Here there is question of a creation of a
material existing: in what connection does it indeed stand with the determination of the
ethical? Martensen postulates the connection, but does not even make an attempt to show what
is the connecting thing, and what, in the ethical concept of God, is to be able to explain
the possibility of such a creation. Against the postulate that the ethical determinateness in
God is to be the explaining thing, it holds that what is in truth to be ethical must be
determined by God's essence; were the moment of nature accordingly excluded from God's essence,
a creation of the real would have no positive relation to the ethical. And above it was shown
that the moment of nature was in reality excluded, and that one had to exclude it if one was
not, instead of a creation, to get a development of nature. But when therefore nature in God
is reduced to an unreal, to a mere fantasy-image, then one acquires no point of attachment for
the possibility of a creation of the material; but where the possibility is lacking, it is
meaningless to speak of the necessity. Even if one would let love, which yet, according to
Martensen's deduction, is already to have its actual object, hence its real satisfaction, in
the relation of the Trinity, be the ground of explanation for a setting of created spirits,
one would not have come a step further with respect to the creation of the material. And such
a representation of a purely spiritual creation would abolish itself. That which is set by a
creation would have to be a finite; but a finite is a determined by time and space, and time
and space, which cannot be thought without each other, can only be thought where a material
is. A creating of purely spiritual, an existing of created spirits, in actual distinction,
with actual beginning, with an actual life-history without the real existence's fundamental
forms, would accordingly be self-contradictory. However the working of creation is
represented, creation remains equally incomprehensible.\footnote{It naturally explains nothing
that it is said (in the Dogmatics) that ``the Logos releases out of its spiritual depth the
whole manifold of life-forces to self-movement''.}

\subsection{The Miracle}
\label{ssec:miracle}

What holds with respect to \emph{creation}, the universal miracle, must also have
validity with respect to \emph{the miracle}, the partial \emph{creation}.

If we analyze the representation of the miracle, we will accordingly obtain results
analogous to those to which the analysis of the representation of creation led. The
miracle, as the working of the divine, absolutely free, almighty will, is a working
in which the nothingness of the natural is set. It is set inasmuch as omnipotence in
the miracle works without all that is real, and inasmuch as in the miracle the
existing real is set as a nothing. Only apparently does the miracle-working
omnipotence, according to its concept \emph{sole-working}, use means, only apparently
does that which has come to be by the miracle have a real presupposition. The apparent
means is no means, the apparent presupposition is no actual presupposition, for there
is no essence-connection between the means and the goal, between the presupposition and
that which follows after; the former's determinateness is in relation to the latter's
determinateness a nothing. The miracle is a new creation and as creation a creation out
of nothing: only the fact that it is in the partial creation gives the semblance of an
actual presupposition. For a new creation the preceding becomes a stuff whose essence is
negated in the moment of creation, i.e. it becomes a nothing in relation to the new;
that which precedes becomes, in other words, not a positive presupposition, but such a
one as is only negated in order to give place to the new. ---From the miracle's relation
to omnipotence it then follows that the miracle, like creation, must be the
incomprehensible, an explanation of it an impossibility. Of this the religious
consciousness too is clear: for it the miracle is omnipotence's immediate working (``he
spoke and it happened,'' ``he commanded and there it stood''), and that which surpasses
all human thought.

Against the indicated apprehension of the miracle Martensen makes objection from the
same standpoint from which he made objection against the conceptual determinations given
above with respect to creation. The miracle is not to be without a presupposition of
nature, which in the miracle itself acquires a relative validity, and the science of
faith is to be able to explain, not indeed the miracle's how, but the miracle's
possibility in general, inasmuch as the miracle is seen in connection with God's
revelation's economy and the freely self-revealing God is seen as ``the Lord of nature,''
who ``can do what nature cannot.'' The miracle's relative comprehensibility rests
accordingly, like creation's, on its relation to the ethical concept of God, more nearly:
on its connection with the ethical teleology.

In order to assert for the miracle a presupposition of nature, it is maintained that the
ethical world-teleology, in connection with which the miracle is to be apprehended, is not
to abolish the lower circles of existence's autonomy and independence, or contradict
natural science's presupposition of a law-determined order of existence. An ethical view of
the world is declared entirely impossible without a relatively independent physis, and the
miracle itself impossible without an ordered system of natural laws, since it is precisely
to presuppose an order of nature in order therein to be able to come forward as the miracle.
Likewise it is demanded with respect to the miracle that the higher divine activity must
always be logically conditioned, inasmuch as God cannot contradict himself, in which lies
that God's miracle-working must not be annihilating in relation to the natural set by God;
which is also indirectly demanded by the positing of physical conditions for the divine
will's working. ---But these general determinations show themselves not only as
impracticable as soon as the \emph{determinate} miracle is drawn forth; they also acquire
closer determinations by which they are again abolished. The miracle is set in relation to
the teleological. But that which is designated as miracle is determined as that which cannot
be derived from the preceding stages of creation: it is a specifically higher, not only in
the meaning in which a teleological development's more advanced stages of development are a
higher in relation to the lower, and in reality are that by which these are to be
comprehended, but in the meaning in which the supernatural is higher than the natural.
\emph{Inasmuch as the miracle therefore comes forward, the teleology is broken, the order of
nature negated}; the preceding is a nothing in relation to the miracle and this comes
forward as the purely immediate and arbitrary. It is correct that the miracle, as this
single, in time limited miracle, in a certain sense involuntarily presupposes, \emph{though
only in order to abolish it}, an order of natural existence, in so far as it presupposes a
continuity that is broken by the miracle. Were there not a relative continuity, the miracle
could not become visible as miracle; for inasmuch as the miracle was set in all single
discrete points of time, the miracle would become the law, the very order of existence:
\emph{the miracle's infinite discreteness would become the law's infinite continuity}. But
that relative continuity does not designate an order grounded in nature's substantial
essence; it is not an expression for the substantiality itself, but is only the expression
for something arbitrarily set and arbitrarily repeated by the creating will. It is precisely
a ``stage of creation,'' to use Martensen's expression. It is abolished by the miracle and
again arbitrarily set: God's freedom is for the religious consciousness only absolute
inasmuch as it is unconditionally unbound by everything, also by the presuppositions set by
itself. Consistently this consciousness would make everything a miracle; that it, in its
immediacy, does not do so and cannot do so without letting the miracle become invisible,
shows that it never fills the human spirit wholly, that it has the natural consciousness
alongside itself, and only is by constantly negating this.

From this latter there originally presses into the former the representation of the
relatively constant in nature; but inasmuch as the relatively constant is seen as that which
can be broken through by the miracle, it has already ceased to be a law of nature, and when
the religious consciousness sees it in its own illumination, it becomes for the religious
reflection the arbitrarily set by omnipotence, hence miracle, like the specific, striking
miracle, so that the miracle again embraces the whole of existence, but at the same time
becomes invisible as miracle, for which reason consciousness is involuntarily led back to
that point of departure, but only to run through the same cycle anew.

As regards the miracle's comprehensibility, it is attempted to be asserted partly through a
polemic against the representation of the immutability of the laws of nature, partly through
the attachment to the ethical, through the determination of the natural's serving relation
to this. In order to procure a place for the miracle, the representation of the immutability
of the laws of nature is sought to be abolished. But inasmuch as the apprehension of the laws
of nature as immutable is sought to be refuted, there comes forward a striking unclarity, both
with respect to the difference between a change of the laws of nature and the same unchanged
laws of nature's different working under changed conditions, and with respect to the relation
between the laws of nature and the mathematical laws, whose inseparability the laws of
mechanics show. The attempt at a refutation of that representation rests essentially on the
fact that the way by which we come to the determination of the laws of nature must determine
their character, and since the way is that of experience, the laws of nature are to have no
other validity than that which experience can secure for them. Inasmuch as the laws of nature
are determined by the way of experience, through induction and analogy, they can, it is
said\footnote{Martensen: Om Tro og Viden. p. 108 ff.}, only acquire relative universality and
necessity, their immutable validity can never acquire more than the highest degree of
probability. And they would, even if within this determinate circle they had unconditional
validity, as laws of experience yet only have validity in a given facticity, within this
determinate, this present existence, but not have it for all existence. Natural science can,
as a mere science of experience, decide nothing about the in an absolute sense possible or
impossible; it does not even have, with respect to the reality to which it relates itself, a
cognition that embraces beginning and end: these are for it enveloped in an impenetrable
darkness; it moves only as it were in a middle, within a circle of conditions that have once
for all factually been made, it does not itself even know how. It can therefore not
contradict that there can be a higher existence than nature's, which can penetrate this and
work upon it, be it redemptively or transfiguringly, set it as organ and means for a goal
that is higher than nature's own. Since the miracle does not refer itself to the world of
pure concepts and pure forms, but to existence, it is precisely the question whether nature
itself is the highest existence, which as such must not be able to be exchanged, or whether
that higher existence exists. Only from a physical concept of God as the highest could the
immutability of the laws of nature be demonstrated.

In this argumentation it is \emph{first overlooked} that the fact \emph{that the laws of
nature can come to our knowledge only by the aid of experience} does not abolish the
possibility of a comprehension that gives them a higher necessity than that \emph{attainable
by the way of mere experience}; in other words: does not exclude the possibility \emph{that
the empirical law can be transformed into a rational one}, a transformation that is precisely
the task of natural cognition, and whose possibility it is incumbent on the miracle's
defenders to disprove. Natural cognition is and remains indeed always \emph{an experiential
cognition}, but not in the meaning that it is to rest on \emph{an isolated experience}; it
takes its point of departure in experience in order from it to reach a comprehension whose
possibility is conceded when there is talk of reason in nature, of nature's doing the works
of thought. ---When it is then concluded that the laws of nature, as only valid in a given
facticity, within this determinate existence, cannot have unconditional immutability and
necessity, then here partly the confusion touched upon above between a change in the laws of
nature and in the conditions is to be removed, partly it is to be urged that the conclusion's
validity depends on the relation that is set between this form of existence and the supposed
higher one. Is this---to which the expression \emph{world-teleology} points---\emph{a higher
stage of development}, then the laws retain, under the new conditions of activity, under the
new complications, unchanged their validity, just as we see the laws of the inorganic retain
it within the organic; they hold for all existence, i.e. for all natural existence, and we
get in the higher not a miracle, at most a \textit{mirabile}. \emph{Nature's primitive laws
must, since nature must always be thought as a whole, be the same at all stages of nature's
development; for the laws are precisely the expression for the unfolded essence thought in
activity}. In order to come to the abolition of the law of nature's validity one would have
to presuppose that the higher had come forward by a break, by an absolute abolition of the
continuity. But then the higher would have come to be by a new creation, and the relation
would---what is precisely to be avoided---be incomprehensible. When the higher, the miracle,
is not to be thought as a new stage of development, but as a re-creating, then it leads to the
consequence that the miracle, the partial creation, is to enter into a reality that is changed
only at this point, and that accordingly this reality, thus partially changed, would become
the contradictory: a defective reality, an incomplete whole, a discordant harmony. Is there
posited ``an ordered system of natural laws,'' then the break at this one point is a break on
the whole of the system. And even if one looks most nearly to the single real that is negated
in the miracle, then an annihilation of even the most insignificant real, an abolition of it
against the law, is a break on the whole order of existence, a disturbance of the whole.
\emph{And that change in existence reaches from existence into the world of pure concepts and
pure forms.} Analyze whatever miracle one will, it will bring a contradiction into the sphere
of concepts, inasmuch as it is a working without means, a causing of that which negates its
cause's essence, a setting of a real without all that is real, a setting that negates that set
by the eternal activity; and it will, inasmuch as \emph{all physical existence is
mathematically determined}, bring a contradiction into the world of pure forms: that of the
mathematical; for it is a setting of a something, a by quantum determined, out of a nothing,
of a form-determined out of the formless. It becomes not only the merely physical concept of
God, from which the miracle is negated, and the immutability of the laws of nature, which is
only the expression for nature's essentiality, asserted; as soon as there is posited a ground
of nature, a physis in God, the miracle is set as impossible. And is the moment of nature in
God negated, openly or covertly, then the miracle is again, though for another reason, the
impossible, the purely incomprehensible, like creation. Because the miracle is denied, nature
as nature is not set as the highest, but the higher being than nature: the being of spirit, is
seen as nature's fulfillment, not as nature's arbitrary negation. The laws that, as holding at
a \emph{determinate} stage of nature's development, hold for nature in general, hold for nature
as the nature ensouled by spirit, as the nature that has its goal in spirit. Nature is spirit's
abiding, essential foundation.

From the \emph{negative proof for the miracle's comprehensibility} we turn to \emph{the
positive}, it rests on the ethical concept of God. By this there is, it is said\footnote{Martensen:
Om Tro og Viden. p. 111 f.}, not only set a living relation of interaction between God and the
creature, but the world-teleology is at the same time determined as ethical teleology; that is,
the ethical is, as the fundamental purpose, thus also the innermost fundamental-determining in
existence, that for whose sake all else is, and which all else must serve. This is yet not to mean
that the ethical teleology at all points of existence would come forward entirely arbitrarily and
purely immediately, and everything without further ado be able to be stamped with the ethical and
religious relation of dependence. As indicated above, the ethical world-teleology is precisely to
presuppose and confirm a relatively independent physis, a law-determined order of existence. The
meaning is to be that the ethical view of the world and the ethical concept of God necessarily
entail that nature in its arrangement cannot make impossible the divine world-plan, or that the
concept of nature cannot deny the concept of creation, that nature cannot deny its dependence on
the Creator. For the ethical view of the world nature is to be an intermediate being that has the
significance, not only of being itself and of fulfilling its own purposes, but also of being organ,
means for freedom, and of serving purposes that are far higher than its own. It is to be receptive
to freedom's workings, of which it is said that, when the God of freedom is presupposed, it would be
most strange to restrict them to human freedom's. And when now the divine world-plan entailed that
in mankind's history there entered periods, determinate times, in which God in his relation to
mankind would in a pregnant and entirely unambiguous way reveal himself as the one he always is, as
the Lord of nature, in order through workings that cannot be explained from the created to bring
himself nearer to human consciousness, because the content of his world-plan was an economy for
man's salvation, then such times would in a peculiar and strictest meaning be \emph{the times of
the miracle or the great times of revelation}. And the science of revelation or of religion would
have to seek a closer understanding, seek to see these miracles' significance in and with the great
total connection, naturally without wishing to comprehend the miracle's how.

By the reference to the ethical concept of God, however, nothing is gained. This lies already in
what was developed with respect to creation; it becomes, when account is taken of the form in which
the argumentation comes forward in Martensen, clear in the following way.

\emph{Spirit's primacy in relation to nature}; the fact that nature serves spirit and is receptive
to freedom's workings, does not make the miracle comprehensible. Nature serves spirit \emph{as
nature}, i.e. in accordance with its laws, and freedom's workings are not, like the miracle's,
arbitrary and purely immediate, but through natural means, i.e. again in accordance with the
cognized reason-determined laws of nature. Inasmuch as nature serves the human rational freedom, the
free spirit bows beneath the infinite reason that is at once its own and nature's essence; only the
fantastic freedom dreams of a magical, immediate dominion over the subjugated nature. Nature, which
in the living immediately shows the power of the ideal, does not deny its determinateness by the
ideal, which is precisely its essence; but the concept of nature must deny the concept of creation,
hence also the concept of the miracle, because the concept of creation denies the concept of
nature, negates all of nature's reality. The realization of reason-determined freedom through
reason-determined nature includes no difficulty for thought, but a revelation of God as \emph{the
Lord of nature in the indicated meaning}, as the one who can do what nature cannot, and, mark well,
\emph{can do it within nature, in relation to a natural, without all that is natural}, is the
incomprehensible.

\emph{The miracle can next not become explicable by any historical connection}, for \emph{just as
the miracle has no nature, so it has no history either}. It is, as expression for the unconditioned,
almighty will, the absolutely sudden, discrete, connectionless; its comprehension out of a total
connection is impossible, is a contradiction.

The miracle finally does not become comprehensible either by being supposed to be a condition for
God's unambiguous revelation. God is in the miracle to reveal himself in an entirely unambiguous way,
by workings that cannot be explained from the created: the miracle was thus to \emph{create faith}.
But, according to Martensen's own premises, it cannot be known that this working, which takes place
in nature, cannot have natural causes. Experience, which, according to Martensen's determination, is
\emph{natural cognition's only source}, cannot say it: since it is indeed to know nothing about
possibility or impossibility. But in what way then does God's revelation become unambiguous? And how
is it to be understood that God in this determinate working is to reveal himself more clearly than in
his natural workings, bring himself nearer to human consciousness than through his
essence-revelation in nature and reason? It is only to be understood thus, \emph{that the miracle
unambiguously reveals God when man has by a miracle been led to see it as miracle}. That is, in other
words: \emph{the miracle is acknowledged as miracle only through faith: through a miracle; it
presupposes faith: the primitive miracle, the miracle in us, by which we first see the miracle
outside us}; not the outer miracle, but the wondrous working in the spirit becomes accordingly the
decisive thing with respect to the unambiguous revelation. That the miracle is apprehended only by a
miracle, not understood by thought, lies also in what Martensen says, that the sought understanding
of the miracle is unthinkable without the mind receptive to revelation; for revelation and its
reception is indeed precisely a miracle.

Inasmuch as the miracle thus presupposes itself, and is as it were the only way to itself, every
attempt to make it comprehensible, also in that more indeterminate meaning, must fail, and the
attempt to make it comprehensible in this meaning leads inevitably, inasmuch as the miracle is to be
fitted into a teleological connection that at the same time becomes a natural connection, to
involuntary attempts also to determine its how, which was excluded by the presuppositions.

\begin{center}---\end{center}

\subsection{The General View of Existence}
\label{ssec:general-view}

\emph{The general view of existence}, to which the apprehension of creation and the
miracle points back, is pronounced in the Dogmatics. It is set as the task to bring
about a Christian apprehension of existence, which yet \emph{must not come into conflict
with nature and reason according to their truth}. The standpoint for this apprehension
of existence is to be taken from that which is existence's highest, its τέλος, from
Christ's person, who is the midpoint of the holy history and Christianity's proper
content. The holy history is to explain history, history nature. But by the
representation of a holy history and especially of Christ there is pointed to the
miracle, to an absolutely new beginning, a spiritual creation in mankind, which is to
have not only ethical, but also cosmic significance, the deepest significance for nature.
The miracle of the Incarnation, and the miracle of Inspiration attached to it, by which a
new sense arises in mankind, a new beginning is set in mankind's consciousness, is
Christianity's fundamental miracles, which the supernaturalistic view of the world is to
assert against naturalism and rationalism, but in such a way that nature and reason see
their right, that no disturbing dualism is brought in.

With respect to nature it is then set as the Christian intuition, what was developed
above: that it is not an in-itself eternal and concluded system of laws and forces, but a
system that finds itself in a continued teleological development, a continued creation, so
that there must be able to be thought to enter new potencies, new laws and forces, whose
\emph{entrance is prepared and prefigured by the preceding stages of creation, but which
cannot be derived from these}. Christ's revelation, as the last potency in the work of
creation, is therefore to be thought to work \emph{teleologically changing and
determining} upon the lower potencies, inasmuch as these do not have an eternal and
concluded significance in themselves, but have only a temporal intermediate significance.
The point of unity of the natural and supernatural is thus to lie in nature's
teleological determination for God's kingdom and in the receptivity and formability for
the supernatural creative activity thereby given. \emph{Nature is not to deny the concept
of creation, but just as little is creation to deny the concept of nature}. For although
the new creation in Christ is an abolition of this nature's laws, it is yet by no means to
be an abolition of the very concept of nature. For it is precisely nature's concept to be,
not the hampering barrier for freedom, but freedom's organ.

By freedom there is pointed to the kingdom of history, but within this is already
intimated the repetition of the same relation that obtained in the kingdom of nature. For
rationalism history is the unfolding of human nature, of the human essence's
nature-determined dispositions; the universal reason (κοινὸς λόγος) is for it the only
source of cognition; it acknowledges no other truths than those that are explications of
mankind's innate reason. Christianity, on the contrary, must demand that to the new source
of life that is opened in Christ there corresponds a new source of cognition, a kingdom of
divine, hitherto concealed counsels, a kingdom of new cognitions that do not let
themselves be derived from the concept of mankind's advancing development of reason. But
just as the miracle did not conflict with nature, so these new, supernaturally
communicated cognitions are by no means to conflict with mankind's general cognitions of
reason, although they in many ways \emph{more closely determine these}; for partly they
relate themselves \emph{completingly and concludingly} to the proper cognitions of reason,
partly \emph{redemptively}, inasmuch as they free them from the darkness with which they,
on account of the all-embracing sinfulness, are afflicted. By them there is just as little
to be posited a dualism as by the doctrine of the two creations: \emph{there is only one
kingdom of reason, although in this there are two potencies of reason's revelation}.
Objectively the unity is to lie in this, that it is \emph{the same Logos} that reveals
itself in the two creations, so that the revelation of the Logos in Christ is the
revelation in higher potency, therein different from the universal revelation of the Logos,
that it is the world-completing and world-redeeming revelation, while that universal one is
merely the world-creating and sustaining. Subjectively the unity is to lie in this, that
human reason is receptive to Christ's spirit as the spirit of the world's completion and
the world's redemption: a receptivity through which reason is to be raised to a higher
stage of productivity. That the revelation in Christ is to conflict with reason's laws is
only to be conceded in the meaning that it, in the same way as the moral laws are abolished
in their abstract independence in order to be confirmed in love, abolishes the laws of
reason in their abstraction in order to confirm them in the Christ-wisdom that is the
fullness of the law of reason.

The relation that this ``Christian intuition'' of existence posits can accordingly in
brief be determined thus: \emph{the miracle is nature's higher potency, revelation
reason's; the miracle is prefigured and prepared in nature, revelation in reason, but the
first cannot be derived from the last.} The relation has thus a duality in itself:
\emph{on the one hand there is pointed to the higher unity, on the other hand the unity is
again abolished and the break enters between the higher set by a new creation and the
preceding lower.} The impracticability of this duality has already been shown by the
doctrine of creation and the miracle; what there has been developed with respect to the
concept of creation, which embraces the concept ``new creation,'' ``continued creation,''
and with respect to teleology's relation to the absolutely new beginning, finds, just as
that which was developed about the lack of an attachment for the physical in the ``ethical
concept of God,'' about the immutability of the laws of nature, about freedom's relation to
nature, etc., its application here too. The turns of phrase peculiar to the given
development do not lead closer to the aimed-at goal. When there is thus talk of a
\emph{receptivity and formability of nature for the supernatural}, and in analogy herewith
of \emph{the general human reason's receptivity and formability for the supernatural
revelation}, then \emph{receptivity and formability presuppose essential unity}: only what
agrees with my essence can I receive; for receptivity is inseparable from a self-activity
corresponding to the received, and the self-activity expresses the essence. Without
self-activity the reception would be an annihilation of the receiving thing. When there is
therefore talk of the natural's receptivity for the supernatural, reason's for revelation,
then we must either set the difference as a merely quantitative one, or, when we will hold
it fast as a decisive difference of quality, let the receptivity become an empty word
without significance, let the ``teleological reaction,'' i.e. the afterward set teleology,
be an actual re-creation, the reception of revelation into the I a break in the I, a
re-creation that makes self-consciousness's identity incomprehensible.

It is on the whole characteristic of the given development that the determinations suffer
from ambiguity and vagueness. When it is thus said of the higher, which cannot be derived
from the lower, that it is ``prepared and prefigured'' by this, then there lies in the
``preparation'' an \emph{objective connection}, in the ``prefiguration'' a \emph{merely
figurative, merely subjectively valid} one. By the combination of these expressions it is
accordingly held in suspense of what kind the relation is; it becomes at once a real and a
not-real relation, at once inner connection and break on the connection. And when there is
talk of the lower's \emph{heightening and confirmation in the heightening}, and of the
higher potencies' reaction, then these expressions acquire a significance only for the
fantasy so long as they are held in the universal, not attempted to be brought into
relation with the single and concrete. The change and determination of the lower potencies
is to coincide with the realization of their receptivity for the supernatural creative
activity. But this realization, which would have to be brought about by removing that which
hampered the possibility, would, since the hampering thing could only be a real, have real
consequences and be visible in nature---unless precisely nature were annihilated by this
working, which is indeed not the presupposition. But let one only attempt in the concrete to
demonstrate the change in nature's laws, in the logical and mathematical laws, that was
supposed to have taken place by the Incarnation, which is indeed to have the deepest cosmic
significance, or the heightening of the forces of nature by the single miracle, and one will
soon discover the hopelessness of the attempt. And even without this: when one attempts to
think the thought of a defective nature and of a defective reason, which were yet nature and
reason---a thought that that representation of higher potencies presupposes---then the
representation's inner contradiction will show itself.

\begin{center}---\end{center}

\subsection{The Result of the Critique of the Concept of God}
\label{ssec:god-result}

\emph{If we summarize the results of the investigation of the ethical concept of God's
relation to the physical}, they become the following:

\begin{enumerate}
\item The ethical concept of God does not in reality acquire the physical in itself as
  a moment; it becomes the concept of a \emph{nature-less God}, and therefore \emph{the
  natural's existence, and divine workings in the natural, become inexplicable} from it.
\item The freedom by which the ``ethical'' concept of God is determined does not give it
  a worth higher than that which belongs to the reason-determined nature apprehended
  according to its true concept; this freedom, which becomes only \emph{essence-less
  arbitrariness}, becomes on the contrary a lower than the reason-necessity manifest in
  nature, and it becomes a self-contradictory and self-abolishing.
\item The apprehension of existence that rests on this ethical concept of God does not
  become the higher. It negates the law of reason in natural existence; it abolishes,
  with nature's laws, \emph{the concept of nature}, and history becomes for it, whether
  it is seen from God's or man's standpoint, only a \emph{necessity-less, i.e. arbitrary
  freedom's working}. (Precisely because this apprehension sees in the human essence the
  \emph{necessity-less will freed from nature's necessity as the first}, and therefore in
  mankind's history sees only that freedom's kingdom, inasmuch as spirit, as ``the free,''
  becomes the one unbound by law, it accentuates ``the historical'' because this, as that
  which is withdrawn from the law, forms the analogy to the divine ``working of freedom''
  that comes forward in the free revelation and in the miracle. Yet this historical
  acquires also for another reason so great a significance for it, namely because it, by
  being the medium in which God immediately comes forward and reveals ``himself,''
  satisfies the religious consciousness's urge to acquire \emph{the eternal for intuition
  in that form of sensibility}\footnote{cf. O. Waage: J. P. Mynster og de philosophiske
  Bevægelser paa hans Tid i Danmark, p. 122, cf. 116, 144, 163.}, which is inseparable
  from this consciousness at the standpoint of positive religion).
\end{enumerate}

\begin{center}---\end{center}

Thus we are led, by the investigation of the determinations of God's essence given by
Martensen, from the as ethical named, but not in reality ethical, concept of God back to
the logical, apprehended according to the logical's true and complete meaning, according
to which it expresses the essence-determinations of nature and spirit. And we are led from
the representation of the natural's coming-to-be by an arbitrary creative act and of the
outer revelation proceeding from the will's arbitrariness back to the concept of a nature
standing in eternal essence-relation to the divine, resting on immutable laws, and of the
eternal essence-communication, determined by the divine essence's free necessity, to the
human spirit grounded in God.

\begin{center}---\end{center}

\subsection{What is Typical for Theology in the Criticized Conception}
\label{ssec:typical}

What in the foregoing has been developed in close attachment to the theological attempt
present in Martensen's writings is now to be seen as holding for theology in general,
inasmuch as Martensen's apprehension, in its \emph{decisive fundamental features}, is shown
to be typical for theology.

In order to demonstrate the typical character of his apprehension, it will not be sufficient
to demonstrate corresponding fundamental features in significant representatives of theology
at various times. By this way, on which one alongside the agreeing would always be able to
find the more or less deviating, one would only be able to reach an approximate. The
apprehension's typical character can only be asserted incontestably by going back to that
which becomes theology's object: the content of the religious consciousness, and when the
form is sought that this content must consistently assume when it is to become content for a
``science of faith.''

It will here, where there is question of the apprehension of the concept of God, be enough
to draw forth \emph{one phenomenon that is as it were the religious consciousness's
ur-phenomenon and runs through all, even the lowest, religious forms}, namely: \emph{prayer}.
In the act of prayer the peculiarity of the positively religious lies concentrated; the
analysis of the concept of prayer gives the essential moments for the determination of the
positive religions' apprehension of the divine.

In the concept of prayer there lie essentially two moments: \emph{1) prayer is prayer to the
personal God, whose will is determined by the prayer; 2) the object of prayer is the
miracle.}

1.~Inasmuch as the praying one prays to God, in the prayer ``says Thou to his God,'' he
relates himself to God as to a \emph{single personality outside himself}, and in the very
form of prayer there lies expressed that this \emph{single personality can be moved by the
prayer, be affected from man's side}. As the one to whom there is prayed, God is the one who
is not bound by any eternal law, but has the possibility of constantly new counsels, and of
counsels that relate themselves particularly to the single one who prays, to his welfare,
his wishes, his blessedness. ``Ask, and it shall be given you,'' is the expression for the
immediate religious consciousness's apprehension of prayer. ``God, who is good, will give you
good gifts when you ask him for them'' expresses \emph{prayer's conditioning power}. ``Give
us this day our daily bread'' expresses that prayer relates itself to a determinate object.
Even when there is prayed: ``thy will be done,'' then the will whose fulfillment is prayed for
is that will of omnipotence which wills \emph{the believing and praying one's best,
omnipotence is goodness's omnipotence}. Is there prayed: ``not my, but thy will be done,''
then the prayer again concerns the same will of omnipotence apprehended in the same way. In
this prayer a beginning reflection is manifest, which is yet bound by the immediate religious
feeling. Inasmuch as God is the religious feeling's object as the absolutely exalted,
knowing, the prayer acquires this form of resignation, but \emph{prayer's determining of
God's will is therefore not negated. To an immutable will one cannot pray}. The prayer leaves
to God in what way he in his goodness will best fulfill the praying one's wishes, but the
latter does not doubt prayer's significance for the fulfillment of these wishes. Precisely
the fact that to the determinate prayer there is added the prayer: see not my, but thy will be
done, includes that the praying one, who represents to himself the possibility of a divine
willing that does not at all refer itself to the first prayer's object, will yet by the prayer
determine God in relation to this. Is prayer apprehended only as the lively expression for
\emph{the contemplation of God}, then the specific for prayer is taken away, another concept
is slipped in under the name of prayer. Is prayer's working to be only the \emph{indirect
subjective reaction} upon the praying one, then prayer's essence is undermined by reflection
and the apprehension comes into manifest conflict with the immediate religious consciousness's
apprehension\footnote{Let one only take account of such a prayer as the prayer to saints or to
the Madonna for intercession: \textit{Ora pro nobis}, or recall e.g. Luther's statement about
his prayer for the recovery of the mortally ill Melanchthon.}. Such a prayer without objective
goal indeed preserves prayer's form, is formally a prayer, but in such a way that, through the
consciousness of its merely subjective working, \emph{prayer as prayer is recalled in the very
prayer: it is in reality not faith's, but unbelief's prayer}. When the concept of prayer is
taken, without distortion, in the meaning in which the religious consciousness takes it, then
it is apprehended as that which \emph{changes God's will}, and this is understood more closely
in such a way that \emph{prayer determines God}, but determines God to an utterance of
omnipotence. The contradiction lying herein the religious consciousness as such does not see:
\emph{prayer's omnipotence is concealed by God's omnipotence}.

2.~That which is prayed for is \emph{the miracle and only the miracle}. The concept of prayer
in the general religious consciousness's apprehension includes that there is prayed, not for
what as an ethical task rests in our will's decision, or for what follows from existence's
universal law or from an immutable divine decision from eternity, but that there is prayed for
what has its ground only in an absolutely free divine will. It is the almighty will's working
that prayer invokes, and it invokes the will as the one that is not bound by any presupposition,
neither by one existing outside God nor by one given in God's essence-determination, not even
by the will's preceding determinations. But the thus apprehended almighty will is the
miracle-working will. And the miracle that is expected of the will is as miracle always a
miracle \emph{``in the strictest sense''}\footnote{cf. Martensen's utterances: Om Tro og Viden,
p. 112.}. \emph{There are just as little degrees in the miracle as there are degrees in
omnipotence; for the miracle is omnipotence's expression}. There can be a difference between
what is more or less \textit{mirabile}; but inasmuch as this \textit{mirabile} is seen as
\textit{miraculum}, and thus it is seen by faith and by faith alone, the difference is
abolished: then omnipotence, which is subjected to no quantification, is seen as sole-working
cause.

But the God to whom prayer points: the God determined by the form of single personality, whose
almighty will transcends his essence and thereby all essence, all law; this God, in whom the
immediate religious consciousness finds satisfaction for its willing and for its representing
inasmuch as both are determined by fantasy\footnote{For thought there is here no reconciliation,
because the spirit's development still remains in the sporadic, the unworked-through.},---this
God is precisely the God who, when the representation of him is made the object of a science of
faith, must yield a concept of God like that set forth by Martensen. In the carrying-out
nuances can occur: the essential fundamental features must be the same where theology with
consciousness holds fast its task.

The latter's task as a science of faith was at an earlier point determined in such a way that
it had to apprehend revelation and God's essence and working in such a way that they were drawn
out of the purely incomprehensible's sphere toward the rational, that there accordingly forced
itself in a gleam of knowledge's and grounding's necessity, while yet the content's
supra-rational character was not given up, and knowledge, which was subordinated to the
revealed's norm, did not with its forms bring its essence-determinateness into the object.

The investigation of how this task is sought to be more closely determined and carried through,
and what influence in general its carrying-through must have on theology's treatment of its
object, now becomes the second main point with which the critique of theology's standpoint has
to occupy itself. The critique of the concept of God acquires for its supplement the critique
of the concept of science that lies at the ground of the apprehension of a science of faith's
possibility and peculiarity.

\begin{center}---\end{center}

\section{Critique of the Concept of Science}
\label{sec:science-concept}

At the beginning of this section (p. 44 ff.) it has been briefly indicated how Martensen
seeks to demonstrate the possibility of a science of faith as an organic unity of faith
and knowledge; how he determines its task, its form, and its scientific character. It has
at the same time been shown how his apprehension of that which is peculiar to the ``higher
science'' and of this science's right to the primacy rested on his apprehension of the
fundamental relation in the human essence and of the concept of God.

Through the critical investigation of the theory's psychological foundation and of the
``ethical concept of God'' the premises for the judgment on such a ``higher'' science have
already in all essentials been won. When the will cannot be asserted the aimed-at primacy
in the human essence, and when in the divine the will's transcending of the essence's and
knowledge's necessity, by which the ethical concept of God was to acquire its peculiarity
and its higher rank in comparison with the logical, cannot be carried through, then that
which is conditioned thereby: a science resting on the free divine personality's
supernatural revelation, in which knowledge's necessity was to be limited, its laws
corrected through the will, is impossible: the logical asserts its right and the divine's
communication to man becomes one determined by the essence's necessity, apprehensible
through the forms of thought. From the untenability of the postulated ``higher'' science's
principle follows as a necessary consequence this science's untenability. This will confirm
itself when we, after having supplemented the earlier-made determinations of the higher
science's essence by a more exact determination of its relation to science in the general
meaning, more closely illuminate the given determinations and by some examples briefly
demonstrate the results to which the higher science leads and must lead. These examples
will then show the correctness of the general formula for the higher science's method which
we set up, and which most clearly reveals the contradictions and ambiguities that are
inseparable from theology as theology and that deprive it of the character of science in the
proper meaning.

\subsection{The Relation of the Higher Science to the Merely Rational}
\label{ssec:higher-science}

The relation between the ``higher'' science and science in the general meaning is
determined\footnote{Martensen: Om Tro og Viden, p. 94 f.} by the opposition, drawn from
\emph{Schelling}, between the \emph{positive science, the positive philosophy}, and the
\emph{merely negative}. This latter is determined as the seeking philosophy, which, through
a continued negation of the merely relative principles, seeks \emph{the true Absolute,
higher than which nothing can be thought or be}, while the positive philosophy takes its
point of departure from the \emph{thus-found principle}, developing its positive richness in
so far as this lets itself be cognized through its factual revelations. The negative
philosophy, whose principle is a merely logical one, is to be the \emph{abstract necessity's,
the impersonal's} philosophy, whose final goal is, from necessity to come to freedom, and
whose categories contain only the negative condition, \textit{conditio sine qua non}, for
the cognition of truth, but do not give us the actual living truth; the \emph{positive
philosophy}, whose principle is a more than a logical one, is to be \emph{freedom's, the
personal's} philosophy, which more closely determines itself as \emph{history's and
revelation's philosophy}, inasmuch as it investigates in history the personal God's
testimony, and from this standpoint it cognizes also nature, which now will show itself in a
new light. \emph{The negative philosophy cannot give an explanation of reality}, since it can
only seek the principle that is the all-explaining; this, on the contrary, the positive
philosophy, the philosophy for life, reality's philosophy, can do.

The positive philosophy is determined namely as \emph{a science of experience, empiricism},
but as \emph{metaphysical empiricism} (cf. Schelling), which through intuition as organ sees
in the working the working cause, in the fact the idea, in that which reveals itself in the
midst of time that which was before the foundation of the world was laid, and seeks to
develop the intuition into concept. It begins, inasmuch as it sets freedom's God as the
truth, with a great affirmation, a great fundamental assumption,---it hereby goes beyond the
logical and the logical's necessity. But it is to furnish a \emph{progressing proof} for its
fundamental assumption, inasmuch as it demonstrates and cognizingly absorbs itself in the
proof which the willing and acting God himself continues to furnish for his existence in his
continued relation to mankind. \emph{Only an uninterruptedly progressing proof and not a
concluded one} is namely to be suited to the \emph{personal God's existence}. Just as we
demand that human individuals who are important to us shall continuingly give us new proofs of
their existence, so we demand it also with respect to God. We demand that the Godhead always
comes nearer to man's consciousness, we require that it not only in its workings, but itself,
becomes an object for consciousness, which can yet only be reached by degrees. The full proof
for the personal God's existence must therefore be furnished through the whole of reality; it
goes back to our race's past and reaches out into its furthest future. \emph{In freedom's
philosophy there can therefore be no question of any concluded system} because \emph{reality's
system is not concluded}, and freedom's God, who has introduced himself into history, is not
only the one who is and who was, but is also the one who comes\footnote{Martensen, op. cit.,
p. 95 f.}.

\emph{To this higher science}: the science of revelation, the ethical science, the science of
the good in the eminent meaning, the science of God, the one Good, in his free
self-communication to man, theology relates itself in such a way that \emph{the former
becomes the universal science, theology the central science}, which exclusively absorbs
itself in the One, in God's kingdom as such.

\subsection{Critique of the Attempt to Ground the Possibility of
  a Science of Faith}
\label{ssec:faith-science-critique}

\emph{The critique of this concept of science must first illuminate the attempt to
demonstrate the possibility of a scientific knowledge of faith's content.}

In what turns of phrase this attempt is made has in essentials already been indicated at an
earlier place\footnote{See above, p. 44 ff.}. What is there stated is more closely determined
by the statements that \emph{revelation does not indeed want to give us a scientific system of
concepts}, but that this by no means is to say that it has wanted to teach us nothing about the
highest questions whose solution metaphysics seeks; for even though it does not give us
\emph{abstract concepts}, it gives us a \emph{comprehensive total intuition that embraces
heaven and earth, nature and history, the past, present, and to-come}\footnote{Martensen: Om
Tro og Viden, p. 26 f.}; and it ``has its own metaphysics, of which it will give up not the
least.''

It will at first glance be evident that the sought proof for that scientific knowledge is
reached only through \emph{subreptions} by the aid of ambiguity in the use of the concepts. It
is the concepts \emph{knowledge, truth, thought, and science} that are used in a double meaning
or in an effaced meaning. The fact that the believer must \emph{know what he believes}, i.e.
\emph{be conscious of what he believes}, is made into his having to be able to know it in the
meaning of having \emph{a scientific cognition of it}. When faith was not itself a knowledge,
it is said, then the believer would not be able to confess his faith and not be able to reject
and combat the erring doctrines. Here is again the same exchange of concepts. The believer must
be conscious of the \emph{faith-content}, it must be for his consciousness in unfolded form in
order that he can confess it; but he can be conscious of it as that which precisely \emph{cannot
be known scientifically}, i.e. comprehended by a rational principle, because its metaphysics,
that is, the metaphysical expression to which its determinations can be reduced, is a
contradictory, a contradictory combination of the metaphysical determinations valid for
thought. The believer can, out of that consciousness, combat opposite doctrines in the way in
which it has factually happened: either by demonstrating inner contradictions in the opponents,
or by arguing out of the religious presuppositions not acknowledged by the opponents. ---It is
further said that when faith was not already in itself a knowledge, then \emph{it was not the
faith of truth}. Faith's truth designates for the religious consciousness that faith's content
is the \emph{true}, i.e. \emph{actual God}, in opposition to the \emph{false}, i.e. \emph{unreal
gods}. But the concept \emph{truth} is transferred from this meaning to \emph{the meaning of
theoretical cognition in the form of science}. The consciousness of the \emph{Christian truth}
in the meaning of that which really is the Christian is by a \emph{subreption} made into the
\emph{Christian idea of truth}, and through the determination \emph{idea} there is reached by a
new \emph{subreption} \emph{the scientific}. In the apocryphal writings' representation, taken
over from the later Greek philosophy, of \emph{wisdom}, ``the divine thought that before the
creation of the world played before God's countenance,'' and which has stamped itself in
revelation, there is sought a possibility for the believing consciousness to be able
\emph{thinkingly} to search out revelation and scrutinize, not only for connection, but also
for ground; there is talk of a divine \emph{thought of revelation and wisdom that in the
believing consciousness grounds a unity of knowledge and conscience}. But when God's essence is
determined by the will transcending all knowledge's necessity, then \emph{wisdom} relates
itself to the ``eternal counsels,'' which, as proceeding from a will grounding only in itself
and therefore groundless, are unfathomable. Therefore the divine wisdom becomes that which
surpasses all human knowledge; revelation reveals the counsels and the counsels' subject as that
which is \emph{incommensurable} with knowledge, as that of which there can only be known a That,
not a How. Faith's content can with inner consistency be linked together, but the principle
becomes the \emph{inaccessible to thought}, therefore \emph{the connected as a whole becomes
unthinkable}; only the \emph{formal connection in the comprehended} becomes object for thought.
That the faith-content does not become object for \emph{a proper knowledge}, i.e. \emph{for a
comprehension}, is also partly indirectly conceded through the determinations that are given of
the science of faith's essence, partly it follows from the results of the attempt to present
such a science.

The determinations of \emph{the science of faith's peculiarity} are so constituted that they,
while they show an uncertainty in the apprehension of its relation to \emph{the purely rational
science}: an uncertainty that is precisely grounded in the impossibility of limiting the
latter's domain, at the same time, precisely inasmuch as they are to secure for the science of
faith a higher rank, \emph{draw it down below the sphere of knowledge and science, deprive it of
the right to the name of science}.

\subsection{Critique of the Determination of its Relation to
  Rational Science}
\label{ssec:rational-science}

As regards the apprehension of that relation, the negative philosophy is first determined
in such a way that it not only \emph{seeks}, but \emph{finds}\footnote{cf. Martensen: Om Tro
og Viden, p. 94: ``the thus-found principle''.} the true principle, higher than which
nothing can be thought or be, hence reaches \emph{cognition's highest, the in-truth
Absolute}, and embraces \emph{the whole domain of metaphysics}; while the positive
philosophy is determined in such a way that it had to coincide with \emph{the philosophy of
nature and of spirit}, so that accordingly both, the negative and the positive, would fall
within \emph{the purely rational philosophy's sphere}, as parts of it. Then, however, \emph{the
purely rational science}, which is identified with \emph{the negative}, is sought to be
distinguished from \emph{the positive}; but the separation acquires only a semblance of
validity when one arbitrarily abstracts from that which is set as the negative philosophy's
conclusion or breaks off where the goal is touched without developing the concept of the
Absolute. Thereupon there is postulated as condition for the positive philosophy's beginning
``a great fundamental assumption.'' But the negative philosophy is indeed already to have
``found the principle''; there is accordingly needed no separate new fundamental assumption.
By these determinations the positive philosophy is therefore not in reality withdrawn from the
purely rational's compass, and when the ``higher'' science's principle, which, in opposition to
the merely logical-physical truth, is designated as the absolute truth, is determined as ``the
Good itself in its revelation for thought as the highest reality,'' then this points in the
same direction. \emph{But thus the positive science would accordingly draw its scientific
character only from the unity with the negative, purely rational one.}

\emph{The involuntary concession of the rational philosophy's universality} is with necessity
grounded in the nature of the case. \emph{Science is science only as system, as the according
to its essence universal thought's organic unfolding. A limitation of science by an other, which
is yet to be an object for a cognition, is an impossibility, because it is a negation of
science's essence}, and because this is determined by the indivisible self-consciousness's
universal nature.

This conclusion one will from Martensen's standpoint reject by saying: it is not science as
such that is limited, but only the lower science, and inasmuch as that which limits this is the
higher science, then \emph{the lower science is precisely in this limitation supplemented into
science's true system, embracing all given reality}.

Hereby we are led to investigate the determinations by which the postulated science was to
assert itself as the higher.

That which in reality is to justify the distinction of the positive science from the rational,
and the former's higher worth, is the postulate of a new kind of cognition, a free cognition
determined by the will, which is first to make itself felt in that ``fundamental assumption,''
which accordingly in reality is not to become an attachment to the purely rational philosophy's
point of conclusion, but a setting of a new, heterogeneous point of departure, in virtue of a
subordination under faith's and revelation's content, of a ``free relation of dependence'' to
revelation. But this higher cognition that is postulated will show itself in reality not to be
cognition, but only \emph{an activity of fantasy determined by craving}.

The new form of cognition, which is set as a higher than the merely logical thinking, is to be
the higher and make the positive science built thereon a higher, because it is \emph{a
cognition of the free personal God's revelations}. But the concept of a divine personality to
which there is pointed is already illuminated, and it has been shown that it has only the finite
single personality's content, whereby there is brought in a determinateness of finitude that
makes that proceeding from this source a lower than that derived from the idea. One is placed
before the dilemma: \emph{when nature and history are to be cognized as God's revelations, then
they must be seen as grounded in God's essence, and this can then, as all existence's absolute
principle, not be determined by personality};---\emph{is, on the contrary, the determination of
personality maintained, then God becomes the absolutely incomprehensible}, and nature and
history cannot by thought be apprehended as God's revelations. \emph{Personality is precisely to
be that which does not go up into the concept}, that which becomes only object for intuition,
feeling, adoration. The \emph{absolute personality} would be the most personal of all
personalities, the absolutely singular being, which had no universal predicates at all, and
therefore no points of attachment for thinking. ``Indeed God acquires, when he is determined as
personal, a determination that also belongs to me, but its absoluteness is its singularity, its
distinction from me. Concerning another finite personality I can, on account of its kinship with
me, indeed guess and surmise something; but about the absolute person I can, since here even the
basis for the presupposition of a possible likeness with me in thinking and acting falls away,
only dream and fantasize.'' (Feuerbach.) ---What is revealed by the free, personal God is to be
able to be cognized by man only \emph{by freedom}, and this \emph{relation of freedom}, which is
opposed to \emph{the relation of thought-necessity}, is then to designate the new cognition's
and the new science's higher worth. But this free cognition, which is to be a higher cognition,
must be thought to come into being by a break, set by the will,---an actual conceptual
determination of the relation Martensen does not give---; so that thought-necessity is
interrupted by the act of will, which sets a new beginning. \emph{The will that sets this then
becomes a will different from cognition}, but for such a will we have no other designation than
\emph{arbitrariness}, and the cognition that is to rest on this act of will becomes such a one as
is borne by ``the will's plumage'' fantasy, and lacks every corrective, every regulator. Its
``freedom'' is a lower than thought's necessity; for the thought-necessity which in the rational
science is to be the sign of the latter's lower worth is precisely the expression for thought's
determinateness by its own essence, its own law; it is therefore \emph{the mark of all proper
science}, and it is a \emph{misunderstanding} when it is restricted to science's most abstract
and contentless forms. The more concrete and richer science's content is, the more difficult
becomes the assimilation that gives the object the form of thought-necessity, but therefore this
assimilation becomes just as fully the task, when the concrete, only successively appropriable,
is not confused with that which, in accordance with its irrational nature, is incommensurable
with thought and absolutely unappropriable for it. And in these domains we get not an abstract
thought-necessity, but \emph{the concrete thought's necessity, which is compelling for thought
because it is thought's own essence}.

The new cognition's and the new science's higher worth is finally motivated by the fact that they
are to give a \emph{cognition of reality}, which the negative philosophy is not to be able to do.
Here, however, the same confusion of concepts is present that was demonstrated at an earlier
point. The thought-necessity that is my thinking essence's necessity has \emph{existential
certainty} for me and is precisely the highest assurance, while the certainty of outer existence
can always become doubtful. But the existence Martensen seeks is precisely an existence in analogy
with the outer, sensible reality's, such an existence as becomes object for intuition. To such a
one there is also pointed when there is talk of the higher science's personal principle and this
principle's revelation. Precisely inasmuch as there is demanded a continuing proof for this
personal God's existence, he is drawn down into likeness with the single individual, whose
existence as finitely determined can be interrupted. The higher cognition, which is designated by
the contradictory concept \emph{metaphysical empiricism}, thus becomes in reality only \emph{an
empiricism that relates itself to a finite, sensible, fantasy-determined object}, and the cognition
of this is to be \emph{cognition of reality}, while the cognition of the idea and of that grounded
in the idea is not to be it. And yet it is surely evident that when it is demanded of the higher
science that it shall relate itself to existence, then existence, as set in thought, even though
this thought is a thought determined by will, \emph{must always become existence's thought, hence
in its truth the idea's being in thought}. Is existence to be something other than the existence
the ideal has in my thought, then it becomes \emph{sensible existence}, and thought, which is to
take it up, becomes only \emph{formal consciousness of a sensible relation, an experiential relation
in the general meaning}. When emphasis is laid on God's revealing himself in a single personality in
order thereby to be immediately revealed, then it is forgotten that such a single personality, in
its separation from the other human individuals and in its being in man's form, yet becomes only a
\emph{mediate revelation} and a revelation through a working of God. \emph{God's revelation of
himself can only be thought as an essence-communication to the spirit}.

When thus everything that was to assert for the positive science a higher rank precisely points in
the opposite direction, and when this higher science, inasmuch as it is not to be able to give its
object the form of necessity, inasmuch as it is not to reach an actual comprehension of it, and
inasmuch as it is to bow beneath an authority that is not thought's, and cannot be higher than
thought's, because, as has shown itself, the highest in the spirit's world finds its expression in
the logical, outside which only the alogical falls---precisely lacks that which constitutes
science's essence and form, that which characterizes science as science, and shows itself as
resting on a foundation that is lower than the concept, then it loses \emph{the right to bear the
name of science and still more to be the higher and ruling science}. It preserves of that which
designates science only this, that from the principle coherent consequences are drawn; but the
principle, which bears the consequences and cannot be transformed by a reaction of the consequences,
becomes an uncomprehended. \emph{Science in the true sense and ruling science becomes the science
that comprehends the Absolute in its true infinity as existence's principle of reality}, and this
science is \emph{philosophy}.

\subsection{The Method of the Science of Faith}
\label{ssec:method}

The judgment that has here been passed with respect to the ``higher'' science: the
science of faith, confirms itself completely when we consider its results. Already above
it has, with respect to essential points in the doctrine of faith: God's personality,
creation, and the miracle, been demonstrated how the attempt to make them the object of a
cognition failed, and how the attempted rationalization, consistently carried through,
would set something quite other in place of the religious doctrines. The same would be
able to be demonstrated with respect to any dogma whatever. Everywhere the development
must be carried forward by postulates, everywhere the conceptual determinations must be
used in a double sense, everywhere the result must become an immediate union of the
contradictory.

The method will, when everything concealing is removed, show itself to be the following.
In order to give the movement of thought a point of attachment, there is a setting out
from determinations drawn from nature's and spirit's reality, but which, inasmuch as they
are to be applied to faith's object, are arbitrarily, in part fantastically, transformed
in order that the original determinateness can be forgotten and the application become
possible. That thus transformed now becomes the ground for new determinations won by
postulates, and inasmuch as the consequences of the rational are arbitrarily broken off,
or a quantification to which no concept can be attached is applied, the faith-content that
is to be grounded for thought is brought forward through an illusory grounding, in reality
only through postulates, and it is concluded therewith that the two opposite elements that
meet in the science of faith both find their expression, but in such a way that the
expressions neutralize each other in indeterminacy, inasmuch as the rational element is
paralyzed by the faith-element, but yet by its presence affects this. The subject in the
proposition that is set up as result is determined by predicates that deprive it of its
content and leave only an indeterminate, vague representation for consciousness, which then
precisely in the indeterminacy, in the lack of the determinate, clear limitation, imagines
itself to have something deep, something speculative, something infinite.

\subsection{Its Results}
\label{ssec:results}

When we seek determinate examples of this method, whose necessity lies in the task that
is set, then we can e.g. find them in the development of the dogma of the Trinity in the
Dogmatics and in the work on Faith and Knowledge. There one will find all the method's
moments united: a setting out from rational determinations: self-consciousness's and
community's determinations; a transformation of them by which they lose their specific
meaning (self-consciousness's ideal doubling becomes a multiplication of personality
positing several self-consciousnesses; community is to be a person); an arbitrary
advancing through postulates (e.g. that the person as absolute, divine person is to be a
plurality of persons; that the persons, in spite of their difference, are yet to be God;
that the unity, although a unity of oppositions, must not be a higher unity); an arbitrary
breaking off (at the threeness of the persons, while the premises lead to an indeterminable
plurality); an immediate linking together of the mutually abolishing. \emph{We can find
them in the development of creation's relation to cosmogony}, where first the representation
of creation is illusorily rationalized by the concept of cosmogony, but then this is
factually abolished by being brought in under the one creative supernatural beginning.
\emph{We can find them in the development of the doctrine of the Fall}, which one is at once
to think as a both unconditioned and conditioned divine counsel, that is, as a counsel not
only determined from eternity, but also as determinable by the creature's freedom, as a
coming-to-be counsel that determines itself to historical movement of life; in other words:
not only as a logical, but as an ethical counsel, the holy will's counsel, in which there is
still something undecided, but which yet in its essence does not cease to be unconditioned.
\emph{We can find them in the development of the relation between grace and freedom}, where
on the one hand, inasmuch as grace, in distinction from the mere omnipotence, is not to work
irresistibly, but to respect human freedom, there is asserted for the freedom of the will a
significance, but this is then again taken back by the added determinations, inasmuch as the
essential freedom is determined as the co-created grace that comes to breakthrough in the
natural will, and is to be regarded as grace in nature, while the merely natural, human will
can only offer resistance. We can finally, in order not unnecessarily to multiply these
examples, which are only to supplement those already analyzed by others, point in conclusion
to the Dogmatics' last section, to \emph{the doctrine of the last things}, where eternal
damnation is maintained alongside God's counsel for the blessedness of all, where there is
ascribed to the damned an unlosable possibility for the good, but in such a way that this
possibility is to be absolutely cut off from all conditions of reality for its development,
and where fantasy creates the representation of a ``spatiality'' within spirituality and
inwardness, of an intermediate state after death with an ``intermediate corporeality,'' etc.

As soon as the ambiguous and indeterminate in the science of faith's designations is removed,
then the faith-determinations and the knowledge-determinations show themselves here, instead
of being combined in organic unity, as disturbing and abolishing each other. \emph{A science
of faith is accordingly here factually not given; whether it can be given at all depends on
whether this theological attempt expresses something universally valid for theology or not.}

\subsection{What is Typical for Theology in the Criticized Conception
  of the Science of Faith}
\label{ssec:typical-faith-science}

This becomes easy to decide when the relation between the form of science that theology
as a science of faith aims at, and the content that is to be taken up into this form, is
considered. \emph{Thought's forms are inseparable from a content of thought, determined by
its essence}; if accordingly the faith-content is really to be formed by these forms, i.e.
\emph{go up as formed-through in these forms, not only partially and in an external way be
determined by them}, it must in the proper sense be able to be a \emph{content of
knowledge}. Can it not be this, such an application of the forms as that by which science
is conditioned will be impossible. Science, as comprehending cognition of an organically
unfolded whole, is conditioned by the fact that \emph{its content in its wholeness is a
rational} one, i.e. \emph{a thought-necessary forming-out of a principle determined by
thought, homogeneous with thought, not merely something formally linked together by
thought-determinations, heterogeneous with thought}.

To the question whether the faith-content can stand in such a relation of unity to forms
and content of knowledge, the answer is already intimated at the beginning of this section,
and by the later developments new moments are given for the answering. It was emphasized in
the first place that the religious consciousness's distinction between the transcendent God
and man had to lead to a relation of opposition, inasmuch as reflection determined the
relation. Faith with its transcendent, supernatural object and its transcendent supernatural
principle must distinguish itself from knowledge's content and knowledge's principle; for
\emph{knowledge, thought, is the determination of human nature, its forms accordingly natural
forms, its content determined by these forms. As expression for human nature, thought cannot
take up the supernatural}; for what can be received is determined by the receiving organ.
\emph{A reception of a qualitatively higher would presuppose a re-creating of the organ}. The
distinction between the supernatural object and the natural organ has then, by the later
developments, acquired the more determinate form that \emph{the faith-content becomes the
content of the free, personal, by no necessity of essence or of knowledge bound God's
supernatural revelation}, a content that concerns God's essence surpassing thought, the works
of omnipotence, the divine counsels. \emph{The supernatural revelation is for the religious
consciousness a miracle: the reception of that communicated by the revelation is a miracle},
by which there is set a new beginning of cognition that does not have a connection and an
essence-unity with the preceding, in-itself concluded knowledge. \emph{The content of
revelation therefore becomes an incomprehensible, an opposite of the content of knowledge}.
That the opposition would just as fully come forward when one began with a quantitatively
sounding difference between human knowledge and that which abidingly surpasses it has been
demonstrated in the foregoing.

The religious consciousness must thus determine the relation as a relation of opposition, but
in this it agrees with the scientific consciousness, which cannot acknowledge any other
principle for science than the knowledge determined by its own essence's law, and which sees
in the positive religion's content that which by the representational form is determined as a
finite, and as such stands in opposition to the content of knowledge determined by thought's
infinity.

Theology shows itself accordingly, inasmuch as it would found a science of faith, to acquire a
task contradictory in principle. For its task is indeed, inasmuch as it has that supernatural
content of revelation, set by the religious consciousness from whose standpoint theology would
reflect on religion, in opposition to human knowledge and its content, for its object, to make
it a content of knowledge. If an attempt is to be made to solve this task without the
contradiction at once becoming striking, then there must everywhere be resorted to what can
make the relation untransparent: arbitrary restrictions and double-sounding determinations must
be applied everywhere where boundaries are to be indicated, hence precisely where the opposition
would come forward, and the two opposites that are to be united must both receive a transformation
conflicting with their nature. \emph{This is to say, in other words: on the one hand the concept
of knowledge and science must be distorted, while a semblance of their subsistence is preserved;
on the other hand the concept of the God whose essence is will, whose working is the miracle, must
acquire a distorting admixture of an apparently rational:} i.e. \emph{the rational must be made
half irrational, the irrational half rational.}

As regards the first, the concept of science, to be the cognition of that which with rational
consistency is derived from rational presuppositions, is extended to designate a derivation of
formal consequences from an anti-rational presupposition. \emph{Science accordingly acquires a
double meaning. The name is preserved as a common name for the essentially different, indeed in
essence opposite, determined by an opposite principle.} The science of faith's knowledge becomes
a consciousness of that \emph{formal connection, not a comprehension of the whole out of the
principle}. And inasmuch as theology, in the consciousness of the difference between the
thus-transformed science and science in its general meaning, would determine the former's
peculiarity and, instead of that science's freedom which is the expression for its
determinateness by its own principle, by thought's necessity, ascribes to it a dependence on a
principle heterogeneous with thought, and yet ``correcting thought,'' then it must, by a
postulate, in order to preserve the semblance of what characterizes science, determine this
dependence as free dependence, and by an ambiguity use the mediate, scientific cognition's
relation to its immediate point of departure in order to bring about a semblance of likeness. But
along neither of these ways is the aimed-at goal reached. That postulate of freedom cannot be
justified. Science, which relates itself to a \emph{real object}, can set itself in a free
relation of dependence to this, and must do so in order through the assimilation of the object in
its reality to become master of the objectivity; but it makes itself dependent only because it in
the object acknowledges \emph{a rational}, hence such a one in which its own essence is affirmed.
But knowledge's subordination under the faith-content's authority is the subordination under the
heterogeneous, hence not science's true dependence on reason's universal essence-law, but an
unfree dependence abolishing science's essence. \emph{Knowledge's subordination under the foreign
authority must take place by an act of an activity different from knowledge, by an act of will not
determined by knowledge}; therefore all theology must consistently accentuate the will as the
theological science's organ. ---The ambiguity in the appeal to the immediate lies in this, that
inasmuch as it is emphasized that all science rests on an immediate point of departure, and inasmuch
as it is concluded therefrom that there is accordingly no contradiction in a science resting on the
faith-relation's immediacy, but that on the contrary it must be demanded that faith be the first,
the primitive truth, science the derived, the concept: the immediate is at once used of that which
is object for cognition's phenomenologically first form: sensation, experience, and which is to be
converted into a comprehended, deduced and justified by being comprehended, and of that which is
object for faith in this word's proper---not in the effaced Jacobian---meaning, and which cannot and
shall not be converted into knowledge, or justified by thought.

It is thus only by ambiguities and distorting transformations of the decisive concepts that the
relation of opposition can be concealed. In whatever forms the science of faith can be presented,
it must in reality constantly, because it rests on a ground absolutely incommensurable with
thought, unconvertible into knowledge, lack that which constitutes knowledge's essence and form; it
preserves only the name of science, a name that loses all its content by being supposed to comprise
the in essence opposite, and gives only a semblance of science by partially preserving the merely
formal. \emph{The rational in the science of faith becomes the formal connection; an only apparently
rational becomes the grounding}, whose necessity becomes illusory, on account of the principle's
groundlessness for thought. Science according to the concept's true meaning, true science is only
one: that which draws its content from the idea, the universal science, which by its principle is in
unity with the personal life's principle. It is science without particularizing predicates, not
``Christian'' or ``believing'' science, but science alone.

To theology's distortion of the concept of science corresponds the distortion of the positive
religion's apprehension of that which is peculiar to the Godhead's essence and working: of \emph{the
miracle}, the expression for the almightily willing God's form of working and form of revelation.
The miracle, which expresses the absolute opposition to thought and to the law of existence by which
thought is borne, must apparently be made accessible to cognition when a knowledge of faith is to
become possible. That is to say: the concept of the miracle must be used ambiguously, so that it at
once designates what is a \textit{mirabile} and what is a \textit{miraculum}, so that there is applied
a \emph{quantification} that draws it toward the merely wonderful, astonishing, and through this
brings about an illusory continuity, a ``teleological'' continuity, with the universal and natural.
But with the distortion of the miracle \emph{faith's essence} is distorted, for in the miracle the
positively religious's peculiarity is concentrated and with it all the dogmas' peculiar
determinations stand in indissoluble connection. To the distortion of the concept of science there
corresponds accordingly in theology the distortion of the faith-content.

\begin{center}---\end{center}

\section{Final Result}
\label{sec:theology-result}

The final result of these investigations is accordingly that the defects in scientific
respect which were demonstrated in Martensen must become typical for theology in general,
and that the discord between the faith-determinations and the knowledge-determinations,
which was factually present in Martensen, must, inasmuch as both determinations are seen
from the standpoint of cognition, become inevitable for any theological attempt whatever.

There then remains only the possibility that the relation could be changed if the
standpoint were changed and the two territories: that of knowledge and that of faith, were
seen as separated in principle, and as, by this separation, which as it were made the
content incomparable, withdrawn from the collision. In Martensen, in spite of the polemic
against this view as it is stamped in Prof. Nielsen, such a separation in principle is
present. Inasmuch as he lets revelation ``from first to last be practical,'' and inasmuch as
he accentuates the will's absolute priority and its difference from cognition, he in reality
gives Nielsen the right with respect to his fundamental separation. That the will is set by
Martensen as highest unity and constitutive principle does not remove the principled
separation; for it is not demonstrated whence comes the difference that constitutes
cognition, how this accordingly proceeds from the constitutive.

From the critique of theology's unachievable attempt at unification we then proceed, in order
to test that possibility, to the critical investigation of the theory that, precisely by
positing the principled separation of the territories of faith and knowledge, would vindicate
both in their validity; and we occupy ourselves first with the theory in the form in which it
has been presented by its founder: S.\ Kierkegaard.

\begin{center}---\end{center}


%% ============================================================
%%  THE PRINCIPLED SEPARATION (pp. 118--224)
%% ============================================================
\chapter{The Principled Separation of Faith and Knowledge in the
  Interest of Reconciliation}
\label{ch:separation}

\section{The Main Features of S.\ Kierkegaard's Theory}
\label{sec:kierkegaard}

In the historical introduction the main features of S.\ Kierkegaard's
conception of the relation between faith and knowledge have already
been indicated, with particular attention to the comparison with
Nielsen's theory. At this point these main features are to be
developed more fully.

Kierkegaard's conception is developed through a polemic against
speculation, particularly as it appears in Hegelian philosophy; but the
polemic is directed against speculation because Kierkegaard sees in it
the strongest expression of the whole age's misdirection: having
forgotten to exist, forgotten what inwardness is. For speculative
philosophy the highest task of the human spirit was to transpose itself
into a medium of eternity, to live a pure life of eternity, and this
task was to be realized through pure thought, in that the individual's
thinking of the Absolute was posited as one and the same with the
Absolute's thinking of itself, with the absolute thought existing in
the element of eternity. The life of the individual, which found its
exhaustive expression in logical thought, through its thinking of the
Absolute's eternal life-process encompassing all truth and reality,
became a moment in this eternal process, and was freed from the
determination of finitude. Against this conception of human life and its
task Kierkegaard protests. Against speculation's volatilization of
existence, against its forgetting of existence's claim upon the
individual who is subject to the demands of the ethical, against its
overloking of the conditions of existence, against the illusory movement
in a sphere of abstract eternity, Kierkegaard accentuates existence,
gathers everything around the question: how do I, as existing, express
the truth? As existing, the human being is not merely eternal, as
speculation would have it, but is a synthesis of the temporal and the
eternal, of the finite and the infinite; as existing, the spirit that
is eternal according to its nature is subject to conditions of finitude,
is in becoming; as existing, the thinker must in all his thought think
along with it that he is existing; he must set all his thought into
becoming. But becoming is the alternation between being and non-being;
through becoming the determination of negation is brought at every
moment into the spirit's being and thinking. From this it follows with
necessity that the human being is prevented by existence from reaching
the positive, secure rest in eternity; for since the existing being
cannot hold fast the infinite and eternal---which is the only thing
certain---in the element of infinity, but must hold it fast in the
element of finitude, in existence, the relation does not become one of
secure certainty, but must express the determination of negation that
lies in becoming, and the eternal can therefore be had by the existing
being---for whom it converts itself into something future---only in the
form of uncertainty for cognition. Herein lies, then, that objective
knowledge in none of its forms can attain such security that it could
become the foundation for the infinite decision for the subject, for a
decision with respect to what concerns the subject's deepest need.
Sensory certainty is a deception, historical knowledge is an
approximation, mathematical cognition is an abstraction, speculative
thinking likewise an abstraction and in the last instance a hypothesis
in relation to existence; upon none of these forms of knowledge can the
decision that infinitely interests the individual therefore be built.
But if the goal cannot be reached by the path of objectivity, then it
must be reached by that of subjectivity: in a relation of subjectivity
the eternal must be grasped, and this relation of subjectivity is the
one in which the certainty of the eternal is grasped with the passion
of faith out of objective uncertainty. Faith is ``objective uncertainty,
held fast in the most passionate appropriation of inwardness.''

The task for the individual is therefore: to become subjective; and
since what is posited as a task is what every human being might seem,
without further ado, already to be, the becoming subjective acquires
the more determinate meaning: to become what the human being according
to its nature is, to let subjectivity develop, be transformed in itself
over against the eternal. Only for such a subjectivity, which is the
developed possibility of subjectivity's first possibility, does truth
come into being. The subject, which is referred to itself, receives the
task of being infinitely concerned about itself; but this infinite
concern about oneself is the direction toward the eternal in relation to
existence; it leads toward the ethical and the ethico-religious, the
essential truth and the only essential cognition. The claim of existence
prevents the individual from remaining in the metaphysical, sets it the
task: to understand itself in existence. In being referred to a relation
of subjectivity, the subject is thus referred to the fact that it
reaches it only through a choice between this and the relation of
objectivity, and this choice is the choice between the personal and the
impersonal relation to truth. For the existing being there is no
mediation's unity of the subjective and the objective relation.
Mediation, which would connect them, leads us only back to abstraction;
it does not explain how the eternal truth is to be understood in the
determination of time by one who, by existing, is himself in time.

In the subjective relation, the how of the relation becomes the
essential. In the objective relation the direction goes away from the
individual, strives toward the objective, finished result, and its
maximum is the self-contradiction that the existing subjectivity
vanishes, while the objectivity that comes into being is, subjectively
considered, either a hypothesis or an approximation, because all
eternal decision lies in subjectivity. In the subjective relation,
objectivity drops out, and the how of the relation becomes the
decisive. At its maximum this how is the passion of infinity, and the
passion of infinity is truth itself. But the passion of infinity is
precisely subjectivity, and thus subjectivity is truth.

The relation of faith, the expression for the relation of subjectivity,
for the relation of infinitely passionate interest, is thus the
highest, because in existence it becomes the only possible relation to
truth, the form of truth in the actuality of existence. Truth in the
conception of objective thinking, as the unity of thought and being,
becomes merely a chimera of abstraction and in truth only a creature's
longing---not because truth is not this, but because the knower is an
existing being, and thus truth cannot be this for him so long as he
exists, because existence spaces apart thought and being and for the
existing spirit sets truth itself into becoming. The passion of faith,
the infinitely passionate interest, thus becomes the form in which the
eternal is grasped; in it the eternal is present in a negative way in
subjectivity, and the negative infinity becomes in existence the one
true expression for the infinite; the infinite negative resolution
becomes individuality's form for God's being in it. And faith in the
strictest sense is faith in the absolute paradox, which is constituted
not, as the Socratic, by the eternal truth's relating itself to the
existing being, but by the eternal truth itself becoming the paradox by
coming into being in time. As faith in the paradox, faith is the
infinitely passionate holding fast of what is absurd for thought, and
precisely through the repulsion of the absurd is passion potentiated to
its highest. The absolute paradox is the central point of Christianity,
and the question of how I as existing express the truth thus coincides
with the question: how do I become a Christian? The answer is: through
faith in the paradox; and what in the last instance bears the burden of
the paradox and compels the individual to hold it fast in faith despite
its absolute unthinkability is the despair of the consciousness of sin,
which does not find reconciliation in the ethical, nor in the
religiosity that still rests in the unbroken relation of immanence, but
only in the Christian, in faith in ``the god in time.''

In this theory there is thus posited a self-limitation of knowledge,
grounded in the fundamental relation of the existing spirit. Knowledge
receives its limitation---conditioned by the contradiction of existence
and acknowledged by knowledge itself---a limitation in a double respect:
in respect of certainty and in respect of concretion; and through this
double limitation there is established an incommensurability-relation
between objective knowledge and the striving and need of personality.
But under the limitation, knowledge is vindicated a justified place,
though a subordinate one, and a validity, though only a relative one.
It receives its validity and security within the sphere of the abstract
and the impersonal, within the mathematical and logical; it finds its
limit where it reaches the concrete and existential; it does not, on
account of this limit, become decisive for personality's deepest and
highest interests, for the relation to eternal blessedness; the
decisive conclusion does not become possible in its territory. Its
sphere extends to the point where the sphere of faith begins, where the
relation of faith is posited through a breach with objective knowledge;
and within the sphere of faith, thought, in entering into a serving
relation to faith, again acquires a significance through the task of
unfolding the de-

% [text continues on p. 123]

\noindent terminations in which one is to exist, and to defend the paradox as
paradox, to protect it against every attempt to comprehend it. A
comprehension of the paradox would precisely be an annihilation of its
distinctive significance, of faith's specific object; therefore theology,
which posits such a comprehension as its goal, rests on a fundamental
contradiction.

In that faith comes into being through a breach with objective knowledge,
it vindicates its distinctive territory and its principled validity; it
becomes---while it is unattainable for knowledge---inaccessible to
knowledge's attack. Knowledge and faith each receive their own sphere,
to which they relate as distinctive principles. Faith as religious faith
rests on a choice that in a higher potency expresses the same decision
of subjectivity as the choice by which the individual's ethical life is
constituted. The reconciliation between faith and knowledge is grounded
in the separation of the spheres and principles, and in the
subordination of knowledge under faith. The certainty of faith, which
victoriously rises above doubt, is the higher certainty compared to that
of cognition, which, in its relation to the concrete and actual,
continually remains afflicted with doubt.

\section{Its Relation to Other Contemporary Conceptions}
\label{sec:contemporary}

Through what has been developed, the relation between Kierkegaard's
conception of religion and faith and other contemporary conceptions
becomes clear.

Speculative philosophy of religion had let religion be a stage in the
spirit's development of cognition, which found its completion in
philosophy, in speculative thought that with logical necessity unfolded
itself from the lower form. Only in the concept did it let the content
of faith find its true form; by being converted into the form of
philosophical cognition, the content of faith received a higher truth;
for the truth of thought, of the concept, is the highest truth. In
Kierkegaard the relation of the spheres is seen from a different point
of view. Science becomes one thing, life another; the unfolding of
thought in the form of objectivity, in the medium of possibility,
according to the law of necessity, is separated from truth's
appropriation in the individual in the form of subjectivity, in the
medium of actuality, according to the law of freedom. The personal
existence of truth is accentuated as the highest; the relation of
subjectivity becomes the form of this existence; and within the personal
forms of existence the religious becomes the highest, and within it
again the relation of faith to what is absolutely inaccessible to
thought: the paradoxical, the specific determination of Christianity.
The task of thought here becomes not to assimilate the object of faith,
to dissolve it in the concept, but to discover thought's opposite: the
absolutely unthinkable, the paradox, and to hold it fast as paradox so
that one may then believe by virtue of the absurd.

Ordinary theology immerses itself objectively in learned investigations;
it aims through these to bring out what Christianity is as objective
doctrine; by the path of history and criticism it aims to secure
certainty on this point, and in this striving, which by the nature of
the case loses itself in an infinite approximation, it lets it be the
presupposition that the investigator is a Christian and that what
matters is the scientific result, whose appropriation is supposed to
follow of itself. For Kierkegaard the infinite decision of existence is
the main thing: the approximation, by which precisely the decision is
prevented, is the deception; appropriation is the task, and what is
essential is not a what but a how. The question is not what Christianity
historically is; but the question is: how do I become a Christian? Truth
is the inwardness of appropriation; truth is subjectivity.

In Feuerbach, as in Kierkegaard, the spheres of knowledge and faith are
separated, but in that Feuerbach determines religion as the irrational,
he vindicates the right of thought against faith and posits religion as
the untrue, which is to be abandoned. The opposition between Feuerbach
and Kierkegaard in the conception of religion and especially of
Christianity is therefore the greatest possible, and yet there is no
contemporary with whom Kierkegaard has more points of contact in the
determination of what Christianity is. This has not escaped him either,
and he has, precisely on account of his conception of Christianity, not
been able to wonder at it. If the individual's relation to Christianity
is a relation of subjectivity, and more precisely one in which
subjectivity is at its highest potentiation, at its most energetic
passion, then besides the believer, who in passion holds it fast as
truth, only the one who in passion rejects it can grasp its nature. The
latter, standing outside Christianity, will be able decisively and
sharply to say what Christianity is; he will be able to portray it in
such a way that the believer can in everything recognize his portrayal.
It is undoubtedly Feuerbach whom Kierkegaard has had in mind when he (in
the \textit{Concluding Unscientific Postscript}, p.\ 473), after having
discussed theology's and revivalism's confusion of the Christian, says:
``on the other side: a scoffer attacks Christianity and presents it at
the same time so admirably that it is a pleasure to read him, so that
one who is at a loss to have it determinately presented almost has to
resort to him.'' We also find an agreement in the conception on a series
of the most essential points. Feuerbach and Kierkegaard agree in
determining the relation of religion as a subjective one, even if the
nature of the subjectivity is conceived differently; they agree in the
conception of the religious standpoint as a practical one, and of
faith's opposition to thought; in the denial of the possibility of a
knowledge of a being higher than the human; in vindicating the category
of the single individual for the religious; in the conception of
religion's absolute relation as separating itself from and therefore in
reality negating the relative relations; in holding fast the equal
relation of the Incarnation to all subsequent ages; in conceiving
Christianity's relation to the world as unchangeable, as an abiding
relation of opposition; in the conception of suffering as Christianity's
mark, of dying away as the believer's wish; in the determination of the
Christian's relation to the natural, in the limitation of Christians'
sympathy to Christians, and in the close connection in which they place
the concepts: God as the absolute τέλος, and eternal
blessedness. These parallels could be further multiplied, but those
cited are the most prominent.

The relation between Kierkegaard's and Nielsen's conceptions has been
treated in the foregoing and determined in such a way that the two forms
of the theory agree in all the main points significant for the
development of the problem, while what Nielsen has added does not
advance the problem and is self-contradictory and untenable. On account
of this relation, the connected presentation of the theory's secondary
form is immediately appended to the presentation of the original, and
the more detailed critique of the standpoint is attached to the former,
to be then completed by a critique of what is distinctive to the theory
in its primitive form.

\section{The Main Features of Prof.\ R.\ Nielsen's Theory}
\label{sec:nielsen-theory}

In that Nielsen, like Kierkegaard, determines the fundamental relation
between faith and knowledge as one of separation, and more precisely as
a relation of principled separation, and precisely through the
separation seeks the possibility of vindicating both principles and both
spheres in their validity and letting them subsist as reconciled in the
same consciousness, he concentrates his theory in the proposition: the
absolutely heterogeneous principles, faith and knowledge, can, precisely
on account of their absolute heterogeneity, be united in the same
consciousness. The principled opposition, which is here determined as an
opposition between faith and knowledge, is also brought under other
determinations already found in Kierkegaard: under the determinations
life (existence) and knowledge, subjectivity and objectivity, personal
and impersonal truth, or under the determinations: will and knowledge.

% [Nielsen's theory continues pp. 127--130]

\noindent this last opposition, and through it the others, is then led by
Nielsen back to the metaphysical opposition, the opposition between the
fundamental ideas: Power and Knowledge, or the Good and the True.

The opposition between will and knowledge, between the principle of will
and the principle of thought, becomes the predominant one in the later,
more detailed communications from Nielsen's own hand and from his
expositors. Both the opposed activities of spirit are sought to be
conceived as totalities, such that in the activity of cognition a will
is present, but everywhere as serving thought, as a will-to-know, while
in the activity of will a thought is said to be present, but as
subordinate to and serving will, as a will-thought. Faith is led back to
will; the religious, which is conceived as exclusively practical, is
through the principle of will inseparably connected with the ethical,
which has its highest form in the anti-rational religio-ethical.

From the psychological starting point, the underivable distinctiveness
of will in relation to knowledge is then accentuated. Will is said
indeed to presuppose thought, but not theoretical, only practical
thinking, and in this presupposing it is said to be just as original as
thought. In its distinctiveness will has its principled significance and
its ideal just as knowledge does, and as the determinations of will and
the determinations of knowledge develop in their independence, governed
by their distinctive principle, they form two self-enclosed totalities
that do not come into conflict with each other. This lies precisely in
the absolute heterogeneity. The form of the opposition is the expression
of the factual dualism in the human spirit. In the Absolute, the
principle of knowledge and the principle of will---the Idea of the Good
and the Idea of the True---are in concrete, absolutely transparent
unity; in the finite spirit, the psychological consciousness, a dualism
is present that lets what is united in the Idea step forth in absolute
opposition. The task for the individual is precisely to comprehend the
opposition as such and thereby to win the reconciliation. Philosophy's
main task is to discover the limits of human knowledge---that is, its
own limits---and, in discovering them, to acknowledge the other
principle, which, together with the whole sphere of determinations
attached to it, is mirrored by knowledge but not assimilated, not
comprehended. Since the two spheres are contained in the same
consciousness, they must indeed touch each other and to that extent have
a boundary against each other; but the boundary is not a hampering,
restricting boundary, for it is set from within by each principle's
distinctive determinateness, and ``the absolute mystery'' prevents the
collision. While knowledge thus, mirroring its opposite, clarifies the
relation of the principles, it unfolds itself unhindered within its own
sphere, following only its own laws: it unfolds a science on the
foundation of the Idea that encompasses all the humane forms of life and
gives a true expression of the concrete, real existence. And at the same
time faith unfolds its content for the believing consciousness, whose
cognition is practical, not theoretical. The gap between religion and
philosophy is to be filled by ethics in its double form as rational and
religious ethics. The border-science is the philosophy of religion,
which as science belongs to the circle of knowledge, while, as having
the religious for its content, it is to provide orientation in the
sphere of faith. It is to show how one and the same concept can become
heterogeneous with itself depending on whether it is developed into a
``knowledge-concept'' or a ``faith-concept.'' It is to demonstrate an
all-sidedly articulated reciprocal determination, an exact mutual
correspondence, between the heterogeneous concepts, which occurs in such
a way that the knowledge-concepts serve as intellectual substratum,
while the faith-concepts step into the foreground and are developed
systematically.

As will, which has its autonomy by virtue of its distinctive principle
and in choice asserts its independence, autonomously potentiates itself
through ever more intensive concentrations, the progressive series of
stages of will appears---the stages of the ethical---whose advance is
marked by the cognitive moment being more and more absorbed and becoming
invisible, and by the regard for theory being more and more displaced by
the energy of will. The maximum of the rationally ethical is present
where subjectivity, by virtue of a decision of will, turns its back on
theory. But even beyond this standpoint one advances into the
paradoxical, anti-rationally ethical. ``The concentration is here so
decisive that all regard for theory vanishes.'' The energy of will is so
intensive that it not only takes no account of theory but even breaks
with it. Indeed not even rational ethics at its maximum can here come
into consideration, because its tasks are immanent, its faith a faith of
reason. At the height of the paradoxical ethics, the human being is
addressed in nothing but imperatives; it acts by virtue of the absurd.

The ethical is here placed in the closest relation to the paradoxically
religious, just as in general its principle is posited in unity with the
principle of religion; and the ethical is conceived as by itself
pointing toward the religious, because the Good is to be comprehended
only through the Good One. Already in the territory of the ethical there
is therefore talk of faith, but this faith is a ``rational faith''
before it, through the progressive concentration of will, becomes a
faith in relation to the paradox. Religion has its stages just as the
ethical does, and they are marked, like those of the ethical, by the
deepening of will and faith---a deepening that becomes clear when, for
example, Greek religion with its myths is compared with Judaism with its
conception of Jehovah as holy will, and with Christianity.

As the concluding stages of the ethical and the religious are reached,
the original apparent relation of coordination between faith and
knowledge transforms itself into a relation of subordination. The
relation of faith and the ideal of will become the highest; knowledge
the subordinate. There is talk of ``religion's supra-scientific
truths,'' of ``something higher than science,'' of a ``supra-scientific
teleology.'' To this corresponds exactly the conception of will as the
true constitutive element in the human being. At the concluding stage of
the ethical, where the breach with theory has occurred, the whole of
life is so absorbed in acts of will, in action, in a working upon the
task of becoming an ethical character, that no room remains for an
independent cognition, for a life in knowledge.

Through the given presentation of the main points in Kierkegaard's and
Nielsen's theory, the earlier demonstration of the agreement between
them will have been supplemented, while the real and apparent deviations
also emerge more determinately. Through the developments that follow,
the relation between the original theory and the derived one will be
still more completely clarified.

\section{Critique of Nielsen's Theory}
\label{sec:nielsen-critique}

The critique of the theory in the form that Nielsen has given it must
first confine itself to the determination of the form of the principled
opposition. But in several of the expressions that have been given for
the opposition, especially those taken from Kierkegaard, there appears
an ambiguity and an arbitrariness in the conception that render them
unfit to serve as starting points for the critique; they need only be
illuminated in order to make the relation between Nielsen and
Kierkegaard clearer.

What has been said here applies to the opposition between existence
(life) and knowledge (theory), and between the subjective and the
objective. As for the first opposition, the concept of existence is
taken in an arbitrary double meaning, upon which a sophistical inference
is built. Existence is taken, first, in order to be asserted as a
universal presupposition and necessary condition for every concrete
determination---thus also for knowledge---in the meaning of individual
being in general; and then, in order that through the concept of
existence the practical may be claimed a priority and a preponderance,
in the meaning of volitional being (``existence is practical''). Without
this subreption the argumentation, which is supposed to lead to the
conclusion that it is possible to place existence highest and yet retain
theory, but impossible to place theory highest and yet retain existence,
and further to establish the superiority of faith over knowledge, could
not advance at all. And in the proof of this proposition, the first
premise---that a consciousness can exist without being theoretical, but
cannot be theoretical without existing---also contains an untruth. All
being in the form of individuality is existence; there can therefore be
an existence that is not consciousness, not theoretical. But when
consciousness becomes the subject of existence, existence includes
theory, for consciousness is only existing as being theoretical.

The unclarity arises from Nielsen's taking up Kierkegaard's
determinations but altering the presuppositions. Kierkegaard accentuated
in the determination ``existence'' the alternation between being and
non-being, which set the existing person's thought into becoming and cut
off objective knowledge's certainty from it. He then obtained the
opposition between existential (subjective) thinking and objective
thinking, and subjective thinking became for him the thinking whose goal
was existential inwardness, existing in what is thought. Nielsen lets
knowledge unfold itself into a system, with a different certainty and a
different relation to the system of existence than in Kierkegaard.
Existence is then conceived not as that which, as a determination for
the thinker, prevents objectively certain knowledge; rather, existence
receives the determination of particularity in opposition to
universality, and since universality is the determination of knowledge,
existence is posited as practical. The fact that existence, in
Kierkegaard's conception, becomes that to which regard must be paid,
from which one cannot abstract without conflict with the ethical and
without inner contradiction, becomes for Nielsen the claim that
existence, i.e.\ the practical, is the presupposition; but in order to
be asserted as such, it must be taken in the double meaning demonstrated
above.

A similar unclarity and arbitrariness appears where the opposition is
determined by the subjective and the objective, of which the former
relates to life, the latter to science. Where the opposition is
determined by these forms (e.g.\ in the \textit{Philosophisk
Prop\ae deutik} \S15), ``subjectivity'' and ``the subjective'' are taken in
a double meaning, so that it now comes to express the accidentally
individual and the egoistic, which falls outside the objectively valid,
and now the universal form of subjectivity. Here too the relation to
Kierkegaard is determinative. For Kierkegaard, subjective truth
designated appropriated truth. This truth is, in the territory of the
ethical, the universal that is realized in existence and that must be
thought in order that one may exist in it. In the territory of the
paradoxically religious, the truth that is to be appropriated has the
determination of the paradox. In all spheres, the appropriation of truth
is posited, because existence prevents objective certainty, as
conditioned by an act that everywhere can be designated as an act of
faith. In Nielsen, where it is not the uncertainty of existence but the
abstract opposition between the universal and the particular that
conditions the disjunction between knowledge and existence, the
subjective, in relation to the objectivity of knowledge, is placed
under the determinateness of a one-sided opposition, as being in
individuality's opposition to the universal, and from this comes the
double meaning, in that partly the Kierkegaardian determination
continues to operate, and partly the new meaning exclusively comes
forward. But in that the concept of subjectivity is conceived in this
meaning through a one-sided and abstract opposition, it is overlooked
that the form of subjectivity also contains what is objectively valid,
that the subject takes it up as a norm in itself without the content
being changed, that subjective action can have objective principles,
that there is an objectively true and valid element in
reason-determined subjectivity. The standpoint of subjectivity, of life,
comes in that conception to coincide with the standpoint of egoism.

In the same way as the concept of subjectivity, the concept of
objectivity also receives a double meaning. Where the main stages of
subjectivity are developed, the objective now means objective evil or
the merely abstract universal, and now the universal and universally
valid that encompasses the totality of the subjective. And as the
abstract opposition is definitively maintained, it is arbitrarily
ignored that at the different stages the subject's determinateness---
which especially appears clearly at the stage designated as that of
infinite subjectivity---is precisely by virtue of what has been
determined as the objective principle of existence, and that it
expresses this. It is overlooked that subjective truth is precisely the
appropriated and, in the appropriation, essentially identical objective
truth.

The opposition between life and science, subjective and objective, is
led by Nielsen back to the opposition between will and cognition, which
becomes the psychological foundation for the opposition between faith
and knowledge, just as this last opposition has its metaphysical
foundation in the opposition of the fundamental ideas, in the
ontological fundamental determinations of absolute subjectivity. The
opposition between will and cognition, the fundamental psychological
opposition, must become the critique's proper starting point; the
examination of the determination of their relation must be the
fundamental one; and from the psychological we shall then, through the
ethical, be led to the metaphysical, where the definitive decision is
to take place and the final result to be won.

\subsection{The Psychological Opposition between Will and Knowledge}
\label{ssec:psych-opposition}

What is rightly urged by Nielsen is that will, as a primitive function
of spirit, has a distinctiveness that makes it impossible to reduce it
to thought, in the sense that it should become a mere modification of
the latter. From the fact that it, like cognition, has its
distinctiveness, there must then follow a total activity of spirit, a
presence of all the spirit's functions within each of the two
activities of spirit, and a distinctive fundamental form for these;
and they must in this distinctive form be objectified for
consciousness, which encompasses both the theoretical and the
practical.

But when Nielsen rightly asserts the distinctiveness, and rightly polemicizes against a
one-sidedness which, in the form in which he combats it, has no other representative than
Spinoza, he loses his right when by a subreption he transforms the distinctiveness into an
\emph{opposition} that is carried over from the form to the content, and which is more
closely determined as an \emph{absolute heterogeneity between the two activities of spirit
with respect to content, goal, and ideal}; and he becomes inconsistent both when he posits
the absolute heterogeneity while demanding the therewith incompatible unity in the absolute
principle, and when he holds fast that absolute heterogeneity, although he, by one-sidedly
setting the will as the properly constitutive in the human essence, lets it become a higher
than knowledge and ground the latter.

The untruth in these further-reaching determinations will become clear by a closer
investigation of the way in which Nielsen apprehends the relation.

When we set out from the concession that cognition and will must be set as distinctive,
self-valid forms of utterance of the essence, which cannot be reduced to modifications of
each other, then it at once becomes the task to determine how much lies in this concession
and whether it must lead to the postulated consequences: the relation of subordination and
the absolute heterogeneity. Nielsen determines the will's relation to the self in the
proposition: ``the most subjective of all in subjectivity is the self; but the selfish in the
self is self-determination, is will.'' By using the concept self-determination in such a way
that it at once becomes the essential expression for the selfhood-energy and is identified
with the will, the emphasis falls on this latter; it appropriates---even though thinking's,
like also the artistic fantasy's, original distinctiveness is pronounced---in reality the
primacy, and the nearest and immediate consequence becomes that the placing of will and
knowledge side by side, to which the principled opposition leads, becomes the subordination
of the knowledge-principle under the will-principle: a subordination that during the theory's
development comes forward in many forms. To this apprehension contributes also the above-
demonstrated one-sided apprehension of the concept existence (life), which identifies it
exclusively with the practical, with the will (``life is in the will''). The relation
presents itself otherwise when that arbitrariness is avoided, when we, in accordance with
what was intimated at an earlier point, in the apprehension of the human set out from the
determination that categorically distinguishes human life from life's lower stages: from the
determination of essence-universality, which expresses itself in the being-for-itself in the
form of liberated universality that is self-consciousness's, and when we apprehend the
essence's fundamental determination as \emph{energy}, more closely as selfhood-energy. Both
forms: cognition and will, then become primitive expressions for the energy of the essence
revealed in self-consciousness as unity, for the self's self-determination through which it,
determining itself, realizes itself: in both functions selfhood and energy assert themselves
under distinctive fundamental forms. The task then becomes to determine more closely the
distinctiveness of the forms of utterance, and the determination of it will decide the
question about the relation of opposition and more closely illuminate the postulated relation
of subordination. The distinctiveness is to be determined in such a way that in cognition the
selfhood-energy has most nearly the form of universality and ideality, in the will that of
singularity and reality. But these determinations, given as points of departure, point to
each other and determine themselves toward each other. The reality that is constituted by the
will is a realization of the ideal essence's determinations in the self or outside the self.
The ideality that is realized through cognition is an idealizing of a real conditioned by the
self's side of reality, and is such even in the ideal self's self-cognition. And in the same
way the universality that characterizes cognition is a universality that strives toward
concretion, punctuality, and whose final goal is the concrete cognition of existence, in
which the cognition of the self according to its individualized essence-determinations and
essential relations is comprised; in other words: the individual-universal essence's
transparency to itself as a member in a totality of existence. Conversely the determination of
singularity in the will is such a one as with necessity unfolds the universal in itself,
because the will is only will as the self-conscious, i.e. according to its concept universal
and general, essence's will. The will must unfold itself from the will determined by the
single, positing the single, to the will determined by the universally rational that expresses
the essence, realizing the universal. The will-utterance's form of singularity is here not
abolished, but the determinateness of singularity has here become transparent in the universal.
In a manner similar to the determinations of universality and singularity it stands with the
therewith so closely connected ones: those of objectivity and subjectivity. Cognition has most
nearly the stamp of objectivity; through the determinateness by otherness, through the
subject's relation to the for all cognizing individuals common, universally valid forms of
thought and laws of thought, which determine it with necessity. But cognition completes itself
only in the cognition embracing the subject as it were concentrated about the self clear to
itself, and it is, as realization of an essence-determination of the subject, as ideal
extension of the self, object for an interest that can have the most intensive stamp of
subjectivity, as a \emph{passion} of cognition, and in it there is in general, because it is
\emph{essence-utterance}, \emph{the subjective essence active according to its wholeness, its
undivided personality}. On the other side the will's subjectivity has, the more the will-
realization removes itself from the instinctual, drive-like, and egoistic, the more it becomes
a realization of the universally rational, of the essence's law, an essence-willing, the
determinateness of objectivity, of the universally valid, in itself.

At the will's and cognition's imperfect stage of development they diverge and can come forward
in isolation; at their culminating stage, i.e. at the stage of rational determinateness, there
arises, while the distinctiveness is preserved, the most indissoluble connection and the most
complete harmony. The cognition that has reached its goal would, as expression for the in-itself
transparent essence's energy, in accordance with the one essence's necessary being in both
forms, immediately convert itself into the will's energy determined by the essence-cognition,
the essence's self-revealedness, and the will that with the most intensive energy, in the most
embracing meaning, was essence-willing: realization of the essence according to its concretion,
would be a will that immediately converted itself into consciousness, into cognition, and into a
cognition that was determined by knowledge's universal, objectively valid laws.

By what is here developed there is thus abolished the sharp opposition posited by Nielsen between
will and knowledge, and between the corresponding ``practical'' and ``theoretical
consciousness.'' The practical consciousness is to be personal, interested, in conflict with
existence, moved by fear and hope, full of self-regard; the theoretical impersonal, contemplative,
disinterested, lost in the universal, moved by necessity, forgetting itself in the grounding of the
objective. The relative truth in these oppositions is acknowledged through what is developed above,
but the abstract and absolute opposition is denied.

And when the question is raised: whether, without the two essence-utterances' distinctiveness and
primitivity being given up, there can be ascribed to one of them an ideal primacy, then this
question must be answered in another way than Nielsen has done it by ascribing the primacy to the
will; which he does partly for the reason stated above, and partly because cognition is to stop at
the universal and at the essence-necessity and not reach the single as such and freedom. From the
point of departure indicated above the primacy must be ascribed to cognition. The essence's
universality expresses itself \emph{immediately} in the form of liberated universality that is
cognition's fundamental form; the determinateness of singularity, which is the will's form, rests
on the ground of universality, and the will is only will as the cognizing essence's will. The
ethical willing as primitive decision, which creates the relation to the good, in which the will
constitutes itself as ethical, and as concrete willing of the good in its determinateness, is
conditioned by knowledge and is only possible through it. The consciousness of the \emph{obligating
presupposes in its clarity and certainty thought}.

But when thus the primacy is ascribed to cognition, which in the human essence thereby acquires
the same relation to the will as in the Absolute the determination of ideality to the determination
of reality; then, in accordance with what is developed above, both essence-utterances'
distinctiveness is to be held fast in order that the will not be reduced to a modification of
thinking. In the transition from the one essence-utterance to the other there takes place a
\emph{conversion, a transformation of form}, whose kind definitively determines the relation.
Inasmuch as the content of cognition is converted into a content of will, the conversion is most
nearly to be characterized as a conversion from the sphere of reflection and of the universal to
that of immediate being and of singularity. Inasmuch as the content of will is converted into a
content of cognition---and this means not only: formally reproduced by consciousness---the
conversion takes place in the opposite direction. In the moment of decision that separated by
reflection is closed together into unity, reflection goes to rest in the decision's act of will,
and there is acted in this sense reflectionlessly. Conversely, inasmuch as the act of will with its
content is converted into cognition, reflection is again called forth, the immediate unity is
dissolved into reflection's distinction, the action's single is translated into cognition's
universal.

What then is the categorical determinateness of the relation in accordance with this conversion?
The distinctiveness of the essence-utterances brings a \textit{differentia specifica}, a difference
of kind, in between will and cognition, but how is the heterogeneity to be more closely determined?
We will here see the matter from the psychological standpoint alone, and leave aside the postulate
that an absolute heterogeneity is to be grounded in the relations of the ideas. The question then
acquires the form: whether the heterogeneity becomes such a one that it, in Nielsen's expression,
allows only a reflection of the will in knowledge as a foreign, only by itself determined, which
does not become transparent for knowledge, or whether the difference is governed by the
essence-unity. The answer depends essentially on the apprehension of the determination of
singularity in the will. Is it overlooked, in the apprehension of this, that the will is only will
as the self-conscious, i.e. according to its nature universal and general, essence's will, then the
will's singularity and the moment of immediacy in it is traced back to the immediate
nature-determinateness's utterance of activity in the drive, the sensible-egoistic craving, and
inasmuch as this is most nearly seen from the side through which it hangs together with the
individual nature-being's unopened singularity, which stands in \emph{immediate opposition to
knowledge's universal}, then the activity that proceeds from it is apprehended as everywhere
determined by this form of singularity, and, where it does not in general withdraw itself from
consciousness, as that which at most is reproduced by consciousness as the in its immediacy with
cognition's universal absolutely incommensurable. Is, on the contrary, this apprehension of
singularity abandoned---which would little agree with freedom's being ascribed to the will---and is
it held fast that the will's singularity is the universal essence's form of appearance, is
\emph{spirit-determinateness}, the relation becomes an essentially other. In the thus determined
single the universal goes to rest only in such a way that the single becomes transparent: the
conversion is a conversion from thinking's, and expressly from the ``theoretical'' thinking's,
universality to the personal living's singularity; but the individual as rationally willing
expresses in the singularity the universal, gives it being in its form, which transparently reveals
the universal. The difference of form is accordingly governed by the essential unity; the content of
consciousness becomes under both forms: as content of cognition and content of will, essentially the
same. I cognize duty---and that is again: I cognize it theoretically---e.g. the duty to fight for
the fatherland; I resolve to exist in the cognized, and inasmuch as the resolve is realized, the
content, as realized in the single will, is essentially the same. The unity comes forward most
determinately in the relation where cognition itself becomes the will's object.

% [translation continues from p. 138]

\subsection{The Ethical and the Religious in Relation to Cognition}
\label{ssec:ethical-religious}

% [text to be added: pp. 161--171]

\subsection{The Presuppositions of the Psychological Determinations
  in the Metaphysical}
\label{ssec:metaphysical}

% [text to be added: pp. 171--178]

\section{The Theory's Consequences}
\label{sec:consequences}

\subsection{With Respect to the Psychological}
\label{ssec:conseq-psych}

% [text to be added: pp. 178--195]

\subsection{With Respect to the Ethical}
\label{ssec:conseq-ethical}

% [text to be added: pp. 195--199]

\subsection{With Respect to the Religious}
\label{ssec:conseq-religious}

% [text to be added: pp. 199--210]

\section{Result of the Critique of Nielsen's Theory}
\label{sec:nielsen-result}

% [text to be added: pp. 210--212]

\section{What is Common to Kierkegaard and Nielsen and What is
  Distinctive to Each}
\label{sec:common-distinctive}

% [text to be added: pp. 212--216]

\section{Critique of the Determinations Distinctive to Kierkegaard}
\label{sec:kierkegaard-critique}

% [text to be added: pp. 216--224]

\section{Final Result}
\label{sec:separation-result}

% [text to be added: p. 224]


%% ============================================================
%%  CONCLUSION (pp. 225--226)
%% ============================================================
\chapter*{Conclusion}
\addcontentsline{toc}{chapter}{Conclusion}
\label{ch:conclusion}

% [text to be added: pp. 225--226]

\end{document}
